%% Volume I, Chapter 21: The Singular Value Decomposition
\chapter{The Singular Value Decomposition}
\label{ch:svd}

\abstract{Every matrix, square or not, symmetric or not, becomes diagonal once orthonormal coordinates are chosen well in the space where it acts and in the space where its values lie. The numbers on that diagonal, the singular values, answer at once questions that earlier chapters approached one at a time: what the rank of a matrix is when its entries carry errors, how much it can stretch a vector, how far it is from a singular matrix, which matrix of lower rank approximates it best, and which least squares solution is shortest. For data, the decomposition sorts the directions of variation by importance; in computation, it is the dependable measure of how nearly singular a matrix is.}

Many threads of this part of the book end here. \Cref{ch:several-variables} defined the operator norm $\|A\|_2$ as a supremum and could only bound it by the Frobenius norm, and \cref{sec:trace-norms} gave the value $\|A\|_2=(15+5\sqrt5)^{1/2}$ for one $2\times2$ matrix without proof. \Cref{ch:vectors-matrices} introduced the condition number $\kappa(A)=\|A\|\|A^{-1}\|$, singled out the left inverse $(1,2)/5$ of the column $(1,2)$ among all its left inverses, and read a matrix product as a sum of outer products; \cref{ch:linear-systems} found the condition number in the errors of the Hilbert systems; \cref{ch:vector-spaces} computed the rank of matrices in floating-point arithmetic and found that the library routine reports a \emph{numerical} rank, whose meaning was left open; \cref{ch:linear-maps} reduced every matrix to the rank normal form $N_r$ by a change of bases; \cref{ch:determinants} warned that a determinant says little about nearness to singularity; \cref{ch:inner-products} characterised the shortest least squares solution and observed, for one example, that the normal equations square the condition number; and \cref{ch:eigenvalues,ch:symmetric-matrices} drew the image of the unit circle under a matrix, an ellipse, whose axes are eigenvector lines only when the matrix is symmetric. All these questions have one answer.

The answer is the singular value decomposition $A=U\Sigma V^{\mathsf H}$ of an arbitrary matrix $A\in\mathbb F^{m\times n}$, with unitary matrices $U$ and $V$ and a matrix $\Sigma$ that is zero except for the nonnegative numbers $\sigma_1\ge\sigma_2\ge\cdots$ on its diagonal. The columns of $V$ and of $U$ are orthonormal bases of $\mathbb F^n$ and of $\mathbb F^m$, and $A$ maps the $i$th vector of the first basis to $\sigma_i$ times the $i$th vector of the second. From this one statement follow the norm $\|A\|_2=\sigma_1$, the rank as the number of nonzero $\sigma_i$, orthonormal bases of the four fundamental subspaces, the best approximation of $A$ by matrices of lower rank (the theorem of Eckart and Young), the distance to the nearest singular matrix, the pseudoinverse and the minimum norm least squares solution, and the condition number $\sigma_1/\sigma_r$. The singular values are perfectly conditioned, like the eigenvalues of Hermitian matrices and unlike those of general matrices, and the decomposition yields the polar decomposition $A=WP$, the analogue of $z=e^{i\theta}|z|$, and the principal components of a table of data.

As in the preceding chapters, $\mathbb F$ denotes $\mathbb R$ or $\mathbb C$. For $A\in\mathbb F^{m\times n}$ we write
\[
p\coloneqq\min(m,n)
\]
for the number of diagonal positions of an $m\times n$ matrix. The singular values are numbered in \emph{decreasing} order, $\sigma_1(A)\ge\sigma_2(A)\ge\cdots\ge\sigma_p(A)\ge0$, as is customary, and it is convenient to set $\sigma_i(A)\coloneqq0$ for $i>p$. \Cref{ch:symmetric-matrices} numbered the eigenvalues of a Hermitian matrix in \emph{increasing} order, $\lambda_1\le\cdots\le\lambda_n$, so that the two conventions run in opposite directions; the relation between them, which \cref{prop:singular-eigen} proves, is
\begin{equation}
\sigma_i(A)^2=\lambda_{n+1-i}(A^{\mathsf H}A)\qquad(i=1,\ldots,n),
\label{eq:sigma-lambda}
\end{equation}
so that the largest singular value belongs to the largest eigenvalue $\lambda_n(A^{\mathsf H}A)$ and the smallest to the smallest. The Schur and Jordan forms of \cref{ch:canonical-forms} are not used. In this chapter, as in \cref{ch:symmetric-matrices}, $H_n$ denotes the Hilbert matrix with the entries $1/(i+j-1)$, not the harmonic number of the Notation.

The chapter is organised as follows. \Cref{sec:svd-geometry} finds the decomposition of a $2\times2$ matrix from the picture of the image of the unit circle. \Cref{sec:svd-existence} proves that every matrix has a singular value decomposition, in two ways, describes how unique it is and relates singular values to eigenvalues, and \cref{sec:svd-subspaces} reads off the four fundamental subspaces and refines the rank normal form. \Cref{sec:svd-norms} expresses norms by singular values, \cref{sec:low-rank} proves the theorem of Eckart and Young on low-rank approximation and gives the numerical rank its precise meaning, and \cref{sec:svd-perturbation} shows that singular values are perfectly conditioned. \Cref{sec:pseudoinverse} treats the pseudoinverse, its discontinuity and its regularised forms, and \cref{sec:svd-condition} the condition number. \Cref{sec:polar-pca} derives the polar decomposition, the nearest unitary matrix and principal component analysis, and \cref{sec:svd-computation} tests several statements with the stardust library.

\section{The Geometry of a Linear Map}
\label{sec:svd-geometry}

\subsection{The Image of the Unit Circle}

\Cref{fig:unit-circle-image} showed the image of the unit circle under the two matrices of \eqref{eq:two-by-two-examples}. For the symmetric matrix $A$ the image is an ellipse whose axes are the eigenvector lines, a fact that the spectral theorem explains (\cref{fig:spectral-action}). For the matrix
\[
B=\begin{bmatrix}1&1\\0&2\end{bmatrix}
\]
the image is again an ellipse, but its axes are not the eigenvector lines, which are not even perpendicular. We now find the axes.

A point $B\mathbf x$ of the image lies farthest from the origin, or closest to it, when $\|B\mathbf x\|^2$ is largest or smallest among the unit vectors $\mathbf x$. This squared length is the quadratic form of the symmetric matrix
\[
B^{\mathsf T}B=\begin{bmatrix}1&1\\1&5\end{bmatrix},\qquad\|B\mathbf x\|^2=\langle B\mathbf x,B\mathbf x\rangle=\langle\mathbf x,B^{\mathsf T}B\mathbf x\rangle ,
\]
and by \cref{thm:rayleigh} its extreme values on the unit circle are the eigenvalues of $B^{\mathsf T}B$, attained at the eigenvectors. The trace is $6$ and the determinant $4$, so the eigenvalues are $3\pm\sqrt5$ by \eqref{eq:eigen-2x2}. With the golden ratio $\phi=(1+\sqrt5)/2$, which satisfies $\phi^2=\phi+1$, one checks that $2\phi^2=3+\sqrt5$ and $2/\phi^2=3-\sqrt5$, so the square roots of the eigenvalues are
\[
\sigma_1=\sqrt2\,\phi=2.2882\ldots,\qquad\sigma_2=\sqrt2/\phi=0.8740\ldots,
\]
with $\sigma_1\sigma_2=2=\det B$. The longest image of a unit vector has the length $\sigma_1$ and the shortest the length $\sigma_2$. The unit eigenvectors of $B^{\mathsf T}B$ are
\[
\mathbf v_1=\frac{(1,\phi^3)}{\sqrt{1+\phi^6}}=(0.2298,0.9732),\qquad\mathbf v_2=(-0.9732,0.2298),
\]
up to signs; for the first, the equation $(1-\lambda)x_1+x_2=0$ with $\lambda=3+\sqrt5$ gives $x_2=(2+\sqrt5)x_1=\phi^3x_1$. They are orthogonal, as eigenvectors of a symmetric matrix for different eigenvalues must be (\cref{prop:hermitian-eigen}). Their images are orthogonal as well, since $\langle B\mathbf v_1,B\mathbf v_2\rangle=\langle\mathbf v_1,B^{\mathsf T}B\mathbf v_2\rangle=(3-\sqrt5)\langle\mathbf v_1,\mathbf v_2\rangle=0$, and they have the lengths $\sigma_1$ and $\sigma_2$. Writing them as $B\mathbf v_1=\sigma_1\mathbf u_1$ and $B\mathbf v_2=\sigma_2\mathbf u_2$ defines a second pair of orthonormal vectors,
\[
\mathbf u_1=\frac{(1,\phi)}{\sqrt{1+\phi^2}}=(0.5257,0.8507),\qquad\mathbf u_2=(-0.8507,0.5257).
\]

Every unit vector is $\mathbf x=c_1\mathbf v_1+c_2\mathbf v_2$ with $c_1^2+c_2^2=1$, and its image is $B\mathbf x=c_1\sigma_1\mathbf u_1+c_2\sigma_2\mathbf u_2$. In the coordinates $y_1=\sigma_1c_1$ and $y_2=\sigma_2c_2$ along $\mathbf u_1$ and $\mathbf u_2$, the image of the circle is the curve $y_1^2/\sigma_1^2+y_2^2/\sigma_2^2=1$: an ellipse whose semi-axes are $\sigma_1\mathbf u_1$ and $\sigma_2\mathbf u_2$. The axes of the ellipse are the lines spanned by $\mathbf u_1$ and $\mathbf u_2$, the lengths of the semi-axes are $\sigma_1$ and $\sigma_2$, and the unit vectors that are mapped onto the ends of the axes are $\pm\mathbf v_1$ and $\pm\mathbf v_2$. The numbers $\sigma_i$ are the \term{singular values} of $B$, the vectors $\mathbf u_i$ its \term{left singular vectors} and the vectors $\mathbf v_i$ its \term{right singular vectors}. The eigenvectors of $B$, the $x_1$-axis and the line spanned by $(1,1)$, play no role in this description.

\subsection{Three Steps}

The relations $B\mathbf v_1=\sigma_1\mathbf u_1$ and $B\mathbf v_2=\sigma_2\mathbf u_2$ say $BV=U\Sigma$, with $V$ and $U$ the orthogonal matrices with the columns $\mathbf v_i$ and $\mathbf u_i$ and $\Sigma=\operatorname{diag}(\sigma_1,\sigma_2)$, and since $V^{-1}=V^{\mathsf T}$,
\[
B=U\Sigma V^{\mathsf T}=\begin{bmatrix}0.5257&-0.8507\\0.8507&0.5257\end{bmatrix}\begin{bmatrix}2.2882&0\\0&0.8740\end{bmatrix}\begin{bmatrix}0.2298&0.9732\\-0.9732&0.2298\end{bmatrix}.
\]
With the signs chosen here both $U$ and $V$ are rotations, $U=R_\alpha$ with $\alpha=\arctan\phi$, about $58.28$ degrees, and $V=R_\beta$ with $\beta=\arctan\phi^3$, about $76.72$ degrees. Read from right to left, $B\mathbf x=U\Sigma V^{\mathsf T}\mathbf x$ applies three maps in turn, as $Q\Lambda Q^{\mathsf T}$ did in \cref{fig:spectral-action}: $V^{\mathsf T}$ rotates the plane by $-\beta$ and takes $\mathbf v_1$ and $\mathbf v_2$ to $\mathbf e_1$ and $\mathbf e_2$; $\Sigma$ stretches the first coordinate by $\sigma_1$ and the second by $\sigma_2$; and $U$ rotates by $\alpha$ and takes $\mathbf e_1$ and $\mathbf e_2$ to $\mathbf u_1$ and $\mathbf u_2$. \Cref{fig:svd-action} follows the unit circle and the point $\mathbf e_2$ through the three steps; the image of $\mathbf e_2$ must be the second column $(1,2)$ of $B$, and indeed $V^{\mathsf T}\mathbf e_2=(0.9732,0.2298)$, $\Sigma V^{\mathsf T}\mathbf e_2=(2.2270,0.2008)$ and $U\Sigma V^{\mathsf T}\mathbf e_2=(1.0000,2.0000)$ to four digits.

\begin{figure}
\newcommand\sdsvdpanel[5]{\begin{tikzpicture}[scale=0.8, font=\scriptsize]
\draw[black!35, line width=0.4pt] (0,0) circle (1);
\draw[line width=0.7pt, smooth, samples=100, domain=0:360] plot ({#1*cos(\x)+#2*sin(\x)}, {#3*cos(\x)+#4*sin(\x)});
\draw[line width=0.4pt, -{Latex[length=1.3mm]}] (-2.45,0) -- (2.55,0);
\draw[line width=0.4pt, -{Latex[length=1.3mm]}] (0,-2.15) -- (0,2.25);
#5
\end{tikzpicture}}
% v1 at 76.72 degrees, v2 at 166.72 degrees; u1 at 58.28 degrees, u2 at 148.28 degrees; the point x = e2
\sdpanel{A}{\sdsvdpanel{1}{0}{0}{1}{%
\draw[line width=0.9pt, -{Latex[length=1.6mm]}] (0,0) -- (0.2298,0.9732);
\node[right, inner sep=2pt] at (0.25,0.62) {$\mathbf v_1$};
\draw[line width=0.9pt, -{Latex[length=1.6mm]}] (0,0) -- (-0.9732,0.2298);
\node[above, inner sep=2pt] at (-0.62,0.32) {$\mathbf v_2$};
\fill (0,1) circle (1.5pt);
\node[above left, inner sep=1pt] at (-0.02,1.02) {$\mathbf e_2$};}}\hfill
\sdpanel{B}{\sdsvdpanel{1}{0}{0}{1}{%
\draw[line width=0.9pt, -{Latex[length=1.6mm]}] (0,0) -- (1,0);
\node[below, inner sep=2pt] at (0.75,0) {$\mathbf e_1$};
\draw[line width=0.9pt, -{Latex[length=1.6mm]}] (0,0) -- (0,1);
\node[left, inner sep=2pt] at (0,0.75) {$\mathbf e_2$};
\fill (0.9732,0.2298) circle (1.5pt);
\node[above right, inner sep=1pt] at (0.98,0.24) {$V^{\mathsf T}\mathbf e_2$};}}
\sdpanelrow
\sdpanel{C}{\sdsvdpanel{2.2882}{0}{0}{0.8740}{%
\draw[line width=0.9pt, -{Latex[length=1.6mm]}] (0,0) -- (2.2882,0);
\node[anchor=north, inner sep=1pt] at (1.55,-0.1) {$\sigma_1\mathbf e_1$};
\draw[line width=0.9pt, -{Latex[length=1.6mm]}] (0,0) -- (0,0.874);
\node[left, inner sep=2pt] at (0,0.5) {$\sigma_2\mathbf e_2$};
\fill (2.2270,0.2008) circle (1.5pt);
\node[above, inner sep=2pt] at (2.05,0.95) {$\Sigma V^{\mathsf T}\mathbf e_2$};
\draw[line width=0.3pt] (2.2270,0.2008) -- (2.12,0.9);}}\hfill
\sdpanel{D}{\sdsvdpanel{1}{1}{0}{2}{%
\draw[black, dashed, line width=0.4pt] (-1.3,-2.1) -- (1.3,2.1);
\draw[black, dashed, line width=0.4pt] (-1.75,1.08) -- (1.75,-1.08);
\draw[line width=0.9pt, -{Latex[length=1.6mm]}] (0,0) -- (1.2030,1.9465);
\node[anchor=west, inner sep=1pt] at (1.5,1.1) {$\sigma_1\mathbf u_1$};
\draw[line width=0.3pt] (1.5,1.1) -- (0.85,1.38);
\draw[line width=0.9pt, -{Latex[length=1.6mm]}] (0,0) -- (-0.7435,0.4595);
\node[anchor=east, inner sep=1pt] at (-1.3,0.2) {$\sigma_2\mathbf u_2$};
\draw[line width=0.3pt] (-1.3,0.2) -- (-0.6,0.37);
\fill (1,2) circle (1.5pt);
\node[left, inner sep=2pt] at (1,2.08) {$B\mathbf e_2$};}}
\caption{The action of $B=U\Sigma V^{\mathsf T}$, the matrix with the rows $(1,1)$ and $(0,2)$, followed on the unit circle (grey) and on the point $\mathbf e_2$ (dot); the black curve is the image of the circle after each step. (A) The right singular vectors $\mathbf v_1$ and $\mathbf v_2$. (B) After $V^{\mathsf T}$, a rotation by about $-76.7$ degrees, they have become $\mathbf e_1$ and $\mathbf e_2$, and the circle is unchanged. (C) $\Sigma$ stretches the coordinates by $\sigma_1=\sqrt2\,\phi\approx2.288$ and $\sigma_2=\sqrt2/\phi\approx0.874$. (D) $U$, a rotation by about $58.3$ degrees, turns the ellipse into the image of the circle under $B$, with the semi-axes $\sigma_1\mathbf u_1$ and $\sigma_2\mathbf u_2$ along the dashed lines, and $\mathbf e_2$ arrives at $B\mathbf e_2=(1,2)$. Compare \cref{fig:spectral-action}, where the first and the last step are inverse to each other.\label{fig:svd-action}}
\end{figure}

The comparison with \cref{fig:spectral-action} shows what is new. For a symmetric positive definite matrix the two rotations are inverse to each other, $U=V=Q$, and the right singular vectors are the eigenvectors. For a general matrix the first and the last step are independent: the orthonormal basis that is mapped onto orthogonal vectors is not, in general, mapped onto multiples of itself. What survives from the spectral theorem is that \emph{some} orthonormal basis is mapped onto an orthogonal set, and that this basis is found from the symmetric matrix $B^{\mathsf T}B$.

\subsection{From the Plane to Every Dimension}

Nothing in this computation used that $B$ was $2\times2$. For $A\in\mathbb F^{m\times n}$ the matrix $A^{\mathsf H}A$ is Hermitian and positive semidefinite, its orthonormal eigenvectors $\mathbf v_1,\ldots,\mathbf v_n$ exist by the spectral theorem, and the images $A\mathbf v_i$ are orthogonal, since $\langle A\mathbf v_i,A\mathbf v_j\rangle=\langle\mathbf v_i,A^{\mathsf H}A\mathbf v_j\rangle=\lambda_j\langle\mathbf v_i,\mathbf v_j\rangle$. The next section turns this observation into a theorem. Geometrically it says that the image of the unit sphere of $\mathbb F^n$ under any matrix is an ellipsoid, flattened in the directions of the zero singular values, whose semi-axes have the lengths $\sigma_i$ and point along the left singular vectors; this is the precise sense in which every linear map is a stretching along orthogonal axes, followed and preceded by isometries.

\section{Existence and Uniqueness}
\label{sec:svd-existence}

\subsection{The Theorem}

A matrix $\Sigma\in\mathbb R^{m\times n}$ is called \term{diagonal} if its entries $s_{ij}$ with $i\ne j$ vanish, whatever its shape; its diagonal entries are $s_{11},\ldots,s_{pp}$. For instance, a $3\times2$ diagonal matrix has the rows $(s_{11},0)$, $(0,s_{22})$ and $(0,0)$.

\begin{definition}[Singular value decomposition]
\label{def:svd}
Let $A\in\mathbb F^{m\times n}$ and $p=\min(m,n)$. A \term{singular value decomposition} (SVD) of $A$ is a factorisation
\[
A=U\Sigma V^{\mathsf H}
\]
with unitary matrices $U\in\mathbb F^{m\times m}$ and $V\in\mathbb F^{n\times n}$ and a diagonal matrix $\Sigma\in\mathbb R^{m\times n}$ whose diagonal entries satisfy $\sigma_1\ge\sigma_2\ge\cdots\ge\sigma_p\ge0$. The numbers $\sigma_i$ are the \term{singular values} of $A$, the columns $\mathbf u_1,\ldots,\mathbf u_m$ of $U$ are \term{left singular vectors} and the columns $\mathbf v_1,\ldots,\mathbf v_n$ of $V$ are \term{right singular vectors}. For a real matrix, a real SVD is one in which $U$ and $V$ are real orthogonal matrices, $A=U\Sigma V^{\mathsf T}$.
\end{definition}

Comparing the columns of $AV=U\Sigma$ and of $A^{\mathsf H}U=V\Sigma^{\mathsf T}$ shows what the decomposition says about the vectors:
\begin{equation}
A\mathbf v_i=\sigma_i\mathbf u_i,\qquad A^{\mathsf H}\mathbf u_i=\sigma_i\mathbf v_i\qquad(i=1,\ldots,p),
\label{eq:singular-pairs}
\end{equation}
while $A\mathbf v_i=\mathbf 0$ for $p<i\le n$ and $A^{\mathsf H}\mathbf u_i=\mathbf 0$ for $p<i\le m$. Conversely, orthonormal bases $\mathbf v_1,\ldots,\mathbf v_n$ and $\mathbf u_1,\ldots,\mathbf u_m$ with $A\mathbf v_i=\sigma_i\mathbf u_i$ for $i\le p$ and $A\mathbf v_i=\mathbf 0$ for $i>p$ give an SVD, because these relations say $AV=U\Sigma$. A pair $(\mathbf u_i,\mathbf v_i)$ with \eqref{eq:singular-pairs} is a \term{singular pair}, or a pair of singular vectors, for $\sigma_i$.

\begin{theorem}[Existence of the singular value decomposition]
\label{thm:svd}
Every matrix $A\in\mathbb F^{m\times n}$ has a singular value decomposition. A real matrix has a real singular value decomposition. The number $r$ of positive singular values is at most $p$, and the decomposition can be constructed so that $\mathbf v_1,\ldots,\mathbf v_n$ are orthonormal eigenvectors of $A^{\mathsf H}A$ for the eigenvalues $\sigma_1^2\ge\cdots\ge\sigma_n^2$, with $\sigma_i=0$ for $i>p$, and $\mathbf u_i=A\mathbf v_i/\sigma_i$ for $i\le r$.
\end{theorem}

\begin{proof}
The matrix $A^{\mathsf H}A\in\mathbb F^{n\times n}$ is Hermitian and positive semidefinite by \cref{thm:psd-criteria}\,(iii), so by \cref{thm:spectral} it has an orthonormal basis of eigenvectors with real nonnegative eigenvalues; real ones if $A$ is real. Number them in decreasing order of the eigenvalues, reversing the order of \cref{thm:spectral} by a permutation, so that $A^{\mathsf H}A\mathbf v_i=d_i\mathbf v_i$ with $d_1\ge\cdots\ge d_n\ge0$, and let $r$ be the number of positive $d_i$ and $\sigma_i\coloneqq\sqrt{d_i}$. For $i\le r$ put $\mathbf u_i\coloneqq A\mathbf v_i/\sigma_i$. These vectors are orthonormal:
\[
\langle\mathbf u_i,\mathbf u_j\rangle=\frac{\langle A\mathbf v_i,A\mathbf v_j\rangle}{\sigma_i\sigma_j}=\frac{\langle\mathbf v_i,A^{\mathsf H}A\mathbf v_j\rangle}{\sigma_i\sigma_j}=\frac{d_j\langle\mathbf v_i,\mathbf v_j\rangle}{\sigma_i\sigma_j}=\delta_{ij}\qquad(i,j\le r).
\]
In particular $r\le m$, since $\mathbb F^m$ contains at most $m$ orthonormal vectors, and $r\le n$, so $r\le p$. For $i>r$, $\|A\mathbf v_i\|^2=\langle\mathbf v_i,A^{\mathsf H}A\mathbf v_i\rangle=d_i=0$, so $A\mathbf v_i=\mathbf 0$. Extend $\mathbf u_1,\ldots,\mathbf u_r$ to an orthonormal basis $\mathbf u_1,\ldots,\mathbf u_m$ of $\mathbb F^m$ (\cref{cor:orthonormal-basis}), real if $A$ is real, and let $U$ and $V$ have these columns and $\Sigma\in\mathbb R^{m\times n}$ be diagonal with the entries $\sigma_1,\ldots,\sigma_p$. The $i$th column of $AV$ is $A\mathbf v_i$, which is $\sigma_i\mathbf u_i$ for $i\le r$ and $\mathbf 0$ for $i>r$; the $i$th column of $U\Sigma$ is $\sigma_i\mathbf u_i$ for $i\le p$, which is $\mathbf 0$ for $r<i\le p$, and $\mathbf 0$ for $i>p$. Hence $AV=U\Sigma$, and $A=U\Sigma V^{\mathsf H}$ because $V$ is unitary.
\end{proof}

The proof is constructive, and it shows what the singular vectors are: the right singular vectors are eigenvectors of $A^{\mathsf H}A$, and the left singular vectors, by the same argument applied to $A^{\mathsf H}$, are eigenvectors of $AA^{\mathsf H}$. It is not, however, how singular value decompositions are computed, because forming $A^{\mathsf H}A$ loses information, as \cref{sec:svd-computation} demonstrates. A second proof works with $A$ itself. It finds the largest singular value directly as the largest stretching factor and then splits off one dimension at a time, in the spirit of the proof of the spectral theorem.

\begin{proof}[Second proof of \cref{thm:svd}]
We prove by induction on $p=\min(m,n)$ that $A$ has an SVD whose largest singular value is $\|A\|_2$; the operator norm $\|\cdot\|_2$ is defined for complex matrices by the same supremum as for real ones (\cref{prop:spectral-radius-norm}). If $A=0$, then $A=I_m0I_n$ is an SVD. Let $A\ne0$. The function $\mathbf x\mapsto\|A\mathbf x\|$ is continuous on the unit sphere of $\mathbb F^n$, which is closed and bounded, so it attains its maximum $\sigma_1=\|A\|_2>0$ at some unit vector $\mathbf v_1$ (\cref{thm:evt-rn}, with $\mathbb C^n$ identified with $\mathbb R^{2n}$ if necessary). Let $\mathbf u_1\coloneqq A\mathbf v_1/\sigma_1$, a unit vector, and extend $\mathbf u_1$ and $\mathbf v_1$ to unitary matrices $U_1=[\,\mathbf u_1\ \tilde U\,]$ and $V_1=[\,\mathbf v_1\ \tilde V\,]$. Since $\tilde U^{\mathsf H}A\mathbf v_1=\sigma_1\tilde U^{\mathsf H}\mathbf u_1=\mathbf 0$,
\[
S\coloneqq U_1^{\mathsf H}AV_1=\begin{bmatrix}\mathbf u_1^{\mathsf H}A\mathbf v_1&\mathbf u_1^{\mathsf H}A\tilde V\\\tilde U^{\mathsf H}A\mathbf v_1&\tilde U^{\mathsf H}A\tilde V\end{bmatrix}=\begin{bmatrix}\sigma_1&\mathbf w^{\mathsf H}\\\mathbf 0&C\end{bmatrix}
\]
with $\mathbf w^{\mathsf H}\coloneqq\mathbf u_1^{\mathsf H}A\tilde V$ and $C\coloneqq\tilde U^{\mathsf H}A\tilde V\in\mathbb F^{(m-1)\times(n-1)}$. Unitary matrices preserve lengths, and $V_1\mathbf x$ runs through the unit vectors as $\mathbf x$ does, so $\|S\|_2=\|A\|_2=\sigma_1$. Apply $S$ to the vector $\mathbf z\coloneqq(\sigma_1,\mathbf w)\in\mathbb F^n$: the first component of $S\mathbf z$ is $\sigma_1^2+\mathbf w^{\mathsf H}\mathbf w$, so
\[
\sigma_1^2+\|\mathbf w\|^2\le\|S\mathbf z\|\le\sigma_1\|\mathbf z\|=\sigma_1\bigl(\sigma_1^2+\|\mathbf w\|^2\bigr)^{1/2},
\]
which gives $(\sigma_1^2+\|\mathbf w\|^2)^{1/2}\le\sigma_1$ and hence $\mathbf w=\mathbf 0$. So $A=U_1\operatorname{diag}(\sigma_1,C)V_1^{\mathsf H}$. If $p=1$, $C$ has no rows or no columns and we are done. Otherwise, by the induction hypothesis $C=U_2\Sigma_2V_2^{\mathsf H}$ with the largest singular value $\|C\|_2$, and $\|C\|_2\le\|S\|_2=\sigma_1$, since $\|C\mathbf y\|\le\|S(0,\mathbf y)\|$ for every $\mathbf y$; so the diagonal of $\Sigma_2$ does not exceed $\sigma_1$. Then
\[
A=U_1\begin{bmatrix}1&\mathbf 0^{\mathsf T}\\\mathbf 0&U_2\end{bmatrix}\begin{bmatrix}\sigma_1&\mathbf 0^{\mathsf T}\\\mathbf 0&\Sigma_2\end{bmatrix}\begin{bmatrix}1&\mathbf 0^{\mathsf T}\\\mathbf 0&V_2\end{bmatrix}^{\mathsf H}V_1^{\mathsf H}
\]
is an SVD, the products of unitary matrices being unitary. For a real matrix every step can be carried out with real vectors.
\end{proof}

The second proof shows directly what the first shows only after a computation: the largest singular value is the operator norm, $\sigma_1=\|A\|_2$, and a right singular vector for it is a unit vector that $A$ stretches most. \Cref{sec:svd-norms} develops this.

\subsection{Uniqueness}

The decomposition is not unique; already for $A=I_n$, every unitary $U$ gives an SVD $I=UIU^{\mathsf H}$. The singular values, however, are determined by $A$, and so are the singular vectors, up to the freedom that this example shows.

\begin{proposition}[Uniqueness]
\label{prop:svd-unique}
Let $A=U\Sigma V^{\mathsf H}$ be any singular value decomposition of $A\in\mathbb F^{m\times n}$.
\begin{enumerate}
\item The singular values are the square roots of the $p$ largest eigenvalues of $A^{\mathsf H}A$, and also of $AA^{\mathsf H}$; in particular, they do not depend on the decomposition, and neither does the number $r$ of positive ones.
\item For every $\sigma>0$, the right singular vectors $\mathbf v_i$ with $\sigma_i=\sigma$ form an orthonormal basis of the eigenspace $E_{\sigma^2}(A^{\mathsf H}A)$, and the corresponding left singular vectors one of $E_{\sigma^2}(AA^{\mathsf H})$. These subspaces do not depend on the decomposition.
\item If $\sigma_j>0$ is a simple singular value, that is, different from all other $\sigma_i$, then in every singular value decomposition the pair $(\mathbf u_j,\mathbf v_j)$ is the same up to a common factor $c$ with $|c|=1$: it is $(c\mathbf u_j,c\mathbf v_j)$.
\end{enumerate}
\end{proposition}

\begin{proof}
(i) and (ii). From $A=U\Sigma V^{\mathsf H}$ we get $A^{\mathsf H}A=V\Sigma^{\mathsf T}\Sigma V^{\mathsf H}$, where $\Sigma^{\mathsf T}\Sigma=\operatorname{diag}(\sigma_1^2,\ldots,\sigma_p^2,0,\ldots,0)\in\mathbb R^{n\times n}$. So the columns of $V$ are orthonormal eigenvectors of $A^{\mathsf H}A$, with the eigenvalues $\sigma_1^2,\ldots,\sigma_p^2$ and $n-p$ zeros. Since $A^{\mathsf H}A$ is similar to $\Sigma^{\mathsf T}\Sigma$, these are all its eigenvalues with their multiplicities, and the $p$ largest of them are the $\sigma_i^2$. For $\sigma>0$, the vectors $\mathbf v_i$ with $\sigma_i=\sigma$ are orthonormal vectors in $E_{\sigma^2}(A^{\mathsf H}A)$, and their number is the algebraic multiplicity of $\sigma^2$, which equals the dimension of the eigenspace by \cref{thm:spectral}; so they form a basis of it. The same argument with $AA^{\mathsf H}=U\Sigma\Sigma^{\mathsf T}U^{\mathsf H}$ gives the statements about $AA^{\mathsf H}$.

(iii) If $\sigma_j>0$ is simple, $E_{\sigma_j^2}(A^{\mathsf H}A)$ is a line by (ii), and the right singular vector $\mathbf v_j'$ of another decomposition lies on it and has length $1$, so $\mathbf v_j'=c\mathbf v_j$ with $|c|=1$. By \eqref{eq:singular-pairs}, $\mathbf u_j'=A\mathbf v_j'/\sigma_j=cA\mathbf v_j/\sigma_j=c\mathbf u_j$.
\end{proof}

A singular value $\sigma>0$ of multiplicity $k$ leaves more freedom: if $U_\sigma$ and $V_\sigma$ hold the $k$ singular pairs for $\sigma$, they may be replaced by $U_\sigma W$ and $V_\sigma W$ for any unitary $W\in\mathbb F^{k\times k}$, since $A(V_\sigma W)=\sigma U_\sigma W$, and by (ii) and \eqref{eq:singular-pairs} every other choice is of this form. The left singular vectors for the singular value $0$, and those beyond the $p$th, can be any orthonormal basis of the orthogonal complement of $\mathbf u_1,\ldots,\mathbf u_r$, and similarly on the right. One says ``the'' singular value decomposition when these choices do not matter, which is the case for everything in this chapter that is built from the products $\sigma_i\mathbf u_i\mathbf v_i^{\mathsf H}$ or from the subspaces of (ii).

\subsection{Reduced Forms and the Outer Product Expansion}

Only the first $r$ columns of $U$ and $V$ meet nonzero singular values, and the others are often dropped. Let $U_r\in\mathbb F^{m\times r}$ and $V_r\in\mathbb F^{n\times r}$ consist of the first $r$ columns of $U$ and $V$, and let $\Sigma_r\coloneqq\operatorname{diag}(\sigma_1,\ldots,\sigma_r)$. Multiplying out $U\Sigma V^{\mathsf H}$ by blocks, the rows of $\Sigma$ beyond the $r$th and its columns beyond the $r$th are zero, so
\begin{equation}
A=U\Sigma V^{\mathsf H}=U_r\Sigma_rV_r^{\mathsf H}=\sum_{i=1}^r\sigma_i\,\mathbf u_i\mathbf v_i^{\mathsf H},
\label{eq:svd-outer}
\end{equation}
where the last form is the outer product reading of the product $(U_r\Sigma_r)V_r^{\mathsf H}$ (\cref{thm:four-readings}\,(iv)). The factorisation $U_r\Sigma_rV_r^{\mathsf H}$ is the \term{compact} SVD. Keeping instead only the first $p$ columns of the larger of the two unitary factors, for instance $A=U_n\Sigma_nV^{\mathsf H}$ with $U_n\in\mathbb F^{m\times n}$ and $\Sigma_n\in\mathbb R^{n\times n}$ when $m\ge n$, gives the \term{thin} or \term{reduced} SVD, the form that library routines return on request; the stardust routine \code{linalg.svd} returns it by default, and the full one with the option \code{mode="complete"}. The full SVD is the right form for statements about bases of $\mathbb F^m$ and $\mathbb F^n$, the compact one for computations.

The expansion \eqref{eq:svd-outer} writes $A$ as a sum of $r$ matrices of rank one, which is the least possible number by \cref{thm:rank-factorisation}, since $r$ will turn out to be the rank. It is the expansion that \cref{ch:vectors-matrices} promised after \cref{thm:four-readings}: the terms are ordered by their size, since $\|\sigma_i\mathbf u_i\mathbf v_i^{\mathsf H}\|_{\mathrm F}=\sigma_i$ by \cref{prop:matrix-norms}\,(iv), and they are orthogonal to each other in the Frobenius inner product, since
\[
\operatorname{tr}\bigl((\mathbf u_i\mathbf v_i^{\mathsf H})^{\mathsf H}\mathbf u_j\mathbf v_j^{\mathsf H}\bigr)=\operatorname{tr}\bigl(\mathbf v_i(\mathbf u_i^{\mathsf H}\mathbf u_j)\mathbf v_j^{\mathsf H}\bigr)=\delta_{ij}\operatorname{tr}(\mathbf v_i\mathbf v_j^{\mathsf H})=\delta_{ij}
\]
by \cref{prop:trace}. \Cref{sec:low-rank} shows that cutting the sum after $k$ terms gives the best approximation of rank $k$.

\subsection{Examples}

For the $3\times2$ matrix with the rows $(1,1)$, $(1,0)$ and $(0,1)$,
\[
A=\begin{bmatrix}1&1\\1&0\\0&1\end{bmatrix},\qquad A^{\mathsf T}A=\begin{bmatrix}2&1\\1&2\end{bmatrix},
\]
the eigenvalues of $A^{\mathsf T}A$ are $3$ and $1$, with the unit eigenvectors $\mathbf v_1=(1,1)/\sqrt2$ and $\mathbf v_2=(1,-1)/\sqrt2$. So $\sigma_1=\sqrt3$, $\sigma_2=1$, and the construction of \cref{thm:svd} gives $\mathbf u_1=A\mathbf v_1/\sqrt3=(2,1,1)/\sqrt6$ and $\mathbf u_2=A\mathbf v_2=(0,1,-1)/\sqrt2$. A third vector completing these to an orthonormal basis of $\mathbb R^3$ is orthogonal to both, hence a multiple of the cross product of $(2,1,1)$ and $(0,1,-1)$, which is $(-2,2,2)$; so $\mathbf u_3=(-1,1,1)/\sqrt3$, and
\[
A=\begin{bmatrix}2/\sqrt6&0&-1/\sqrt3\\1/\sqrt6&1/\sqrt2&1/\sqrt3\\1/\sqrt6&-1/\sqrt2&1/\sqrt3\end{bmatrix}\begin{bmatrix}\sqrt3&0\\0&1\\0&0\end{bmatrix}\begin{bmatrix}1/\sqrt2&1/\sqrt2\\1/\sqrt2&-1/\sqrt2\end{bmatrix}.
\]
The vector $\mathbf u_3$ spans $\mathcal N(A^{\mathsf T})$: the three rows of $A$ satisfy $-\mathbf r_1+\mathbf r_2+\mathbf r_3=\mathbf 0$.

For a wide matrix it is easier to start from $AA^{\mathsf H}$, which is smaller. For
\[
A=\begin{bmatrix}3&2&2\\2&3&-2\end{bmatrix},\qquad AA^{\mathsf T}=\begin{bmatrix}17&8\\8&17\end{bmatrix},
\]
the eigenvalues of $AA^{\mathsf T}$ are $25$ and $9$, with the eigenvectors $\mathbf u_1=(1,1)/\sqrt2$ and $\mathbf u_2=(1,-1)/\sqrt2$, so $\sigma_1=5$ and $\sigma_2=3$. The right singular vectors follow from the second relation of \eqref{eq:singular-pairs}, $\mathbf v_i=A^{\mathsf T}\mathbf u_i/\sigma_i$: $\mathbf v_1=(5,5,0)/(5\sqrt2)=(1,1,0)/\sqrt2$ and $\mathbf v_2=(1,-1,4)/(3\sqrt2)$. A unit vector orthogonal to both is $\mathbf v_3=(2,-2,-1)/3$, and $A\mathbf v_3=\mathbf 0$, as it must be; it spans $\mathcal N(A)$. The outer product expansion reads
\[
A=5\,\mathbf u_1\mathbf v_1^{\mathsf T}+3\,\mathbf u_2\mathbf v_2^{\mathsf T}=\frac52\begin{bmatrix}1&1&0\\1&1&0\end{bmatrix}+\frac12\begin{bmatrix}1&-1&4\\-1&1&-4\end{bmatrix}.
\]

The column $\mathbf a=(1,2)$ of \cref{sec:inverse}, a $2\times1$ matrix, has $\mathbf a^{\mathsf T}\mathbf a=5$, so its only singular value is $\sqrt5=\|\mathbf a\|$, with $\mathbf v_1=[1]$ and $\mathbf u_1=\mathbf a/\sqrt5$; the second left singular vector $\mathbf u_2=(-2,1)/\sqrt5$ spans the left null space. In general a nonzero column $\mathbf a$ has the single singular value $\|\mathbf a\|$, and a nonzero row $\mathbf a^{\mathsf H}$ the same.

\subsection{Singular Values and Eigenvalues}

For a square matrix both eigenvalues and singular values are defined, and they are different numbers. The nilpotent matrix $N$ with the rows $(0,1)$ and $(0,0)$ has the double eigenvalue $0$, but it maps $\mathbf e_2$ to $\mathbf e_1$ and $\mathbf e_1$ to $\mathbf 0$, so $N=I\,\operatorname{diag}(1,0)\,P$ with the permutation matrix $P$ that exchanges the coordinates, and its singular values are $1$ and $0$. The matrix $B$ of \cref{sec:svd-geometry} has the eigenvalues $1$ and $2$ and the singular values $2.288$ and $0.874$. The next proposition collects what can be said in general.

\begin{proposition}[Singular values and eigenvalues]
\label{prop:singular-eigen}
Let $A\in\mathbb F^{m\times n}$ have the singular values $\sigma_1\ge\cdots\ge\sigma_p$, and $\sigma_i\coloneqq0$ for $i>p$.
\begin{enumerate}
\item $\sigma_i^2=\lambda_{n+1-i}(A^{\mathsf H}A)$ for $i=1,\ldots,n$, and $\sigma_i^2=\lambda_{m+1-i}(AA^{\mathsf H})$ for $i=1,\ldots,m$, where $\lambda_1\le\lambda_2\le\cdots$ is the increasing order of \cref{ch:symmetric-matrices}.
\item The matrices $A^{\mathsf H}$, $cA$ with $c\ne0$, and $QAZ$ with unitary $Q\in\mathbb F^{m\times m}$ and $Z\in\mathbb F^{n\times n}$ have the singular values $\sigma_i$, $|c|\sigma_i$ and $\sigma_i$.
\item If $A$ is normal, in particular if $A$ is Hermitian, its singular values are the moduli $|\lambda_j|$ of its eigenvalues, arranged in decreasing order.
\item If $m=n$, every eigenvalue $\lambda$ of $A$ satisfies $\sigma_n\le|\lambda|\le\sigma_1$, and $|\det A|=\sigma_1\sigma_2\cdots\sigma_n$.
\item $\|A\|_{\mathrm F}^2=\operatorname{tr}(A^{\mathsf H}A)=\sigma_1^2+\cdots+\sigma_p^2$.
\end{enumerate}
\end{proposition}

\begin{proof}
(i) By the proof of \cref{prop:svd-unique}, $A^{\mathsf H}A$ has the eigenvalues $\sigma_1^2\ge\cdots\ge\sigma_n^2$, with the convention for $i>p$; the $i$th largest is the $(n+1-i)$th smallest. The same holds for $AA^{\mathsf H}$ with $m$ in place of $n$. (ii) If $A=U\Sigma V^{\mathsf H}$, then $A^{\mathsf H}=V\Sigma^{\mathsf T}U^{\mathsf H}$, $cA=(e^{i\theta}U)(|c|\Sigma)V^{\mathsf H}$ with $c=|c|e^{i\theta}$, and $QAZ=(QU)\Sigma(Z^{\mathsf H}V)^{\mathsf H}$ are singular value decompositions. (iii) By \cref{thm:normal-unitary}, $A=Q\Lambda Q^{\mathsf H}$ with unitary $Q$ and $\Lambda=\operatorname{diag}(\lambda_j)$. Write $\lambda_j=|\lambda_j|e^{i\theta_j}$, with $\theta_j=0$ if $\lambda_j=0$, and $D\coloneqq\operatorname{diag}(e^{i\theta_j})$. Then $A=(QD)\,|\Lambda|\,Q^{\mathsf H}$ with $|\Lambda|\coloneqq\operatorname{diag}(|\lambda_j|)$, and $QD$ is unitary; permuting the diagonal into decreasing order as in the proof of \cref{thm:spectral} gives an SVD. (iv) Let $A\mathbf x=\lambda\mathbf x$ with a unit vector $\mathbf x\in\mathbb C^n$ (for a real $A$, a real SVD is also a complex one). With $\mathbf y\coloneqq V^{\mathsf H}\mathbf x$, a unit vector,
\[
|\lambda|^2=\|A\mathbf x\|^2=\|U\Sigma\mathbf y\|^2=\|\Sigma\mathbf y\|^2=\sum_{i=1}^n\sigma_i^2|y_i|^2 ,
\]
a weighted average of the $\sigma_i^2$, which lies between $\sigma_n^2$ and $\sigma_1^2$. Further $|\det A|=|\det U|\,|\det\Sigma|\,|\det V^{\mathsf H}|=\sigma_1\cdots\sigma_n$ by \cref{thm:det-product,prop:det-unitary}. (v) The trace of $A^{\mathsf H}A$ is the sum of the squared lengths of the columns, and the sum of the eigenvalues (\cref{thm:trace-det-eigen}), which are the $\sigma_i^2$ by (i).
\end{proof}

Part (i) is the relation \eqref{eq:sigma-lambda} announced in the introduction, and the reversal of the index is the price of the two customary orders. Part (iii) shows that for Hermitian matrices nothing new appears: the singular values are the absolute values of the eigenvalues, so that for the second difference matrix $T_n$, which is positive definite, they are its eigenvalues $2-2\cos(j\pi/(n+1))$ (\cref{thm:second-difference-eigen}), in decreasing order. For a matrix that is far from normal, eigenvalues and singular values can differ greatly. The matrix $A_2$ of \cref{sec:matrix-powers}, with the rows $(4/5,2)$ and $(0,4/5)$, has the double eigenvalue $4/5$. Its singular values satisfy $\sigma_1\sigma_2=|\det A_2|=16/25$ and $\sigma_1^2+\sigma_2^2=\|A_2\|_{\mathrm F}^2=4+32/25$, so $(\sigma_1-\sigma_2)^2=4$ and $(\sigma_1+\sigma_2)^2=4+64/25$, which gives
\[
\sigma_1=1+\frac{\sqrt{41}}5=2.2806\ldots,\qquad\sigma_2=\frac{\sqrt{41}}5-1=0.2806\ldots;
\]
the eigenvalue $4/5$ lies between them, as (iv) requires, but far from both. \Cref{ch:eigenvalues} used only $\|A_2\|_2>2$, which is now the exact value $\|A_2\|_2=\sigma_1$ of \cref{sec:svd-norms}.

\begin{note}{Weyl's inequalities between eigenvalues and singular values}
Part (iv) of \cref{prop:singular-eigen} has a sharper form. If the eigenvalues of $A\in\mathbb C^{n\times n}$ are ordered by decreasing modulus, $|\lambda_1|\ge\cdots\ge|\lambda_n|$, then
\[
|\lambda_1\lambda_2\cdots\lambda_k|\le\sigma_1\sigma_2\cdots\sigma_k\qquad(k=1,\ldots,n),
\]
with equality for $k=n$, where both sides are $|\det A|$. Weyl proved these inequalities in 1949~\cite{weyl1949}, and A.~Horn showed in 1954 that they are the only relations between the two sets of numbers: any numbers that satisfy them are the moduli of the eigenvalues and the singular values of some matrix~\cite{horn1954}. The case $k=1$ is $\rho(A)\le\|A\|_2$ of \cref{prop:spectral-radius-norm}, and the general case follows by applying that inequality to a larger matrix built from $A$, the $k$th compound matrix, whose entries are the $k\times k$ minors of $A$ and whose eigenvalues and singular values are the products of $k$ eigenvalues and of $k$ singular values of $A$~\cite[Chap.~3]{horn1991topics}.
\end{note}

\section{The Four Subspaces Again}
\label{sec:svd-subspaces}

\subsection{Orthonormal Bases of the Four Subspaces}

\Cref{thm:four-bases} found bases of the four fundamental subspaces by elimination, and \cref{thm:four-subspaces-orthogonal} showed that they form two pairs of orthogonal complements. The singular value decomposition provides orthonormal bases of all four at once, and it pairs them: $A$ maps the basis of the row space onto the basis of the column space, vector by vector, up to the factors $\sigma_i$.

\begin{theorem}[The four subspaces and the singular value decomposition]
\label{thm:svd-subspaces}
Let $A=U\Sigma V^{\mathsf H}\in\mathbb F^{m\times n}$ be a singular value decomposition, and let $r$ be the number of positive singular values. Then $r=\operatorname{rank}A$, and
\begin{enumerate}
\item $\mathbf u_1,\ldots,\mathbf u_r$ is an orthonormal basis of the column space $\mathcal C(A)$;
\item $\mathbf u_{r+1},\ldots,\mathbf u_m$ is an orthonormal basis of the left null space $\mathcal N(A^{\mathsf H})$;
\item $\mathbf v_1,\ldots,\mathbf v_r$ is an orthonormal basis of the row space $\mathcal C(A^{\mathsf H})$;
\item $\mathbf v_{r+1},\ldots,\mathbf v_n$ is an orthonormal basis of the null space $\mathcal N(A)$.
\end{enumerate}
Consequently the orthogonal projections onto the four subspaces are
\[
P_{\mathcal C(A)}=U_rU_r^{\mathsf H},\qquad P_{\mathcal N(A^{\mathsf H})}=I-U_rU_r^{\mathsf H},\qquad P_{\mathcal C(A^{\mathsf H})}=V_rV_r^{\mathsf H},\qquad P_{\mathcal N(A)}=I-V_rV_r^{\mathsf H}.
\]
\end{theorem}

\begin{proof}
Multiplying by invertible matrices does not change the rank (\cref{prop:rank-product}), so $\operatorname{rank}A=\operatorname{rank}\Sigma$, and the rank of a diagonal matrix is the number of its nonzero diagonal entries, which is $r$. By \eqref{eq:singular-pairs}, $\mathbf u_i=A\mathbf v_i/\sigma_i\in\mathcal C(A)$ for $i\le r$, and these are $r$ orthonormal, hence independent, vectors in a space of dimension $r$; so they form a basis (\cref{thm:basis-extension}). For $i>r$, $A\mathbf v_i=\mathbf 0$, either by \eqref{eq:singular-pairs} with $\sigma_i=0$ or because $i>p$, so $\mathbf v_{r+1},\ldots,\mathbf v_n$ are $n-r$ orthonormal vectors of $\mathcal N(A)$, whose dimension is $n-r$ by the rank--nullity theorem; they form a basis. Applying these two arguments to $A^{\mathsf H}=V\Sigma^{\mathsf T}U^{\mathsf H}$, which has the same $r$ positive singular values, gives (iii) and (ii). The projections follow from \cref{thm:projection-matrix}\,(ii), and from \cref{prop:projection-properties}\,(v) for the complements.
\end{proof}

\Cref{fig:svd-subspaces} shows the theorem. It is the picture of \cref{fig:four-subspaces-orthogonal} with orthonormal bases drawn in, and it says that a matrix, viewed through its SVD, does three simple things: it annihilates the null space, it maps the orthonormal basis $\mathbf v_1,\ldots,\mathbf v_r$ of the row space onto the orthogonal vectors $\sigma_1\mathbf u_1,\ldots,\sigma_r\mathbf u_r$ of the column space, and it misses the left null space entirely. The conjugate transpose runs the same arrows backwards, $A^{\mathsf H}\mathbf u_i=\sigma_i\mathbf v_i$, and annihilates the left null space. Restricted to the row space, $A$ is an invertible map onto the column space, and its inverse there is the map $\mathbf u_i\mapsto\mathbf v_i/\sigma_i$; \cref{sec:pseudoinverse} calls it the pseudoinverse.

\begin{figure}
\sdfloat{%
\begin{tikzpicture}[font=\scriptsize, x=1cm, y=0.62cm]
\draw[line width=0.5pt, rounded corners=3pt] (-0.75,-0.6) rectangle (0.75,6.6);
\draw[line width=0.5pt, rounded corners=3pt] (4.25,-0.6) rectangle (5.75,6.6);
\draw[black!40, line width=0.4pt] (-0.75,3.0) -- (0.75,3.0);
\draw[black!40, line width=0.4pt] (4.25,3.0) -- (5.75,3.0);
\node at (0,7.1) {$\mathbb F^n$};
\node at (5,7.1) {$\mathbb F^m$};
\foreach \y/\l in {6/{\mathbf v_1}, 5/{\vdots}, 4/{\mathbf v_r}, 2.2/{\mathbf v_{r+1}}, 1.2/{\vdots}, 0.2/{\mathbf v_n}} {\node at (0,\y) {$\l$};}
\foreach \y/\l in {6/{\mathbf u_1}, 5/{\vdots}, 4/{\mathbf u_r}, 2.2/{\mathbf u_{r+1}}, 1.2/{\vdots}, 0.2/{\mathbf u_m}} {\node at (5,\y) {$\l$};}
\draw[line width=0.6pt, -{Latex[length=1.6mm]}] (0.45,6.15) -- (4.55,6.15);
\draw[line width=0.6pt, -{Latex[length=1.6mm]}] (0.45,4.15) -- (4.55,4.15);
\node[above] at (2.5,6.15) {$A\mathbf v_1=\sigma_1\mathbf u_1$};
\node[above] at (2.5,4.15) {$A\mathbf v_r=\sigma_r\mathbf u_r$};
\draw[line width=0.5pt, dashed, -{Latex[length=1.6mm]}] (4.55,5.75) -- (0.45,5.75);
\draw[line width=0.5pt, dashed, -{Latex[length=1.6mm]}] (4.55,3.75) -- (0.45,3.75);
\node[below] at (2.5,5.75) {$A^{\mathsf H}\mathbf u_1=\sigma_1\mathbf v_1$};
\node[below] at (2.5,3.75) {$A^{\mathsf H}\mathbf u_r=\sigma_r\mathbf v_r$};
\node (z1) at (1.9,1.2) {$\mathbf 0$};
\node (z2) at (3.1,1.2) {$\mathbf 0$};
\draw[line width=0.6pt, -{Latex[length=1.6mm]}] (0.5,2.2) -- (z1);
\draw[line width=0.6pt, -{Latex[length=1.6mm]}] (0.5,0.2) -- (z1);
\node[above left, inner sep=1pt] at (1.25,1.75) {$A$};
\draw[line width=0.5pt, dashed, -{Latex[length=1.6mm]}] (4.5,2.2) -- (z2);
\draw[line width=0.5pt, dashed, -{Latex[length=1.6mm]}] (4.5,0.2) -- (z2);
\node[above right, inner sep=1pt] at (3.75,1.75) {$A^{\mathsf H}$};
\draw[decorate, decoration={brace, amplitude=3pt}, line width=0.4pt] (-0.9,3.5) -- (-0.9,6.5) node[midway, left=3pt, align=right] {$\mathcal C(A^{\mathsf H})$\\$\dim r$};
\draw[decorate, decoration={brace, amplitude=3pt}, line width=0.4pt] (-0.9,-0.5) -- (-0.9,2.5) node[midway, left=3pt, align=right] {$\mathcal N(A)$\\$\dim n-r$};
\draw[decorate, decoration={brace, amplitude=3pt, mirror}, line width=0.4pt] (5.9,3.5) -- (5.9,6.5) node[midway, right=3pt, align=left] {$\mathcal C(A)$\\$\dim r$};
\draw[decorate, decoration={brace, amplitude=3pt, mirror}, line width=0.4pt] (5.9,-0.5) -- (5.9,2.5) node[midway, right=3pt, align=left] {$\mathcal N(A^{\mathsf H})$\\$\dim m-r$};
\end{tikzpicture}}{The singular vectors as orthonormal bases of the four fundamental subspaces (\cref{thm:svd-subspaces}). The matrix $A$ (solid arrows) maps $\mathbf v_i$ to $\sigma_i\mathbf u_i$ for $i\le r$ and the basis of the null space to $\mathbf 0$; its conjugate transpose (dashed) maps $\mathbf u_i$ to $\sigma_i\mathbf v_i$ and the basis of the left null space to $\mathbf 0$. Compare the rank normal form of \cref{fig:rank-normal-form}, where the bases need not be orthonormal and all the factors $\sigma_i$ are $1$.\label{fig:svd-subspaces}}
\end{figure}

For the matrix of \cref{sec:four-subspaces}, with the rows $(1,2,0,3)$, $(2,4,1,8)$ and $(-1,-2,1,-1)$, of rank $2$, the singular values are $10.12\ldots$, $1.912\ldots$ and $0$, as \cref{sec:svd-computation} confirms. The SVD replaces the bases of \cref{thm:four-bases}, whose vectors were read off from the reduced echelon form and were neither orthogonal nor normalised, by orthonormal ones; it cannot be computed by finitely many rational operations, since already a $2\times2$ matrix with integer entries can have irrational singular values, as $B$ of \cref{sec:svd-geometry} illustrates, but its singular values change little when the entries are perturbed (\cref{sec:svd-perturbation}).

\subsection{Unitary Equivalence}

\Cref{thm:rank-normal-form} showed that two matrices are equivalent, $B=MAN$ with invertible $M$ and $N$, if and only if they have the same rank, and \cref{ch:linear-maps} remarked that the bases of the rank normal form can be refined to orthonormal ones at the price of keeping the singular values. With orthonormal bases on both sides the rank is no longer the only invariant.

\begin{theorem}[Unitary equivalence]
\label{thm:unitary-equivalence}
\leavevmode
\begin{enumerate}
\item Two matrices $A,B\in\mathbb F^{m\times n}$ satisfy $B=QAZ$ for some unitary $Q\in\mathbb F^{m\times m}$ and $Z\in\mathbb F^{n\times n}$ if and only if they have the same singular values.
\item Let $T\colon V\to W$ be a linear map between finite-dimensional inner product spaces over $\mathbb F$ of dimensions $n$ and $m$. There are orthonormal bases $\mathcal B$ of $V$ and $\mathcal D$ of $W$ such that $[T]_{\mathcal D\leftarrow\mathcal B}$ is a diagonal matrix with the diagonal entries $\sigma_1\ge\cdots\ge\sigma_p\ge0$; these numbers do not depend on the choice of the bases, and exactly $\operatorname{rank}T$ of them are positive.
\end{enumerate}
\end{theorem}

\begin{proof}
(i) If $B=QAZ$, the singular values agree by \cref{prop:singular-eigen}\,(ii). Conversely, if $A=U_A\Sigma V_A^{\mathsf H}$ and $B=U_B\Sigma V_B^{\mathsf H}$ with the same $\Sigma$, then $B=U_BU_A^{\mathsf H}AV_AV_B^{\mathsf H}$, a product of the required form. (ii) Choose orthonormal bases $\mathcal B_0$ of $V$ and $\mathcal D_0$ of $W$ (\cref{cor:orthonormal-basis}) and let $A\coloneqq[T]_{\mathcal D_0\leftarrow\mathcal B_0}=U\Sigma V^{\mathsf H}$. Let $\mathcal B$ consist of the vectors of $V$ whose coordinate vectors in $\mathcal B_0$ are the columns of $V$, and $\mathcal D$ likewise with $U$ and $\mathcal D_0$. They are orthonormal, because coordinates in an orthonormal basis preserve inner products (\cref{thm:orthonormal-expansion}), and by \cref{thm:change-of-basis-map} $[T]_{\mathcal D\leftarrow\mathcal B}=U^{-1}AV=\Sigma$. If $[T]_{\mathcal D'\leftarrow\mathcal B'}=\Sigma'$ for other orthonormal bases, the same theorem gives $\Sigma'=Q\Sigma Z$ with the unitary change-of-basis matrices, and (i) shows $\Sigma'=\Sigma$. The rank statement is \cref{thm:svd-subspaces} together with \cref{prop:rank-of-matrix}.
\end{proof}

The singular values are therefore a complete set of invariants for \term{unitary equivalence}, just as the rank is for equivalence. Part (ii) is the statement promised after \cref{thm:rank-normal-form}: with orthonormal bases one cannot reduce a map to $N_r$, which would mean replacing every $\sigma_i$ by $1$, but one can reduce it to the diagonal matrix of its singular values, and the remaining numbers measure how much the map stretches in each of the $r$ directions it does not annihilate. A map between inner product spaces is thus determined, up to isometries of the domain and of the codomain, by $p$ numbers. Since the singular values of $T$ can be computed in any pair of orthonormal bases, they belong to the map, and we write $\sigma_i(T)$.

\section{Norms}
\label{sec:svd-norms}

\subsection{The Spectral Norm}

\begin{theorem}[Norms and singular values]
\label{thm:two-norm}
Let $A\in\mathbb F^{m\times n}$ have the singular values $\sigma_1\ge\cdots\ge\sigma_p$ and the rank $r$.
\begin{enumerate}
\item $\|A\|_2=\max_{\|\mathbf x\|=1}\|A\mathbf x\|=\sigma_1$, attained at $\mathbf x=\mathbf v_1$; equivalently, $\|A\|_2^2$ is the largest eigenvalue of $A^{\mathsf H}A$.
\item If $m\ge n$, then $\min_{\|\mathbf x\|=1}\|A\mathbf x\|=\sigma_n$, attained at $\mathbf x=\mathbf v_n$; if $m<n$, the minimum is $0$. In particular $\|A\mathbf x\|\ge\sigma_n\|\mathbf x\|$ for all $\mathbf x$ when $m\ge n$, and $\sigma_n>0$ if and only if $A$ has full column rank.
\item $\|A\|_{\mathrm F}=(\sigma_1^2+\cdots+\sigma_r^2)^{1/2}$.
\item Both norms are unitarily invariant, $\|QAZ\|=\|A\|$ for unitary $Q$ and $Z$, and $\|A^{\mathsf H}\|_2=\|A\|_2$, $\|A^{\mathsf H}A\|_2=\|AA^{\mathsf H}\|_2=\|A\|_2^2$.
\item $\|A\|_2\le\|A\|_{\mathrm F}\le\sqrt r\,\|A\|_2$.
\end{enumerate}
\end{theorem}

\begin{proof}
For $\mathbf x\in\mathbb F^n$ let $\mathbf y\coloneqq V^{\mathsf H}\mathbf x$, so that $\|\mathbf y\|=\|\mathbf x\|$ and $\mathbf x$ runs through the unit vectors when $\mathbf y$ does. Then
\[
\|A\mathbf x\|^2=\|U\Sigma\mathbf y\|^2=\|\Sigma\mathbf y\|^2=\sum_{i=1}^p\sigma_i^2|y_i|^2 .
\]
(i) The sum is at most $\sigma_1^2\sum_i|y_i|^2=\sigma_1^2\|\mathbf x\|^2$, with equality for $\mathbf y=\mathbf e_1$, that is, $\mathbf x=\mathbf v_1$. The largest eigenvalue of $A^{\mathsf H}A$ is $\sigma_1^2$ by \cref{prop:singular-eigen}\,(i). (ii) If $m\ge n$, then $p=n$, and the sum is at least $\sigma_n^2\|\mathbf y\|^2$, with equality for $\mathbf y=\mathbf e_n$. If $m<n$, then $\mathbf v_n\in\mathcal N(A)$ by \cref{thm:svd-subspaces}, since $r\le m<n$. By the same theorem, $\sigma_n>0$ means $r=n$. (iii) is \cref{prop:singular-eigen}\,(v). (iv) $QAZ$ has the same singular values as $A$ by \cref{prop:singular-eigen}\,(ii), so the same norms by (i) and (iii); the same holds for $A^{\mathsf H}$, and $A^{\mathsf H}A=V(\Sigma^{\mathsf T}\Sigma)V^{\mathsf H}$ has the singular values $\sigma_i^2$ by \cref{prop:singular-eigen}\,(iii), since it is Hermitian with these eigenvalues. (v) $\sigma_1^2\le\sum_{i\le r}\sigma_i^2\le r\sigma_1^2$.
\end{proof}

The operator norm $\|A\|_2$ is also called the \term{spectral norm}. The theorem fulfils the promises of \cref{ch:several-variables,ch:vectors-matrices}: the supremum that defined $\|A\|_2$ is attained, and it is the largest singular value, the length of the longest semi-axis of the image of the unit sphere. For the matrix with the rows $(1,-2)$ and $(3,4)$ of \cref{sec:trace-norms},
\[
A^{\mathsf T}A=\begin{bmatrix}10&10\\10&20\end{bmatrix},
\]
with the trace $30$ and the determinant $100$, whose eigenvalues are $15\pm\sqrt{125}=15\pm5\sqrt5$; so
\[
\|A\|_2=\sigma_1=(15+5\sqrt5)^{1/2}=5.1166\ldots,\qquad\sigma_2=(15-5\sqrt5)^{1/2}=1.9543\ldots,
\]
which is the value stated there. As checks, $\sigma_1^2+\sigma_2^2=30=\|A\|_{\mathrm F}^2$ and $\sigma_1\sigma_2=\sqrt{225-125}=10=|\det A|$, as \cref{prop:singular-eigen} requires. The theorem also completes \cref{prop:hermitian-norm}, which found $\|A\|_2=\rho(A)$ for Hermitian and normal matrices: for these matrices the singular values are the moduli of the eigenvalues (\cref{prop:singular-eigen}\,(iii)). For every matrix, $\|A\|_2^2$ is the largest eigenvalue of $A^{\mathsf H}A$, as \cref{ch:symmetric-matrices} announced, and $\rho(A)\le\|A\|_2$ may be far from an equality, as the matrix $A_2$ with $\rho(A_2)=4/5$ and $\|A_2\|_2=2.28$ shows.

Part (ii) gives the shortest image of a unit vector, which \cref{ch:inner-products} used to define the condition number provisionally; \cref{sec:svd-condition} takes this up. Part (v) is sharp on both sides: a matrix of rank one has $\|A\|_{\mathrm F}=\|A\|_2$, and a matrix whose $r$ positive singular values are equal, such as a matrix with orthonormal columns, has $\|A\|_{\mathrm F}=\sqrt r\|A\|_2$; for $A=I_n$ the two norms are $\sqrt n$ and $1$.

\subsection{Comparison with Other Norms}

The spectral norm is not given by a formula in the entries, but it is bounded in terms of the norms $\|\cdot\|_1$ and $\|\cdot\|_\infty$ of \cref{thm:norm-formulas}, which are. That theorem was proved for real matrices; for complex ones the same proof works with $x_j\coloneqq\bar a_{i^*j}/|a_{i^*j}|$ in place of $\operatorname{sign}a_{i^*j}$, and we use it in this form.

\begin{proposition}[The spectral norm and the norms $\|\cdot\|_1$, $\|\cdot\|_\infty$]
\label{prop:norm-comparison}
For $A\in\mathbb F^{m\times n}$,
\begin{enumerate}
\item $\|A\|_2^2\le\|A\|_1\|A\|_\infty$;
\item $\|A\|_\infty/\sqrt n\le\|A\|_2\le\sqrt m\,\|A\|_\infty$;
\item $\|A\|_1/\sqrt m\le\|A\|_2\le\sqrt n\,\|A\|_1$.
\end{enumerate}
\end{proposition}

\begin{proof}
(i) $\|A\|_2^2$ is the largest eigenvalue of $A^{\mathsf H}A$, so by \cref{prop:spectral-radius-norm} and \cref{prop:operator-norms}\,(iii) it is at most $\|A^{\mathsf H}A\|_1\le\|A^{\mathsf H}\|_1\|A\|_1$, and $\|A^{\mathsf H}\|_1=\|A\|_\infty$ because the columns of $A^{\mathsf H}$ are the conjugated rows of $A$. (ii) and (iii) follow from the comparison of vector norms in \cref{prop:standard-norms}, whose proof holds for complex vectors. For every $\mathbf x$,
\begin{gather*}
\|A\mathbf x\|_2\le\sqrt m\,\|A\mathbf x\|_\infty\le\sqrt m\,\|A\|_\infty\|\mathbf x\|_\infty\le\sqrt m\,\|A\|_\infty\|\mathbf x\|_2,\\
\|A\mathbf x\|_\infty\le\|A\mathbf x\|_2\le\|A\|_2\|\mathbf x\|_2\le\sqrt n\,\|A\|_2\|\mathbf x\|_\infty ,
\end{gather*}
which gives (ii), and in the same way
\begin{gather*}
\|A\mathbf x\|_2\le\|A\mathbf x\|_1\le\|A\|_1\|\mathbf x\|_1\le\sqrt n\,\|A\|_1\|\mathbf x\|_2,\\
\|A\mathbf x\|_1\le\sqrt m\,\|A\mathbf x\|_2\le\sqrt m\,\|A\|_2\|\mathbf x\|_2\le\sqrt m\,\|A\|_2\|\mathbf x\|_1
\end{gather*}
gives (iii).
\end{proof}

For the matrix with the rows $(1,-2)$ and $(3,4)$, $\|A\|_1\|A\|_\infty=6\cdot7=42$, against $\|A\|_2^2=26.18\ldots$. Bound (i) is often used as a cheap upper estimate; it is an equality for diagonal matrices and, more generally, whenever $|A|$ has a single nonzero entry in every row and column. Estimating $\|A\|_2$ or $\|A^{-1}\|_2$ more precisely without computing the SVD is the subject of Volume~VIII.

\subsection{The Nuclear Norm}

The two norms of \cref{thm:two-norm} are the largest singular value and the Euclidean length of the vector of singular values. The sum of the singular values,
\[
\|A\|_*\coloneqq\sigma_1+\sigma_2+\cdots+\sigma_p ,
\]
is a third norm, the \term{nuclear norm} or \term{trace norm}; for a positive semidefinite matrix it is the trace. That it satisfies the triangle inequality is not obvious from the definition, and the note below indicates the proof. For the matrix with the rows $(1,-2)$ and $(3,4)$ it is $\sigma_1+\sigma_2=(\sigma_1^2+\sigma_2^2+2\sigma_1\sigma_2)^{1/2}=\sqrt{30+20}=\sqrt{50}$. The nuclear norm plays for the rank the part that the norm $\|\cdot\|_1$ plays for the number of nonzero entries of a vector: minimising it under linear constraints tends to produce matrices of low rank, and under suitable conditions this recovers a matrix of low rank exactly from a small fraction of its entries~\cite{candes2009}.

\begin{note}{Unitarily invariant norms}
Every norm that depends only on the singular values is unitarily invariant, and von Neumann showed that the converse holds: the unitarily invariant norms on $\mathbb F^{m\times n}$ are exactly the functions $g(\sigma_1,\ldots,\sigma_p)$ for the norms $g$ on $\mathbb R^p$ that are invariant under permutations and sign changes of the coordinates~\cite{vonneumann1937},~\cite[Sect.~7.4]{horn2013}. The \emph{Schatten norms} $(\sum_i\sigma_i^q)^{1/q}$, $1\le q\le\infty$, include the nuclear norm ($q=1$), the Frobenius norm ($q=2$) and the spectral norm ($q=\infty$), and the \emph{Ky Fan norms} $\sigma_1+\cdots+\sigma_k$ interpolate between the spectral and the nuclear norm~\cite{fan1951}. The triangle inequality for the nuclear norm follows from the duality
\[
\|A\|_*=\max\bigl\{\operatorname{Re}\operatorname{tr}(B^{\mathsf H}A):\|B\|_2\le1\bigr\},
\]
since a maximum of linear functions is subadditive. To see the duality, write $A=\sum_i\sigma_i\mathbf u_i\mathbf v_i^{\mathsf H}$; then $\operatorname{tr}(B^{\mathsf H}A)=\sum_i\sigma_i\,\mathbf v_i^{\mathsf H}B^{\mathsf H}\mathbf u_i=\sum_i\sigma_i\langle B\mathbf v_i,\mathbf u_i\rangle$, whose real part is at most $\sum_i\sigma_i$ when $\|B\|_2\le1$, by the Cauchy--Schwarz inequality, and equal to it for $B=U_rV_r^{\mathsf H}$. In the same sense the spectral norm and the nuclear norm are dual to each other, as $\|\cdot\|_\infty$ and $\|\cdot\|_1$ are on vectors.
\end{note}

\section{Low-Rank Approximation}
\label{sec:low-rank}

\subsection{The Truncated Decomposition}

The outer product expansion \eqref{eq:svd-outer} orders the rank-one pieces of a matrix by size. Keeping the first $k$ of them gives the \term{truncated SVD}
\[
A_k\coloneqq\sum_{i=1}^k\sigma_i\,\mathbf u_i\mathbf v_i^{\mathsf H}=U_k\Sigma_kV_k^{\mathsf H}\qquad(0\le k\le r),
\]
with $A_0=0$ and $A_r=A$. Its rank is $k$, by \cref{thm:svd-subspaces} applied to the decomposition $A_k=U\operatorname{diag}(\sigma_1,\ldots,\sigma_k,0,\ldots,0)V^{\mathsf H}$, and the remainder $A-A_k=\sum_{i>k}\sigma_i\mathbf u_i\mathbf v_i^{\mathsf H}$ has the singular values $\sigma_{k+1},\ldots,\sigma_r$ and zeros, so by \cref{thm:two-norm}
\begin{equation}
\|A-A_k\|_2=\sigma_{k+1},\qquad\|A-A_k\|_{\mathrm F}=\bigl(\sigma_{k+1}^2+\cdots+\sigma_r^2\bigr)^{1/2},
\label{eq:truncation-error}
\end{equation}
with $\sigma_{r+1}\coloneqq0$. The theorem of this section says that no matrix of rank $k$ does better. Its proof in the Frobenius norm rests on a maximum principle for Hermitian matrices, which also underlies principal component analysis.

\begin{lemma}[Ky Fan's maximum principle]
\label{lem:ky-fan}
Let $M\in\mathbb F^{n\times n}$ be Hermitian with the eigenvalues $\lambda_1\le\cdots\le\lambda_n$ and orthonormal eigenvectors $\mathbf q_1,\ldots,\mathbf q_n$, and let $1\le k\le n$. For every $Q\in\mathbb F^{n\times k}$ with orthonormal columns,
\[
\operatorname{tr}(Q^{\mathsf H}MQ)\le\lambda_n+\lambda_{n-1}+\cdots+\lambda_{n-k+1},
\]
with equality if the columns of $Q$ span $\operatorname{span}\{\mathbf q_{n-k+1},\ldots,\mathbf q_n\}$. If $k<n$ and $\lambda_{n-k}<\lambda_{n-k+1}$, equality holds only if the columns of $Q$ span this subspace.
\end{lemma}

\begin{proof}
By \cref{prop:spectral-decomposition}, $M=\sum_i\lambda_i\mathbf q_i\mathbf q_i^{\mathsf H}$, so
\[
\operatorname{tr}(Q^{\mathsf H}MQ)=\sum_{i=1}^n\lambda_i\operatorname{tr}(Q^{\mathsf H}\mathbf q_i\mathbf q_i^{\mathsf H}Q)=\sum_{i=1}^n\lambda_iw_i,\qquad w_i\coloneqq\|Q^{\mathsf H}\mathbf q_i\|^2,
\]
by \cref{prop:trace}. Each $w_i$ lies in $[0,1]$, because $\|Q^{\mathsf H}\mathbf y\|=\|QQ^{\mathsf H}\mathbf y\|\le\|\mathbf y\|$ for every $\mathbf y$, $QQ^{\mathsf H}$ being an orthogonal projection (\cref{thm:projection-matrix,prop:projection-properties}). Their sum is $k$: with the unitary matrix $W$ whose columns are the $\mathbf q_i$, $\sum_iw_i=\|Q^{\mathsf H}W\|_{\mathrm F}^2=\|Q^{\mathsf H}\|_{\mathrm F}^2=\operatorname{tr}(Q^{\mathsf H}Q)=k$ by \cref{thm:two-norm}\,(iv). Now split the sum at $n-k$. For $i>n-k$, $w_i-1\le0$ and $\lambda_i\ge\lambda_{n-k+1}$; for $i\le n-k$, $w_i\ge0$ and $\lambda_i\le\lambda_{n-k}$, where for $k=n$ the second group is empty. Hence
\[
\begin{aligned}
\sum_{i=1}^n\lambda_iw_i-\sum_{i>n-k}\lambda_i&=\sum_{i>n-k}\lambda_i(w_i-1)+\sum_{i\le n-k}\lambda_iw_i\\
&\le\lambda_{n-k+1}\sum_{i>n-k}(w_i-1)+\lambda_{n-k}\sum_{i\le n-k}w_i ,
\end{aligned}
\]
and since $\sum_{i>n-k}(w_i-1)=k-\sum_{i\le n-k}w_i-k=-\sum_{i\le n-k}w_i$, the right side equals $-(\lambda_{n-k+1}-\lambda_{n-k})\sum_{i\le n-k}w_i\le0$. If the columns of $Q$ span the subspace of the statement, then $w_i=1$ for $i>n-k$ and $w_i=0$ otherwise, and equality holds. If $\lambda_{n-k}<\lambda_{n-k+1}$ and equality holds, then $\sum_{i\le n-k}w_i=0$, so $Q^{\mathsf H}\mathbf q_i=\mathbf 0$ for $i\le n-k$: the columns of $Q$ are orthogonal to $\mathbf q_1,\ldots,\mathbf q_{n-k}$, hence lie in the span of the remaining $k$ eigenvectors, and $k$ independent vectors of a $k$-dimensional space span it.
\end{proof}

For $k=1$ the lemma is the maximum part of \cref{thm:rayleigh}: $\mathbf q^{\mathsf H}M\mathbf q\le\lambda_n$ for unit vectors $\mathbf q$. In general it says that the sum of the $k$ largest eigenvalues is the largest total of the form over $k$ orthonormal directions. Ky Fan proved it in 1949~\cite{fan1949}.

\subsection{The Theorem of Eckart and Young}

\begin{theorem}[Best approximation of lower rank]
\label{thm:eckart-young}
Let $A\in\mathbb F^{m\times n}$ have the singular value decomposition $A=\sum_{i=1}^r\sigma_i\mathbf u_i\mathbf v_i^{\mathsf H}$, and let $0\le k<p$. For every $B\in\mathbb F^{m\times n}$ with $\operatorname{rank}B\le k$,
\[
\|A-B\|_2\ge\sigma_{k+1},\qquad\|A-B\|_{\mathrm F}\ge\bigl(\sigma_{k+1}^2+\cdots+\sigma_p^2\bigr)^{1/2},
\]
and the truncated SVD $A_k$ attains both bounds; for $k\ge r$ the bounds are $0$ and $A_k=A$. In the Frobenius norm, if $1\le k<r$ and $\sigma_k>\sigma_{k+1}$, then $A_k$ is the only matrix of rank at most $k$ that attains the bound; if $\sigma_k=\sigma_{k+1}$, it is not. In the spectral norm, for $1\le k<r$, the bound is attained by infinitely many matrices of rank at most $k$.
\end{theorem}

\begin{proof}
For $k\ge r$ there is nothing to prove, so let $k<r$; that $A_k$ attains the bounds is \eqref{eq:truncation-error}.

The spectral norm. Since $\operatorname{rank}B\le k$, the null space $\mathcal N(B)$ has dimension at least $n-k$ by the rank--nullity theorem, and $S\coloneqq\operatorname{span}\{\mathbf v_1,\ldots,\mathbf v_{k+1}\}$ has dimension $k+1$. By the dimension formula (\cref{thm:dimension-formula}), $\dim(\mathcal N(B)\cap S)\ge(n-k)+(k+1)-n=1$, so there is a unit vector $\mathbf x=\sum_{i\le k+1}c_i\mathbf v_i$ with $B\mathbf x=\mathbf 0$ and $\sum_i|c_i|^2=1$. Then
\[
\|(A-B)\mathbf x\|^2=\|A\mathbf x\|^2=\Bigl\|\sum_{i=1}^{k+1}c_i\sigma_i\mathbf u_i\Bigr\|^2=\sum_{i=1}^{k+1}|c_i|^2\sigma_i^2\ge\sigma_{k+1}^2 ,
\]
so $\|A-B\|_2\ge\sigma_{k+1}$.

The Frobenius norm. The column space of $B$ has dimension at most $k\le m$, so it lies in some subspace $S'\subseteq\mathbb F^m$ of dimension $k$ (\cref{thm:basis-extension}). Let $Q\in\mathbb F^{m\times k}$ have orthonormal columns spanning $S'$, and $P\coloneqq QQ^{\mathsf H}$, the orthogonal projection onto $S'$. For each column, $\mathbf b_j\in S'$, so $\|\mathbf a_j-\mathbf b_j\|\ge\|\mathbf a_j-P\mathbf a_j\|$ by \cref{thm:best-approximation}, with equality only for $\mathbf b_j=P\mathbf a_j$. Summing the squares over the columns and using Pythagoras' theorem $\|\mathbf a_j\|^2=\|P\mathbf a_j\|^2+\|\mathbf a_j-P\mathbf a_j\|^2$,
\[
\|A-B\|_{\mathrm F}^2\ge\|A-PA\|_{\mathrm F}^2=\|A\|_{\mathrm F}^2-\|PA\|_{\mathrm F}^2 .
\]
Since $Q$ preserves lengths, $\|PA\|_{\mathrm F}^2=\|Q^{\mathsf H}A\|_{\mathrm F}^2=\operatorname{tr}(Q^{\mathsf H}AA^{\mathsf H}Q)$, and by \cref{lem:ky-fan} applied to $M=AA^{\mathsf H}$, whose $k$ largest eigenvalues are $\sigma_1^2,\ldots,\sigma_k^2$ (\cref{prop:singular-eigen}\,(i)), this is at most $\sigma_1^2+\cdots+\sigma_k^2$. With $\|A\|_{\mathrm F}^2=\sum_{i\le r}\sigma_i^2$ the bound follows.

Uniqueness. If $\|A-B\|_{\mathrm F}$ attains the bound, both inequalities in the last paragraph are equalities: $B=PA$, and $\operatorname{tr}(Q^{\mathsf H}AA^{\mathsf H}Q)=\sigma_1^2+\cdots+\sigma_k^2$. The vectors $\mathbf u_1,\ldots,\mathbf u_m$ are orthonormal eigenvectors of $AA^{\mathsf H}$, with $\mathbf u_1,\ldots,\mathbf u_k$ for the $k$ largest eigenvalues. If $\sigma_k>\sigma_{k+1}$, \cref{lem:ky-fan} shows that $S'=\operatorname{span}\{\mathbf u_1,\ldots,\mathbf u_k\}$, so $P=U_kU_k^{\mathsf H}$ and
\[
B=U_kU_k^{\mathsf H}\sum_{j=1}^r\sigma_j\mathbf u_j\mathbf v_j^{\mathsf H}=\sum_{i=1}^k\sigma_i\mathbf u_i\mathbf v_i^{\mathsf H}=A_k .
\]
If $\sigma_k=\sigma_{k+1}$, both are positive, since $k+1\le r$, and the matrix $B'\coloneqq A_k-\sigma_k\mathbf u_k\mathbf v_k^{\mathsf H}+\sigma_{k+1}\mathbf u_{k+1}\mathbf v_{k+1}^{\mathsf H}$, obtained by exchanging the roles of the $k$th and the $(k+1)$st singular pairs, has rank $k$, differs from $A_k$, and has $\|A-B'\|_{\mathrm F}^2=\sigma_k^2+\sum_{i\ge k+2}\sigma_i^2=\sum_{i\ge k+1}\sigma_i^2$. In the spectral norm, for $|t|\le1$, the matrix $A_k-t\sigma_{k+1}\mathbf u_1\mathbf v_1^{\mathsf H}$ has rank at most $k$ and leaves the remainder $t\sigma_{k+1}\mathbf u_1\mathbf v_1^{\mathsf H}+\sum_{i>k}\sigma_i\mathbf u_i\mathbf v_i^{\mathsf H}$, an outer product expansion with orthonormal vectors and the coefficients $|t|\sigma_{k+1},\sigma_{k+1},\ldots,\sigma_r$ up to factors of modulus one, whose spectral norm is $\sigma_{k+1}$.
\end{proof}

In words: the distance of $A$ from the matrices of rank at most $k$ is $\sigma_{k+1}$ in the spectral norm and $(\sum_{i>k}\sigma_i^2)^{1/2}$ in the Frobenius norm, and the truncated SVD is a nearest matrix in both. The Frobenius statement is the one promised in \cref{sec:trace-norms}, and the two statements together give the best approximation of a given rank promised after \cref{thm:rank-factorisation}. The theorem is usually named after Eckart and Young, who stated it for the Frobenius norm in 1936~\cite{eckart1936}; Schmidt had proved the corresponding result for integral operators in 1907~\cite{schmidt1907,stewart1993}, and Mirsky extended it to every unitarily invariant norm in 1960~\cite{mirsky1960}.

For the $3\times2$ matrix of \cref{sec:svd-existence}, with $\sigma_1=\sqrt3$ and $\sigma_2=1$, the best approximation of rank one is
\[
A_1=\sqrt3\,\mathbf u_1\mathbf v_1^{\mathsf T}=\sqrt3\,\frac1{\sqrt6}\begin{bmatrix}2\\1\\1\end{bmatrix}\frac1{\sqrt2}\begin{bmatrix}1&1\end{bmatrix}=\begin{bmatrix}1&1\\1/2&1/2\\1/2&1/2\end{bmatrix},\qquad A-A_1=\begin{bmatrix}0&0\\1/2&-1/2\\-1/2&1/2\end{bmatrix},
\]
and the error has the spectral and the Frobenius norm $1=\sigma_2$, as the theorem says; it is the second term $\sigma_2\mathbf u_2\mathbf v_2^{\mathsf T}$ of the expansion. The approximation replaces the second and the third row, $(1,0)$ and $(0,1)$, by their average, and keeps the first.

\subsection{The Distance to the Nearest Singular Matrix}

For a square matrix, the case $k=n-1$ of the theorem measures how far $A$ is from being singular.

\begin{corollary}[Distance to singularity]
\label{cor:distance-singular}
Let $A\in\mathbb F^{n\times n}$. Then
\[
\min\{\|E\|_2:A+E\text{ singular}\}=\min\{\|E\|_{\mathrm F}:A+E\text{ singular}\}=\sigma_n(A),
\]
attained at $E=-\sigma_n\mathbf u_n\mathbf v_n^{\mathsf H}$. If $A$ is invertible, then $\sigma_n=1/\|A^{-1}\|_2$, and the distance relative to $\|A\|_2$ is
\[
\min\Bigl\{\frac{\|E\|_2}{\|A\|_2}:A+E\text{ singular}\Bigr\}=\frac{\sigma_n}{\sigma_1}=\frac1{\kappa_2(A)},\qquad\kappa_2(A)=\|A\|_2\|A^{-1}\|_2 .
\]
More generally, for $A\in\mathbb F^{m\times n}$ and $1\le k\le p$, the distance of $A$ from the matrices of rank less than $k$ is $\sigma_k(A)$ in the spectral norm.
\end{corollary}

\begin{proof}
A square matrix is singular exactly when its rank is at most $n-1$, so the first statement is \cref{thm:eckart-young} with $k=n-1$ for $B=A+E$, in both norms, since $\sigma_n$ is the last singular value and the Frobenius bound has a single term; $A_{n-1}=A-\sigma_n\mathbf u_n\mathbf v_n^{\mathsf H}$. If $A$ is invertible, $A^{-1}=V\Sigma^{-1}U^{\mathsf H}$ is an SVD up to the order of the diagonal, so $\|A^{-1}\|_2=1/\sigma_n$ by \cref{thm:two-norm}. The last statement is the theorem with $k-1$ in place of $k$.
\end{proof}

This is the statement promised in \cref{sec:det-computation}: $1/\kappa_2(A)$ is the relative distance from $A$ to the nearest singular matrix. A matrix is nearly singular when its smallest singular value is small compared with its largest, whatever its determinant. For the triangular matrix $U_{60}$ of that section, with $\det U_{60}=1$, a change of one entry by $2^{-58}$ made it singular, so $\sigma_{60}(U_{60})\le2^{-58}\approx3.5\times10^{-18}$; \cref{sec:svd-condition} computes the exact value. For $2\times2$ matrices the corollary can be checked by hand: for $\varepsilon>0$ the matrix with the rows $(1,1)$ and $(1,1+\varepsilon)$ of \cref{sec:inverse} has the singular values $(2+\varepsilon\pm(4+\varepsilon^2)^{1/2})/2$, which are its eigenvalues because it is positive definite, so its distance from the singular matrices is $(2+\varepsilon-(4+\varepsilon^2)^{1/2})/2=\varepsilon/2+O(\varepsilon^2)$; the perturbation $-\sigma_2\mathbf u_2\mathbf v_2^{\mathsf T}$ that attains it changes all four entries by about $\varepsilon/4$, less than the obvious perturbation that changes one entry by $\varepsilon$.

\subsection{The Numerical Rank}

\Cref{sec:computing-rank} found that the rank of a matrix of floating-point numbers is almost always full and therefore almost never the number one wants, and that the library routine \code{linalg.matrix\_rank} reports instead the number of singular values above a tolerance. The theorem of Eckart and Young gives this number its meaning.

\begin{proposition}[Numerical rank]
\label{prop:numerical-rank}
Let $A\in\mathbb F^{m\times n}$ and $\tau\ge0$, and let $k$ be the number of singular values of $A$ that exceed $\tau$. Then
\[
k=\min\bigl\{\operatorname{rank}B:B\in\mathbb F^{m\times n},\ \|A-B\|_2\le\tau\bigr\},
\]
and the truncated SVD $A_k$ is a matrix at which the minimum is attained.
\end{proposition}

\begin{proof}
The matrix $A_k$ has rank $k$, and $\|A-A_k\|_2=\sigma_{k+1}\le\tau$ by \eqref{eq:truncation-error}, since $\sigma_{k+1}$ does not exceed $\tau$ (or is $0$ if $k=r$). If $k=0$ this shows that the minimum is $0$. If $k\ge1$ and $\operatorname{rank}B\le k-1$, then $\|A-B\|_2\ge\sigma_k>\tau$ by \cref{thm:eckart-young}, so no matrix of rank less than $k$ lies within the distance $\tau$.
\end{proof}

The number $k$ is the \term{numerical rank} of $A$ at the tolerance $\tau$: the smallest rank of a matrix that cannot be told apart from $A$ by a perturbation of norm $\tau$. If the entries of $A$ are measurements or rounded numbers with errors of norm about $\tau$, the rank of the unknown exact matrix is at least the numerical rank, and a matrix of exactly that rank is consistent with the data. The answer is robust only when the singular values have a gap at the tolerance, $\sigma_k\gg\tau\gg\sigma_{k+1}$; if a singular value lies near $\tau$, the numerical rank changes with small changes of $\tau$ or of $A$, as it must, since the question is then ill posed. \Cref{sec:svd-computation} returns to the experiment of \cref{tab:rank-experiment} with this in mind.

\subsection{Compressing a Picture}

A grey-scale picture of $m\times n$ pixels is a matrix of intensities, and the truncated SVD stores it approximately in $k(m+n+1)$ numbers, the vectors $\mathbf u_i$, $\mathbf v_i$ and the numbers $\sigma_i$ for $i\le k$, instead of $mn$. \Cref{fig:low-rank-pattern} shows this for a synthetic $32\times32$ pattern that we constructed for the purpose: a background whose intensity increases linearly from left to right, and four shapes of constant intensity, a ring, a square standing on a corner, a triangle and an ellipse inside the triangle. The intensities lie between $0.05$ and $0.9$, and the picture shows $100$ times the intensity, rounded to a multiple of $5$, as the percentage of black.

\begin{figure}
% \sdr{x}{y}{w}{h}: a w x h block of pixels with lower left corner (x,y), in millimetres, one
% millimetre per pixel, enlarged by 0.05 mm on each side so that neighbouring blocks overlap;
% \sdlevel{p}{blocks}: one path of blocks, filled with p% black.
\def\sdr#1#2#3#4{\pgfpathrectangle{\pgfqpoint{\dimexpr#1mm-0.05mm\relax}{\dimexpr#2mm-0.05mm\relax}}%
  {\pgfqpoint{\dimexpr#3mm+0.1mm\relax}{\dimexpr#4mm+0.1mm\relax}}}
\def\sdlevel#1#2{\begin{pgfscope}\pgfsetfillcolor{black!#1}#2\pgfusepath{fill}\end{pgfscope}}
\newcommand\sdimg[1]{\begin{tikzpicture}#1\draw[line width=0.3pt] (0,0) rectangle (32mm,32mm);\end{tikzpicture}}
% Data for fig:low-rank-pattern: 100 times the entries of the pattern A (\sdimgA) and of its truncated
% SVDs A_1, A_2, A_4, A_8 (\sdimgB to \sdimgE), clipped to [0,100] and rounded to multiples of 5,
% used as percentages of black; rows are counted from the top, so row i of the matrix occupies
% y from 32-i to 33-i. Equal pixels are merged into rectangles, one \sdlevel per percentage.
\def\sdimgA{%
\sdlevel{5}{\sdr{14}{7}{4}{1}\sdr{12}{4}{8}{3}}%
\sdlevel{10}{\sdr{0}{24}{3}{8}\sdr{0}{0}{3}{22}\sdr{0}{22}{2}{2}}%
\sdlevel{15}{\sdr{6}{24}{2}{2}\sdr{3}{27}{1}{1}\sdr{3}{29}{5}{3}\sdr{3}{28}{2}{1}\sdr{3}{18}{1}{2}\sdr{3}{6}{5}{11}\sdr{3}{0}{5}{2}\sdr{3}{4}{2}{1}\sdr{3}{3}{1}{1}\sdr{6}{21}{2}{1}\sdr{5}{22}{3}{2}\sdr{7}{20}{1}{1}\sdr{3}{17}{3}{1}\sdr{3}{5}{4}{1}}%
\sdlevel{20}{\sdr{8}{25}{3}{1}\sdr{12}{28}{1}{1}\sdr{8}{30}{5}{2}\sdr{8}{21}{4}{4}\sdr{8}{20}{3}{1}\sdr{8}{26}{2}{1}\sdr{8}{11}{5}{6}\sdr{8}{0}{5}{2}\sdr{10}{29}{3}{1}\sdr{12}{17}{1}{1}\sdr{8}{8}{2}{1}\sdr{8}{9}{3}{1}\sdr{8}{10}{4}{1}\sdr{8}{7}{1}{1}}%
\sdlevel{25}{\sdr{15}{24}{3}{1}\sdr{15}{25}{4}{1}\sdr{14}{26}{5}{2}\sdr{13}{28}{6}{4}\sdr{13}{14}{6}{5}\sdr{13}{0}{6}{2}\sdr{14}{19}{5}{1}\sdr{15}{20}{4}{2}\sdr{15}{22}{3}{1}\sdr{15}{23}{2}{1}\sdr{13}{13}{2}{1}\sdr{17}{13}{2}{1}\sdr{13}{12}{1}{1}\sdr{18}{12}{1}{1}}%
\sdlevel{30}{\sdr{19}{26}{1}{1}\sdr{19}{27}{2}{1}\sdr{19}{28}{3}{1}\sdr{19}{29}{5}{3}\sdr{19}{11}{5}{6}\sdr{19}{0}{5}{2}\sdr{19}{18}{3}{1}\sdr{19}{19}{2}{1}\sdr{19}{20}{1}{1}\sdr{19}{17}{4}{1}\sdr{20}{10}{4}{1}\sdr{21}{9}{3}{1}\sdr{22}{8}{2}{1}\sdr{23}{7}{1}{1}}%
\sdlevel{35}{\sdr{28}{24}{1}{1}\sdr{27}{25}{2}{1}\sdr{26}{26}{3}{1}\sdr{25}{27}{4}{1}\sdr{24}{28}{5}{4}\sdr{25}{18}{4}{1}\sdr{27}{20}{2}{1}\sdr{28}{21}{1}{1}\sdr{24}{6}{5}{12}\sdr{24}{0}{5}{2}\sdr{25}{5}{4}{1}\sdr{26}{19}{3}{1}\sdr{27}{4}{2}{1}\sdr{28}{3}{1}{1}}%
\sdlevel{40}{\sdr{29}{0}{3}{32}}%
\sdlevel{55}{\sdr{15}{13}{2}{1}\sdr{14}{12}{4}{1}\sdr{13}{11}{6}{1}\sdr{12}{10}{8}{1}\sdr{11}{9}{10}{1}\sdr{10}{8}{12}{1}\sdr{9}{7}{5}{1}\sdr{18}{7}{5}{1}\sdr{8}{6}{4}{1}\sdr{20}{6}{4}{1}\sdr{7}{5}{5}{1}\sdr{20}{5}{5}{1}\sdr{5}{4}{7}{1}\sdr{20}{4}{7}{1}\sdr{4}{3}{24}{1}\sdr{3}{2}{26}{1}}%
\sdlevel{75}{\sdr{22}{28}{2}{1}\sdr{21}{27}{4}{1}\sdr{20}{26}{6}{1}\sdr{19}{25}{8}{1}\sdr{18}{24}{10}{1}\sdr{17}{23}{12}{1}\sdr{18}{22}{11}{1}\sdr{19}{21}{9}{1}\sdr{20}{20}{7}{1}\sdr{21}{19}{5}{1}\sdr{22}{18}{3}{1}\sdr{23}{17}{1}{1}}%
\sdlevel{90}{\sdr{3}{24}{3}{2}\sdr{11}{25}{4}{1}\sdr{4}{27}{10}{1}\sdr{8}{29}{2}{1}\sdr{5}{28}{7}{1}\sdr{4}{19}{10}{1}\sdr{3}{26}{5}{1}\sdr{10}{26}{4}{1}\sdr{3}{21}{3}{1}\sdr{11}{20}{4}{1}\sdr{12}{21}{3}{4}\sdr{2}{22}{3}{2}\sdr{3}{20}{4}{1}\sdr{4}{18}{9}{1}\sdr{6}{17}{6}{1}}%
}
\def\sdimgB{%
\sdlevel{5}{\sdr{0}{30}{2}{2}\sdr{0}{29}{1}{1}\sdr{0}{11}{1}{1}\sdr{0}{12}{2}{5}\sdr{0}{0}{2}{2}\sdr{0}{5}{1}{2}}%
\sdlevel{10}{\sdr{0}{28}{2}{1}\sdr{2}{30}{1}{2}\sdr{1}{29}{1}{1}\sdr{0}{17}{2}{1}\sdr{1}{11}{1}{1}\sdr{0}{7}{2}{4}\sdr{0}{18}{1}{1}\sdr{2}{12}{1}{5}\sdr{2}{0}{1}{2}\sdr{1}{5}{1}{2}\sdr{0}{4}{2}{1}\sdr{0}{2}{1}{2}}%
\sdlevel{15}{\sdr{15}{30}{2}{2}\sdr{2}{29}{1}{1}\sdr{1}{18}{1}{1}\sdr{15}{14}{2}{3}\sdr{15}{0}{2}{2}\sdr{2}{4}{1}{8}\sdr{0}{19}{2}{9}\sdr{1}{2}{1}{2}}%
\sdlevel{20}{\sdr{2}{28}{1}{1}\sdr{3}{30}{1}{2}\sdr{17}{30}{2}{2}\sdr{29}{30}{2}{2}\sdr{15}{29}{3}{1}\sdr{15}{11}{3}{3}\sdr{15}{10}{2}{1}\sdr{2}{17}{1}{1}\sdr{3}{14}{1}{3}\sdr{17}{14}{2}{3}\sdr{29}{14}{2}{3}\sdr{3}{0}{1}{2}\sdr{17}{0}{2}{2}\sdr{29}{0}{2}{2}\sdr{15}{6}{3}{1}\sdr{2}{2}{1}{2}\sdr{15}{5}{2}{1}}%
\sdlevel{25}{\sdr{6}{30}{5}{2}\sdr{14}{30}{1}{2}\sdr{19}{30}{1}{2}\sdr{26}{30}{3}{2}\sdr{31}{30}{1}{2}\sdr{3}{29}{1}{1}\sdr{28}{29}{4}{1}\sdr{18}{29}{1}{1}\sdr{3}{11}{1}{3}\sdr{6}{12}{4}{1}\sdr{18}{11}{1}{1}\sdr{28}{12}{4}{1}\sdr{15}{7}{3}{3}\sdr{6}{13}{5}{4}\sdr{14}{14}{1}{3}\sdr{19}{14}{1}{3}\sdr{26}{14}{3}{3}\sdr{31}{14}{1}{3}\sdr{6}{0}{5}{2}\sdr{14}{0}{1}{2}\sdr{19}{0}{1}{2}\sdr{26}{0}{3}{2}\sdr{31}{0}{1}{2}\sdr{3}{5}{1}{2}\sdr{7}{29}{2}{1}\sdr{18}{6}{1}{1}\sdr{28}{6}{4}{1}\sdr{2}{18}{1}{10}\sdr{15}{17}{1}{1}\sdr{18}{12}{2}{2}\sdr{27}{13}{5}{1}\sdr{6}{6}{4}{1}\sdr{8}{11}{1}{1}\sdr{29}{11}{3}{1}\sdr{17}{10}{1}{1}\sdr{15}{4}{3}{1}\sdr{17}{5}{2}{1}\sdr{29}{5}{2}{1}}%
\sdlevel{30}{\sdr{4}{30}{2}{2}\sdr{11}{30}{1}{2}\sdr{13}{30}{1}{2}\sdr{20}{30}{3}{2}\sdr{24}{30}{2}{2}\sdr{14}{29}{1}{1}\sdr{26}{29}{2}{1}\sdr{9}{29}{2}{1}\sdr{19}{29}{1}{1}\sdr{3}{9}{1}{2}\sdr{6}{10}{4}{1}\sdr{18}{10}{2}{1}\sdr{28}{9}{4}{2}\sdr{5}{12}{1}{1}\sdr{10}{12}{1}{1}\sdr{14}{11}{1}{3}\sdr{19}{11}{1}{1}\sdr{26}{12}{2}{1}\sdr{27}{11}{2}{1}\sdr{3}{7}{1}{1}\sdr{7}{9}{2}{1}\sdr{18}{8}{1}{2}\sdr{4}{13}{2}{4}\sdr{11}{13}{1}{4}\sdr{13}{14}{1}{3}\sdr{20}{14}{3}{3}\sdr{24}{14}{2}{3}\sdr{4}{0}{2}{2}\sdr{11}{0}{1}{2}\sdr{13}{0}{1}{2}\sdr{20}{0}{3}{2}\sdr{24}{0}{2}{2}\sdr{6}{29}{1}{1}\sdr{9}{11}{2}{1}\sdr{14}{5}{1}{2}\sdr{19}{5}{1}{2}\sdr{26}{6}{2}{1}\sdr{15}{28}{2}{1}\sdr{16}{17}{2}{1}\sdr{20}{13}{2}{1}\sdr{25}{13}{2}{1}\sdr{5}{6}{1}{1}\sdr{10}{6}{1}{1}\sdr{20}{12}{1}{1}\sdr{6}{11}{2}{1}\sdr{27}{5}{2}{1}\sdr{3}{4}{1}{1}\sdr{6}{7}{4}{1}\sdr{18}{7}{2}{1}\sdr{28}{7}{4}{1}\sdr{7}{4}{2}{1}\sdr{18}{4}{1}{1}\sdr{29}{8}{2}{1}\sdr{6}{5}{5}{1}\sdr{31}{5}{1}{1}\sdr{29}{4}{3}{1}}%
\sdlevel{35}{\sdr{12}{30}{1}{2}\sdr{23}{30}{1}{2}\sdr{4}{29}{2}{1}\sdr{11}{29}{1}{1}\sdr{20}{29}{2}{1}\sdr{3}{17}{1}{1}\sdr{5}{10}{1}{1}\sdr{10}{10}{1}{1}\sdr{14}{9}{1}{2}\sdr{26}{10}{2}{1}\sdr{4}{12}{1}{1}\sdr{11}{11}{1}{1}\sdr{13}{29}{1}{1}\sdr{24}{29}{2}{1}\sdr{4}{11}{2}{1}\sdr{20}{10}{1}{1}\sdr{25}{11}{2}{1}\sdr{6}{9}{1}{1}\sdr{9}{9}{2}{1}\sdr{14}{7}{1}{1}\sdr{19}{8}{1}{2}\sdr{27}{8}{2}{1}\sdr{12}{14}{1}{3}\sdr{23}{14}{1}{3}\sdr{12}{0}{1}{2}\sdr{23}{0}{1}{2}\sdr{4}{5}{2}{1}\sdr{11}{5}{1}{2}\sdr{13}{6}{1}{1}\sdr{20}{11}{2}{1}\sdr{24}{6}{2}{1}\sdr{17}{28}{1}{1}\sdr{15}{21}{1}{2}\sdr{15}{18}{2}{1}\sdr{3}{8}{1}{1}\sdr{7}{17}{2}{1}\sdr{18}{17}{1}{1}\sdr{29}{17}{3}{1}\sdr{12}{13}{2}{1}\sdr{22}{13}{3}{1}\sdr{4}{6}{1}{1}\sdr{11}{12}{3}{1}\sdr{21}{12}{5}{1}\sdr{25}{5}{2}{1}\sdr{5}{7}{1}{1}\sdr{10}{7}{1}{1}\sdr{14}{4}{1}{1}\sdr{20}{5}{1}{1}\sdr{26}{7}{2}{1}\sdr{6}{4}{1}{1}\sdr{9}{4}{2}{1}\sdr{19}{4}{1}{1}\sdr{27}{9}{1}{1}\sdr{6}{8}{5}{1}\sdr{27}{4}{2}{1}\sdr{31}{8}{1}{1}\sdr{20}{6}{3}{1}\sdr{15}{2}{3}{2}}%
\sdlevel{40}{\sdr{15}{27}{3}{1}\sdr{15}{23}{3}{3}\sdr{15}{26}{2}{1}\sdr{14}{17}{1}{1}\sdr{4}{10}{1}{1}\sdr{11}{9}{1}{1}\sdr{20}{9}{2}{1}\sdr{24}{10}{2}{1}\sdr{12}{29}{1}{1}\sdr{12}{11}{2}{1}\sdr{4}{9}{2}{1}\sdr{11}{7}{1}{1}\sdr{20}{8}{1}{1}\sdr{25}{9}{2}{1}\sdr{18}{28}{1}{1}\sdr{12}{6}{1}{1}\sdr{22}{29}{2}{1}\sdr{3}{28}{1}{1}\sdr{7}{28}{3}{1}\sdr{18}{3}{1}{1}\sdr{28}{28}{4}{1}\sdr{15}{20}{3}{1}\sdr{15}{19}{2}{1}\sdr{16}{21}{2}{2}\sdr{17}{18}{1}{1}\sdr{6}{17}{1}{1}\sdr{9}{17}{2}{1}\sdr{14}{8}{1}{1}\sdr{19}{17}{1}{1}\sdr{27}{17}{2}{1}\sdr{12}{5}{2}{1}\sdr{22}{11}{3}{1}\sdr{4}{7}{1}{1}\sdr{11}{10}{3}{1}\sdr{21}{10}{2}{1}\sdr{24}{7}{2}{1}\sdr{4}{4}{2}{1}\sdr{11}{4}{1}{1}\sdr{20}{7}{2}{1}\sdr{25}{4}{2}{1}\sdr{5}{8}{1}{1}\sdr{20}{4}{1}{1}\sdr{26}{8}{1}{1}\sdr{23}{6}{1}{1}\sdr{21}{5}{4}{1}\sdr{29}{3}{1}{1}}%
\sdlevel{45}{\sdr{17}{26}{1}{1}\sdr{18}{21}{1}{1}\sdr{12}{9}{2}{1}\sdr{12}{7}{2}{1}\sdr{3}{18}{1}{1}\sdr{19}{28}{1}{1}\sdr{7}{18}{2}{1}\sdr{18}{18}{1}{1}\sdr{6}{28}{1}{1}\sdr{10}{28}{1}{1}\sdr{14}{28}{1}{1}\sdr{19}{3}{1}{1}\sdr{27}{28}{1}{1}\sdr{17}{19}{1}{1}\sdr{18}{2}{1}{1}\sdr{29}{21}{1}{1}\sdr{3}{2}{1}{2}\sdr{7}{2}{2}{1}\sdr{29}{18}{3}{1}\sdr{4}{17}{2}{1}\sdr{20}{17}{1}{1}\sdr{25}{17}{2}{1}\sdr{23}{10}{1}{1}\sdr{12}{4}{2}{1}\sdr{22}{9}{3}{1}\sdr{4}{8}{1}{1}\sdr{11}{8}{1}{1}\sdr{13}{8}{1}{1}\sdr{21}{8}{2}{1}\sdr{24}{8}{2}{1}\sdr{22}{7}{2}{1}\sdr{21}{4}{4}{1}\sdr{6}{3}{4}{1}\sdr{28}{3}{1}{1}\sdr{30}{3}{2}{1}\sdr{28}{2}{4}{1}}%
\sdlevel{50}{\sdr{29}{26}{3}{2}\sdr{28}{24}{4}{2}\sdr{7}{23}{2}{5}\sdr{18}{22}{1}{6}\sdr{29}{23}{3}{1}\sdr{19}{21}{1}{1}\sdr{28}{21}{1}{1}\sdr{6}{18}{1}{1}\sdr{19}{18}{1}{1}\sdr{4}{28}{2}{1}\sdr{20}{28}{1}{1}\sdr{25}{28}{2}{1}\sdr{3}{19}{1}{9}\sdr{7}{20}{2}{1}\sdr{18}{19}{1}{2}\sdr{29}{19}{3}{2}\sdr{28}{22}{4}{1}\sdr{7}{22}{3}{1}\sdr{6}{21}{4}{1}\sdr{19}{2}{1}{1}\sdr{28}{18}{1}{1}\sdr{30}{21}{2}{1}\sdr{6}{2}{1}{1}\sdr{9}{18}{1}{1}\sdr{11}{17}{3}{1}\sdr{21}{17}{2}{1}\sdr{24}{17}{1}{1}\sdr{12}{8}{1}{1}\sdr{23}{8}{1}{1}\sdr{10}{3}{1}{1}\sdr{14}{3}{1}{1}\sdr{27}{2}{1}{2}\sdr{9}{2}{2}{1}}%
\sdlevel{55}{\sdr{9}{27}{2}{1}\sdr{27}{27}{2}{1}\sdr{27}{24}{1}{2}\sdr{6}{22}{1}{6}\sdr{9}{26}{1}{1}\sdr{19}{22}{1}{6}\sdr{27}{23}{2}{1}\sdr{28}{26}{1}{1}\sdr{10}{21}{1}{2}\sdr{14}{21}{1}{1}\sdr{14}{18}{1}{1}\sdr{11}{28}{1}{1}\sdr{13}{28}{1}{1}\sdr{21}{28}{1}{1}\sdr{24}{28}{1}{1}\sdr{6}{20}{1}{1}\sdr{9}{23}{2}{3}\sdr{19}{19}{1}{2}\sdr{27}{20}{2}{1}\sdr{9}{20}{1}{1}\sdr{28}{19}{1}{1}\sdr{27}{21}{1}{2}\sdr{10}{18}{1}{1}\sdr{14}{2}{1}{1}\sdr{6}{19}{4}{1}\sdr{26}{18}{2}{1}\sdr{23}{17}{1}{1}\sdr{4}{3}{2}{1}\sdr{20}{2}{1}{2}\sdr{25}{3}{2}{1}\sdr{5}{2}{1}{1}\sdr{26}{2}{1}{1}}%
\sdlevel{60}{\sdr{26}{27}{1}{1}\sdr{5}{24}{1}{2}\sdr{10}{26}{1}{1}\sdr{14}{22}{1}{6}\sdr{27}{26}{1}{1}\sdr{20}{21}{1}{1}\sdr{11}{18}{1}{1}\sdr{25}{18}{1}{1}\sdr{12}{28}{1}{1}\sdr{22}{28}{2}{1}\sdr{14}{19}{1}{2}\sdr{26}{20}{1}{6}\sdr{10}{19}{1}{2}\sdr{27}{19}{1}{1}\sdr{5}{21}{1}{2}\sdr{20}{18}{1}{1}\sdr{4}{18}{2}{1}\sdr{11}{2}{1}{2}\sdr{25}{2}{1}{1}\sdr{13}{3}{1}{1}\sdr{21}{2}{1}{2}\sdr{24}{3}{1}{1}\sdr{4}{2}{1}{1}}%
\sdlevel{65}{\sdr{11}{27}{1}{1}\sdr{25}{27}{1}{1}\sdr{4}{26}{2}{2}\sdr{4}{24}{1}{2}\sdr{4}{23}{2}{1}\sdr{20}{22}{1}{6}\sdr{25}{26}{2}{1}\sdr{4}{19}{2}{2}\sdr{11}{20}{1}{6}\sdr{20}{19}{1}{2}\sdr{25}{20}{1}{6}\sdr{25}{19}{2}{1}\sdr{4}{21}{1}{2}\sdr{21}{21}{1}{1}\sdr{13}{18}{1}{1}\sdr{21}{18}{2}{1}\sdr{24}{18}{1}{1}\sdr{12}{3}{1}{1}\sdr{22}{3}{2}{1}\sdr{12}{2}{2}{1}\sdr{22}{2}{3}{1}}%
\sdlevel{70}{\sdr{13}{27}{1}{1}\sdr{21}{26}{1}{2}\sdr{21}{24}{2}{2}\sdr{13}{23}{1}{3}\sdr{21}{23}{1}{1}\sdr{24}{22}{1}{6}\sdr{11}{26}{1}{1}\sdr{13}{20}{1}{1}\sdr{21}{19}{1}{2}\sdr{24}{19}{1}{2}\sdr{11}{19}{1}{1}\sdr{21}{22}{2}{1}\sdr{12}{21}{2}{2}\sdr{22}{21}{3}{1}\sdr{12}{18}{1}{1}\sdr{23}{18}{1}{1}}%
\sdlevel{75}{\sdr{12}{27}{1}{1}\sdr{22}{26}{2}{2}\sdr{23}{24}{1}{2}\sdr{12}{23}{1}{3}\sdr{22}{23}{2}{1}\sdr{12}{26}{2}{1}\sdr{12}{20}{1}{1}\sdr{22}{19}{2}{2}\sdr{12}{19}{2}{1}\sdr{23}{22}{1}{1}}%
}
\def\sdimgC{%
\sdlevel{5}{\sdr{14}{28}{1}{1}\sdr{0}{30}{2}{2}\sdr{0}{29}{1}{1}\sdr{0}{11}{1}{1}\sdr{0}{12}{2}{5}\sdr{0}{0}{2}{2}\sdr{0}{6}{1}{1}\sdr{2}{28}{2}{1}\sdr{14}{17}{1}{1}\sdr{8}{22}{1}{2}\sdr{2}{4}{1}{3}}%
\sdlevel{10}{\sdr{2}{27}{1}{1}\sdr{0}{28}{2}{1}\sdr{0}{17}{2}{1}\sdr{1}{11}{1}{1}\sdr{0}{7}{2}{4}\sdr{1}{29}{2}{1}\sdr{0}{4}{2}{2}\sdr{2}{18}{1}{1}\sdr{7}{22}{1}{2}\sdr{9}{22}{1}{2}\sdr{1}{6}{1}{1}\sdr{3}{4}{1}{3}\sdr{0}{2}{1}{2}}%
\sdlevel{15}{\sdr{0}{20}{2}{8}\sdr{3}{29}{1}{1}\sdr{2}{30}{1}{2}\sdr{7}{24}{3}{1}\sdr{2}{12}{1}{5}\sdr{2}{7}{1}{1}\sdr{2}{0}{1}{2}\sdr{3}{18}{1}{1}\sdr{19}{28}{1}{1}\sdr{0}{18}{2}{1}\sdr{7}{11}{3}{2}\sdr{6}{22}{1}{2}\sdr{10}{22}{1}{2}\sdr{7}{21}{2}{1}\sdr{0}{19}{3}{1}\sdr{18}{17}{2}{1}\sdr{6}{10}{4}{1}\sdr{8}{9}{1}{1}\sdr{14}{6}{6}{1}\sdr{14}{5}{2}{1}\sdr{17}{5}{3}{1}\sdr{14}{4}{1}{1}\sdr{1}{2}{1}{2}}%
\sdlevel{20}{\sdr{6}{24}{1}{1}\sdr{15}{30}{3}{2}\sdr{6}{30}{5}{2}\sdr{29}{30}{2}{2}\sdr{6}{21}{1}{1}\sdr{10}{24}{1}{1}\sdr{2}{26}{1}{1}\sdr{15}{11}{1}{1}\sdr{2}{8}{1}{4}\sdr{6}{13}{5}{4}\sdr{15}{13}{3}{5}\sdr{29}{14}{2}{3}\sdr{6}{0}{5}{2}\sdr{15}{0}{3}{2}\sdr{29}{0}{2}{2}\sdr{14}{29}{6}{1}\sdr{17}{28}{2}{1}\sdr{3}{27}{1}{1}\sdr{2}{2}{1}{2}\sdr{6}{11}{1}{2}\sdr{10}{10}{1}{3}\sdr{9}{21}{1}{1}\sdr{14}{18}{1}{1}\sdr{15}{12}{2}{1}\sdr{15}{7}{1}{1}\sdr{6}{9}{2}{1}\sdr{9}{9}{2}{1}\sdr{16}{5}{1}{1}\sdr{15}{4}{5}{1}}%
\sdlevel{25}{\sdr{3}{30}{1}{2}\sdr{18}{30}{2}{2}\sdr{3}{19}{1}{1}\sdr{5}{30}{1}{2}\sdr{11}{30}{1}{2}\sdr{18}{18}{2}{1}\sdr{27}{30}{2}{2}\sdr{31}{30}{1}{2}\sdr{15}{28}{2}{1}\sdr{16}{11}{2}{1}\sdr{13}{17}{1}{1}\sdr{26}{29}{6}{1}\sdr{3}{13}{1}{4}\sdr{5}{14}{1}{3}\sdr{11}{10}{1}{7}\sdr{18}{14}{2}{3}\sdr{27}{14}{2}{3}\sdr{31}{14}{1}{3}\sdr{3}{0}{1}{2}\sdr{5}{0}{1}{2}\sdr{11}{0}{1}{2}\sdr{18}{0}{2}{2}\sdr{27}{0}{2}{2}\sdr{31}{0}{1}{2}\sdr{26}{6}{6}{1}\sdr{15}{8}{2}{3}\sdr{7}{25}{3}{1}\sdr{10}{21}{1}{1}\sdr{13}{5}{1}{2}\sdr{18}{13}{1}{1}\sdr{28}{13}{4}{1}\sdr{17}{12}{1}{1}\sdr{29}{11}{3}{2}\sdr{16}{7}{2}{1}\sdr{29}{10}{1}{1}\sdr{6}{8}{5}{1}\sdr{20}{5}{1}{2}\sdr{26}{5}{4}{1}\sdr{27}{4}{1}{1}}%
\sdlevel{30}{\sdr{14}{30}{1}{2}\sdr{20}{29}{1}{1}\sdr{4}{29}{1}{3}\sdr{14}{19}{1}{1}\sdr{20}{30}{7}{2}\sdr{17}{27}{1}{1}\sdr{3}{12}{1}{1}\sdr{18}{12}{2}{1}\sdr{4}{14}{1}{3}\sdr{25}{29}{1}{1}\sdr{30}{10}{2}{1}\sdr{13}{28}{1}{2}\sdr{20}{17}{1}{1}\sdr{28}{7}{4}{1}\sdr{4}{5}{1}{2}\sdr{14}{13}{1}{4}\sdr{20}{14}{7}{3}\sdr{4}{0}{1}{2}\sdr{14}{0}{1}{2}\sdr{20}{0}{7}{2}\sdr{13}{4}{1}{1}\sdr{20}{4}{1}{1}\sdr{25}{5}{1}{2}\sdr{26}{28}{2}{1}\sdr{14}{27}{2}{1}\sdr{17}{8}{1}{3}\sdr{6}{25}{1}{1}\sdr{10}{25}{1}{1}\sdr{2}{20}{1}{1}\sdr{15}{18}{3}{1}\sdr{26}{17}{3}{1}\sdr{4}{13}{2}{1}\sdr{19}{13}{2}{1}\sdr{25}{13}{3}{1}\sdr{3}{7}{1}{1}\sdr{5}{10}{1}{3}\sdr{18}{7}{2}{1}\sdr{27}{12}{2}{1}\sdr{18}{11}{1}{1}\sdr{28}{11}{1}{1}\sdr{30}{5}{2}{1}\sdr{11}{9}{1}{1}\sdr{29}{9}{3}{1}\sdr{29}{8}{2}{1}\sdr{28}{4}{4}{1}\sdr{12}{5}{1}{2}\sdr{21}{6}{1}{1}\sdr{26}{4}{1}{1}\sdr{15}{3}{1}{1}}%
\sdlevel{35}{\sdr{2}{24}{1}{2}\sdr{12}{29}{1}{1}\sdr{12}{30}{2}{2}\sdr{4}{17}{1}{1}\sdr{5}{29}{1}{1}\sdr{28}{28}{1}{1}\sdr{11}{22}{1}{2}\sdr{14}{12}{1}{1}\sdr{20}{28}{1}{1}\sdr{26}{11}{2}{1}\sdr{24}{29}{1}{1}\sdr{21}{29}{1}{1}\sdr{4}{12}{1}{1}\sdr{12}{17}{1}{1}\sdr{25}{17}{1}{1}\sdr{12}{13}{2}{4}\sdr{12}{0}{2}{2}\sdr{5}{8}{1}{2}\sdr{12}{4}{1}{1}\sdr{21}{5}{1}{1}\sdr{24}{5}{1}{2}\sdr{20}{7}{1}{1}\sdr{28}{8}{1}{2}\sdr{16}{27}{1}{1}\sdr{18}{27}{2}{1}\sdr{15}{26}{1}{1}\sdr{2}{21}{1}{1}\sdr{11}{8}{1}{1}\sdr{7}{20}{2}{1}\sdr{15}{19}{5}{1}\sdr{4}{4}{1}{1}\sdr{25}{4}{1}{1}\sdr{29}{17}{2}{1}\sdr{21}{13}{4}{1}\sdr{14}{7}{1}{1}\sdr{20}{12}{7}{1}\sdr{3}{11}{1}{1}\sdr{19}{11}{1}{1}\sdr{26}{7}{2}{1}\sdr{18}{8}{1}{3}\sdr{27}{10}{2}{1}\sdr{31}{8}{1}{1}\sdr{6}{7}{5}{1}\sdr{16}{3}{2}{1}\sdr{15}{2}{3}{1}}%
\sdlevel{40}{\sdr{11}{24}{1}{1}\sdr{15}{25}{2}{1}\sdr{4}{28}{1}{1}\sdr{11}{21}{1}{1}\sdr{21}{17}{1}{1}\sdr{22}{29}{2}{1}\sdr{21}{7}{1}{1}\sdr{3}{8}{1}{2}\sdr{6}{29}{5}{1}\sdr{22}{6}{2}{1}\sdr{4}{11}{1}{1}\sdr{12}{28}{1}{1}\sdr{25}{28}{1}{1}\sdr{29}{28}{3}{1}\sdr{16}{26}{2}{1}\sdr{15}{20}{2}{1}\sdr{11}{7}{1}{1}\sdr{2}{22}{1}{2}\sdr{15}{21}{1}{1}\sdr{6}{20}{1}{1}\sdr{9}{20}{1}{1}\sdr{26}{18}{3}{1}\sdr{21}{4}{1}{1}\sdr{31}{17}{1}{1}\sdr{12}{12}{2}{1}\sdr{14}{11}{1}{1}\sdr{20}{11}{6}{1}\sdr{3}{10}{2}{1}\sdr{19}{10}{8}{1}\sdr{3}{3}{1}{1}\sdr{19}{8}{1}{2}\sdr{27}{8}{1}{2}\sdr{4}{7}{2}{1}\sdr{13}{7}{1}{1}\sdr{24}{7}{2}{1}\sdr{5}{5}{1}{2}\sdr{22}{5}{1}{1}\sdr{24}{4}{1}{1}\sdr{18}{3}{1}{1}\sdr{29}{3}{1}{1}}%
\sdlevel{45}{\sdr{11}{29}{1}{1}\sdr{29}{21}{2}{1}\sdr{14}{10}{1}{1}\sdr{14}{8}{1}{1}\sdr{3}{26}{1}{1}\sdr{18}{26}{1}{1}\sdr{11}{25}{1}{1}\sdr{21}{28}{1}{1}\sdr{27}{27}{2}{1}\sdr{3}{2}{1}{1}\sdr{18}{2}{1}{1}\sdr{29}{22}{1}{4}\sdr{15}{22}{2}{3}\sdr{5}{22}{1}{1}\sdr{16}{21}{1}{1}\sdr{29}{18}{2}{1}\sdr{10}{20}{1}{1}\sdr{17}{20}{1}{1}\sdr{13}{18}{1}{1}\sdr{20}{18}{1}{1}\sdr{24}{17}{1}{1}\sdr{12}{11}{2}{1}\sdr{14}{3}{1}{1}\sdr{4}{8}{1}{2}\sdr{20}{8}{7}{2}\sdr{12}{7}{1}{1}\sdr{22}{7}{2}{1}\sdr{6}{6}{5}{1}\sdr{23}{5}{1}{1}\sdr{22}{4}{1}{1}\sdr{19}{3}{1}{1}\sdr{28}{3}{1}{1}\sdr{30}{3}{2}{1}\sdr{8}{2}{1}{1}\sdr{28}{2}{4}{1}}%
\sdlevel{50}{\sdr{17}{24}{1}{2}\sdr{5}{23}{1}{2}\sdr{29}{27}{3}{1}\sdr{31}{21}{1}{1}\sdr{12}{10}{2}{1}\sdr{5}{21}{1}{1}\sdr{6}{5}{4}{1}\sdr{14}{26}{1}{1}\sdr{19}{26}{1}{1}\sdr{13}{27}{1}{1}\sdr{20}{27}{1}{1}\sdr{26}{27}{1}{1}\sdr{29}{20}{3}{1}\sdr{14}{2}{1}{1}\sdr{19}{2}{1}{1}\sdr{28}{26}{4}{1}\sdr{17}{21}{1}{2}\sdr{30}{22}{2}{4}\sdr{5}{4}{1}{1}\sdr{31}{18}{1}{1}\sdr{26}{19}{5}{1}\sdr{4}{18}{1}{1}\sdr{12}{18}{1}{1}\sdr{25}{18}{1}{1}\sdr{12}{8}{2}{1}\sdr{12}{9}{3}{1}\sdr{11}{6}{1}{1}\sdr{6}{3}{4}{1}\sdr{23}{4}{1}{1}\sdr{26}{3}{2}{1}\sdr{6}{2}{2}{1}\sdr{9}{2}{1}{1}\sdr{27}{2}{1}{1}}%
\sdlevel{55}{\sdr{28}{25}{1}{1}\sdr{20}{19}{1}{1}\sdr{25}{27}{1}{1}\sdr{5}{25}{1}{1}\sdr{24}{28}{1}{1}\sdr{4}{27}{1}{1}\sdr{25}{3}{1}{1}\sdr{27}{26}{1}{1}\sdr{5}{2}{1}{1}\sdr{28}{20}{1}{2}\sdr{17}{23}{1}{1}\sdr{11}{20}{1}{1}\sdr{18}{20}{1}{1}\sdr{13}{19}{1}{1}\sdr{20}{2}{1}{2}\sdr{31}{19}{1}{1}\sdr{21}{18}{1}{1}\sdr{22}{17}{1}{1}\sdr{10}{5}{2}{1}\sdr{6}{4}{1}{1}\sdr{4}{3}{2}{1}\sdr{10}{2}{1}{2}\sdr{26}{2}{1}{1}}%
\sdlevel{60}{\sdr{18}{25}{1}{1}\sdr{12}{27}{1}{1}\sdr{5}{20}{1}{1}\sdr{24}{18}{1}{1}\sdr{4}{19}{1}{1}\sdr{25}{19}{1}{1}\sdr{22}{28}{1}{1}\sdr{12}{19}{1}{1}\sdr{7}{26}{2}{1}\sdr{20}{26}{1}{1}\sdr{26}{26}{1}{1}\sdr{18}{21}{1}{1}\sdr{28}{22}{1}{3}\sdr{3}{20}{1}{1}\sdr{5}{17}{1}{1}\sdr{27}{20}{1}{1}\sdr{4}{2}{1}{1}\sdr{25}{2}{1}{1}\sdr{24}{3}{1}{1}\sdr{23}{17}{1}{1}\sdr{7}{4}{5}{1}\sdr{11}{3}{3}{1}\sdr{21}{2}{1}{2}\sdr{11}{2}{1}{1}}%
\sdlevel{65}{\sdr{23}{28}{1}{1}\sdr{18}{24}{1}{1}\sdr{27}{24}{1}{2}\sdr{21}{27}{1}{1}\sdr{5}{28}{1}{1}\sdr{23}{21}{1}{4}\sdr{21}{19}{1}{1}\sdr{4}{26}{3}{1}\sdr{9}{26}{1}{1}\sdr{13}{26}{1}{1}\sdr{25}{26}{1}{1}\sdr{27}{21}{1}{1}\sdr{18}{22}{1}{1}\sdr{19}{20}{1}{1}\sdr{26}{20}{1}{1}\sdr{22}{3}{2}{1}\sdr{12}{2}{2}{1}\sdr{22}{2}{3}{1}}%
\sdlevel{70}{\sdr{24}{26}{1}{2}\sdr{24}{24}{1}{1}\sdr{19}{25}{1}{1}\sdr{24}{22}{1}{1}\sdr{22}{22}{1}{3}\sdr{24}{19}{1}{1}\sdr{10}{26}{1}{1}\sdr{12}{26}{1}{1}\sdr{21}{26}{1}{1}\sdr{3}{25}{2}{1}\sdr{19}{21}{1}{1}\sdr{21}{25}{6}{1}\sdr{22}{18}{1}{1}\sdr{18}{23}{1}{1}\sdr{27}{22}{1}{2}\sdr{21}{21}{2}{1}\sdr{24}{21}{3}{1}\sdr{4}{20}{1}{1}\sdr{14}{20}{1}{1}\sdr{20}{20}{6}{1}\sdr{5}{18}{1}{1}}%
\sdlevel{75}{\sdr{25}{24}{2}{1}\sdr{3}{24}{2}{1}\sdr{20}{25}{1}{1}\sdr{5}{27}{1}{1}\sdr{22}{27}{1}{1}\sdr{5}{19}{1}{1}\sdr{22}{19}{1}{1}\sdr{11}{26}{1}{1}\sdr{22}{26}{2}{1}\sdr{20}{21}{1}{1}\sdr{3}{21}{2}{1}\sdr{19}{24}{3}{1}\sdr{25}{22}{2}{1}\sdr{21}{22}{1}{2}\sdr{24}{23}{1}{1}\sdr{4}{22}{1}{1}\sdr{19}{22}{1}{1}\sdr{23}{18}{1}{1}}%
\sdlevel{80}{\sdr{23}{27}{1}{1}\sdr{23}{19}{1}{1}\sdr{12}{25}{1}{1}\sdr{14}{25}{1}{1}\sdr{4}{23}{1}{1}\sdr{19}{23}{2}{1}\sdr{25}{23}{2}{1}\sdr{3}{22}{1}{1}\sdr{20}{22}{1}{1}\sdr{12}{20}{2}{1}\sdr{6}{17}{3}{1}\sdr{11}{17}{1}{1}}%
\sdlevel{85}{\sdr{6}{27}{1}{2}\sdr{6}{18}{1}{2}\sdr{13}{25}{1}{1}\sdr{12}{24}{1}{1}\sdr{3}{23}{1}{1}\sdr{12}{21}{3}{1}\sdr{9}{17}{2}{1}}%
\sdlevel{90}{\sdr{7}{27}{3}{1}\sdr{7}{28}{5}{1}\sdr{7}{18}{3}{2}\sdr{13}{24}{2}{1}\sdr{12}{22}{1}{2}}%
\sdlevel{95}{\sdr{10}{27}{2}{1}\sdr{10}{18}{2}{2}\sdr{13}{22}{2}{2}}%
}
\def\sdimgD{%
\sdlevel{5}{\sdr{2}{28}{1}{1}\sdr{2}{19}{1}{1}\sdr{14}{28}{1}{1}\sdr{14}{17}{1}{1}\sdr{3}{29}{2}{1}\sdr{2}{17}{1}{1}\sdr{14}{4}{1}{1}\sdr{7}{23}{1}{1}\sdr{6}{22}{2}{1}\sdr{12}{6}{3}{1}\sdr{17}{4}{1}{3}\sdr{13}{5}{2}{1}}%
\sdlevel{10}{\sdr{0}{29}{2}{3}\sdr{0}{24}{2}{4}\sdr{8}{24}{2}{1}\sdr{0}{28}{1}{1}\sdr{0}{18}{2}{4}\sdr{5}{30}{2}{2}\sdr{5}{29}{1}{1}\sdr{0}{4}{2}{13}\sdr{5}{14}{2}{3}\sdr{0}{0}{2}{2}\sdr{5}{0}{2}{2}\sdr{5}{11}{1}{1}\sdr{0}{17}{1}{1}\sdr{8}{21}{2}{1}\sdr{6}{23}{1}{1}\sdr{8}{22}{1}{2}\sdr{3}{12}{3}{2}\sdr{3}{11}{1}{1}\sdr{18}{4}{1}{3}\sdr{12}{5}{1}{1}}%
\sdlevel{15}{\sdr{2}{29}{1}{1}\sdr{7}{30}{1}{2}\sdr{1}{28}{1}{1}\sdr{2}{30}{3}{2}\sdr{7}{24}{1}{1}\sdr{2}{26}{1}{1}\sdr{19}{28}{1}{1}\sdr{0}{22}{2}{2}\sdr{2}{14}{3}{3}\sdr{7}{14}{1}{3}\sdr{2}{0}{3}{2}\sdr{7}{0}{1}{2}\sdr{2}{11}{1}{3}\sdr{12}{29}{3}{1}\sdr{1}{17}{1}{1}\sdr{19}{4}{1}{3}\sdr{8}{25}{2}{1}\sdr{0}{2}{2}{1}\sdr{9}{22}{2}{2}\sdr{6}{13}{1}{1}\sdr{6}{11}{2}{2}\sdr{4}{11}{1}{1}\sdr{2}{8}{2}{3}\sdr{5}{10}{1}{1}\sdr{7}{10}{3}{1}\sdr{3}{6}{3}{1}\sdr{3}{5}{2}{1}\sdr{13}{4}{1}{1}\sdr{0}{3}{3}{1}}%
\sdlevel{20}{\sdr{10}{24}{1}{1}\sdr{15}{25}{2}{1}\sdr{18}{28}{1}{1}\sdr{11}{30}{2}{2}\sdr{7}{21}{1}{1}\sdr{10}{21}{1}{1}\sdr{4}{17}{1}{1}\sdr{10}{12}{1}{1}\sdr{15}{26}{3}{1}\sdr{2}{7}{1}{1}\sdr{5}{8}{1}{1}\sdr{2}{5}{1}{1}\sdr{11}{14}{2}{3}\sdr{11}{0}{2}{2}\sdr{18}{18}{1}{1}\sdr{15}{7}{3}{1}\sdr{15}{20}{2}{2}\sdr{10}{10}{1}{1}\sdr{7}{13}{1}{1}\sdr{3}{18}{1}{1}\sdr{14}{18}{1}{1}\sdr{4}{10}{1}{1}\sdr{13}{17}{1}{1}\sdr{19}{17}{1}{1}\sdr{8}{12}{1}{1}\sdr{8}{11}{3}{1}\sdr{6}{10}{1}{1}\sdr{5}{9}{5}{1}\sdr{5}{5}{1}{1}\sdr{2}{2}{1}{1}\sdr{2}{4}{2}{1}\sdr{12}{4}{1}{1}}%
\sdlevel{25}{\sdr{6}{29}{1}{1}\sdr{15}{28}{1}{1}\sdr{15}{24}{1}{1}\sdr{6}{24}{1}{1}\sdr{17}{28}{1}{1}\sdr{15}{29}{3}{3}\sdr{19}{29}{1}{1}\sdr{18}{27}{1}{1}\sdr{8}{30}{3}{2}\sdr{13}{30}{1}{2}\sdr{15}{18}{3}{1}\sdr{2}{20}{1}{1}\sdr{6}{21}{1}{1}\sdr{4}{8}{1}{2}\sdr{3}{7}{1}{1}\sdr{8}{14}{3}{3}\sdr{13}{14}{1}{3}\sdr{15}{14}{3}{3}\sdr{8}{0}{3}{2}\sdr{13}{0}{1}{2}\sdr{15}{0}{3}{2}\sdr{6}{6}{1}{1}\sdr{19}{18}{1}{1}\sdr{15}{22}{1}{1}\sdr{17}{20}{1}{1}\sdr{18}{17}{1}{1}\sdr{2}{6}{1}{1}\sdr{8}{20}{2}{1}\sdr{17}{19}{2}{1}\sdr{12}{17}{1}{1}\sdr{8}{13}{5}{1}\sdr{9}{12}{1}{1}\sdr{11}{10}{1}{3}\sdr{22}{11}{1}{2}\sdr{4}{4}{1}{1}\sdr{10}{9}{1}{1}\sdr{6}{8}{3}{1}\sdr{3}{3}{1}{1}\sdr{5}{7}{2}{1}}%
\sdlevel{30}{\sdr{18}{29}{1}{1}\sdr{16}{28}{1}{1}\sdr{16}{24}{1}{1}\sdr{17}{25}{1}{1}\sdr{29}{30}{3}{2}\sdr{19}{27}{1}{1}\sdr{14}{30}{1}{2}\sdr{18}{30}{2}{2}\sdr{21}{30}{5}{2}\sdr{29}{20}{3}{1}\sdr{13}{28}{1}{1}\sdr{11}{22}{1}{2}\sdr{7}{25}{1}{1}\sdr{18}{14}{2}{3}\sdr{14}{14}{1}{3}\sdr{18}{7}{2}{1}\sdr{21}{14}{5}{3}\sdr{29}{14}{3}{3}\sdr{14}{0}{1}{2}\sdr{18}{0}{2}{2}\sdr{21}{0}{5}{2}\sdr{29}{0}{3}{2}\sdr{18}{26}{1}{1}\sdr{13}{13}{1}{1}\sdr{16}{22}{1}{1}\sdr{26}{28}{2}{1}\sdr{3}{27}{1}{1}\sdr{14}{27}{4}{1}\sdr{19}{19}{1}{1}\sdr{28}{26}{4}{1}\sdr{2}{25}{1}{1}\sdr{7}{7}{1}{1}\sdr{10}{25}{1}{1}\sdr{17}{21}{1}{1}\sdr{11}{9}{1}{1}\sdr{15}{19}{2}{1}\sdr{27}{18}{2}{1}\sdr{15}{17}{3}{1}\sdr{20}{17}{1}{1}\sdr{26}{17}{1}{1}\sdr{20}{13}{7}{1}\sdr{12}{12}{1}{1}\sdr{20}{11}{2}{2}\sdr{23}{11}{3}{2}\sdr{21}{10}{4}{1}\sdr{22}{9}{1}{1}\sdr{9}{8}{2}{1}\sdr{4}{7}{1}{1}\sdr{12}{7}{3}{1}\sdr{6}{5}{1}{1}}%
\sdlevel{35}{\sdr{20}{30}{1}{2}\sdr{2}{24}{1}{1}\sdr{29}{25}{3}{1}\sdr{28}{28}{1}{1}\sdr{20}{28}{1}{1}\sdr{20}{18}{1}{1}\sdr{26}{30}{3}{2}\sdr{14}{19}{1}{1}\sdr{26}{18}{1}{1}\sdr{20}{14}{1}{3}\sdr{20}{29}{2}{1}\sdr{29}{21}{3}{1}\sdr{7}{29}{1}{1}\sdr{25}{17}{1}{1}\sdr{29}{18}{1}{1}\sdr{29}{6}{1}{1}\sdr{3}{19}{1}{1}\sdr{20}{10}{1}{1}\sdr{26}{14}{3}{3}\sdr{20}{0}{1}{2}\sdr{26}{0}{3}{2}\sdr{7}{6}{1}{1}\sdr{20}{9}{2}{1}\sdr{24}{29}{3}{1}\sdr{4}{28}{1}{1}\sdr{28}{20}{1}{1}\sdr{27}{27}{3}{1}\sdr{29}{7}{3}{1}\sdr{2}{21}{1}{1}\sdr{17}{24}{1}{1}\sdr{15}{23}{2}{1}\sdr{3}{2}{1}{1}\sdr{14}{13}{1}{1}\sdr{27}{19}{2}{1}\sdr{26}{11}{1}{1}\sdr{29}{4}{1}{1}\sdr{25}{6}{1}{1}\sdr{27}{17}{1}{1}\sdr{19}{13}{1}{1}\sdr{27}{13}{5}{1}\sdr{13}{12}{1}{1}\sdr{26}{12}{6}{1}\sdr{29}{10}{2}{2}\sdr{25}{10}{2}{1}\sdr{23}{9}{3}{1}\sdr{11}{8}{1}{1}\sdr{20}{8}{4}{1}\sdr{10}{7}{1}{1}\sdr{5}{4}{1}{1}}%
\sdlevel{40}{\sdr{11}{29}{1}{1}\sdr{11}{24}{1}{1}\sdr{26}{27}{1}{1}\sdr{29}{24}{3}{1}\sdr{30}{27}{2}{1}\sdr{28}{25}{1}{1}\sdr{19}{26}{1}{1}\sdr{12}{28}{1}{1}\sdr{11}{21}{1}{1}\sdr{30}{18}{2}{1}\sdr{25}{7}{4}{1}\sdr{30}{6}{2}{1}\sdr{11}{6}{1}{2}\sdr{22}{29}{2}{1}\sdr{27}{29}{5}{1}\sdr{12}{11}{1}{1}\sdr{25}{28}{1}{1}\sdr{29}{28}{1}{1}\sdr{20}{27}{1}{1}\sdr{26}{19}{1}{1}\sdr{30}{4}{2}{1}\sdr{19}{12}{1}{1}\sdr{27}{26}{1}{1}\sdr{6}{25}{1}{1}\sdr{28}{17}{1}{1}\sdr{29}{19}{3}{1}\sdr{5}{23}{1}{1}\sdr{17}{22}{1}{1}\sdr{7}{20}{1}{1}\sdr{10}{20}{1}{1}\sdr{18}{20}{1}{1}\sdr{21}{17}{1}{1}\sdr{15}{13}{4}{1}\sdr{27}{10}{2}{2}\sdr{31}{10}{1}{2}\sdr{26}{9}{6}{1}\sdr{24}{8}{8}{1}\sdr{8}{7}{2}{1}\sdr{25}{4}{4}{1}\sdr{24}{6}{1}{1}\sdr{26}{6}{3}{1}\sdr{25}{5}{7}{1}}%
\sdlevel{45}{\sdr{18}{25}{1}{1}\sdr{5}{22}{1}{1}\sdr{24}{17}{1}{1}\sdr{10}{29}{1}{1}\sdr{24}{7}{1}{1}\sdr{20}{19}{1}{1}\sdr{10}{6}{1}{1}\sdr{21}{28}{1}{1}\sdr{30}{28}{2}{1}\sdr{18}{21}{1}{1}\sdr{2}{23}{1}{1}\sdr{17}{23}{1}{1}\sdr{5}{17}{1}{1}\sdr{28}{21}{1}{1}\sdr{20}{3}{1}{2}\sdr{25}{18}{1}{1}\sdr{24}{5}{1}{1}\sdr{29}{17}{3}{1}\sdr{14}{12}{5}{1}\sdr{13}{11}{1}{1}\sdr{19}{11}{1}{1}\sdr{12}{10}{1}{1}\sdr{20}{7}{3}{1}\sdr{20}{5}{2}{2}\sdr{7}{5}{1}{1}\sdr{11}{5}{1}{1}\sdr{6}{4}{1}{1}\sdr{4}{3}{1}{1}\sdr{26}{3}{5}{1}\sdr{29}{2}{1}{1}}%
\sdlevel{50}{\sdr{8}{29}{2}{1}\sdr{14}{26}{1}{1}\sdr{22}{17}{1}{1}\sdr{24}{28}{1}{1}\sdr{8}{26}{2}{1}\sdr{24}{4}{1}{1}\sdr{14}{11}{1}{1}\sdr{26}{26}{1}{1}\sdr{11}{25}{1}{1}\sdr{28}{24}{1}{1}\sdr{29}{23}{2}{1}\sdr{2}{22}{1}{1}\sdr{29}{22}{3}{1}\sdr{27}{20}{1}{1}\sdr{22}{5}{1}{1}\sdr{18}{11}{1}{1}\sdr{13}{10}{1}{1}\sdr{19}{8}{1}{3}\sdr{12}{8}{2}{2}\sdr{23}{7}{1}{1}\sdr{22}{6}{2}{1}\sdr{10}{5}{1}{1}\sdr{7}{4}{1}{1}\sdr{21}{3}{1}{2}\sdr{5}{3}{1}{1}\sdr{7}{3}{3}{1}\sdr{17}{3}{3}{1}\sdr{25}{3}{1}{1}\sdr{31}{3}{1}{1}\sdr{4}{2}{6}{1}\sdr{15}{2}{6}{1}\sdr{26}{2}{3}{1}\sdr{30}{2}{2}{1}}%
\sdlevel{55}{\sdr{25}{27}{1}{1}\sdr{6}{20}{1}{1}\sdr{14}{10}{1}{1}\sdr{21}{18}{1}{1}\sdr{25}{19}{1}{1}\sdr{25}{2}{1}{1}\sdr{20}{26}{1}{1}\sdr{27}{25}{1}{1}\sdr{18}{24}{1}{1}\sdr{31}{23}{1}{1}\sdr{6}{3}{1}{1}\sdr{19}{20}{1}{1}\sdr{13}{18}{1}{1}\sdr{21}{2}{1}{1}\sdr{15}{11}{3}{1}\sdr{14}{2}{1}{1}\sdr{17}{9}{2}{2}\sdr{14}{9}{2}{1}\sdr{14}{8}{5}{1}\sdr{8}{6}{1}{1}\sdr{23}{5}{1}{1}\sdr{11}{4}{1}{1}\sdr{10}{2}{1}{2}\sdr{14}{3}{3}{1}\sdr{22}{3}{1}{1}\sdr{24}{3}{1}{1}}%
\sdlevel{60}{\sdr{27}{24}{1}{1}\sdr{23}{17}{1}{1}\sdr{5}{28}{1}{1}\sdr{22}{28}{1}{1}\sdr{21}{27}{1}{1}\sdr{19}{25}{1}{1}\sdr{27}{21}{1}{1}\sdr{18}{22}{1}{1}\sdr{26}{20}{1}{1}\sdr{24}{18}{1}{1}\sdr{23}{3}{1}{1}\sdr{15}{10}{2}{1}\sdr{16}{9}{1}{1}\sdr{9}{6}{1}{1}\sdr{8}{5}{1}{1}\sdr{10}{4}{1}{1}\sdr{22}{4}{2}{1}\sdr{13}{3}{1}{1}\sdr{11}{2}{1}{1}\sdr{22}{2}{3}{1}}%
\sdlevel{65}{\sdr{23}{28}{1}{1}\sdr{13}{27}{1}{1}\sdr{3}{26}{1}{1}\sdr{7}{26}{1}{1}\sdr{10}{26}{1}{1}\sdr{25}{26}{1}{1}\sdr{26}{25}{1}{1}\sdr{5}{24}{1}{1}\sdr{18}{23}{1}{1}\sdr{28}{22}{1}{2}\sdr{19}{21}{1}{1}\sdr{11}{20}{1}{1}\sdr{20}{20}{1}{1}\sdr{21}{19}{1}{1}\sdr{4}{18}{1}{1}\sdr{12}{18}{1}{1}\sdr{9}{5}{1}{1}\sdr{8}{4}{1}{1}\sdr{11}{3}{2}{1}\sdr{12}{2}{2}{1}}%
\sdlevel{70}{\sdr{19}{24}{1}{1}\sdr{23}{23}{1}{2}\sdr{26}{24}{1}{1}\sdr{24}{27}{1}{1}\sdr{24}{19}{1}{1}\sdr{21}{26}{1}{1}\sdr{20}{25}{1}{1}\sdr{19}{22}{1}{1}\sdr{23}{21}{1}{1}\sdr{26}{21}{1}{1}\sdr{5}{21}{1}{1}\sdr{14}{20}{1}{1}\sdr{25}{20}{1}{1}\sdr{13}{19}{1}{1}\sdr{22}{18}{1}{1}\sdr{6}{17}{1}{1}\sdr{9}{4}{1}{1}}%
\sdlevel{75}{\sdr{22}{27}{1}{1}\sdr{24}{26}{1}{1}\sdr{24}{24}{2}{1}\sdr{22}{23}{1}{2}\sdr{22}{19}{1}{1}\sdr{24}{23}{1}{1}\sdr{21}{25}{5}{1}\sdr{20}{24}{1}{1}\sdr{24}{21}{2}{1}\sdr{4}{22}{1}{2}\sdr{19}{23}{1}{1}\sdr{27}{22}{1}{2}\sdr{22}{22}{3}{1}\sdr{20}{21}{3}{1}\sdr{21}{20}{4}{1}\sdr{23}{18}{1}{1}}%
\sdlevel{80}{\sdr{12}{27}{1}{1}\sdr{14}{25}{1}{1}\sdr{4}{27}{1}{1}\sdr{8}{27}{2}{1}\sdr{12}{22}{1}{2}\sdr{23}{27}{1}{1}\sdr{4}{19}{1}{1}\sdr{8}{19}{2}{1}\sdr{12}{19}{1}{1}\sdr{23}{19}{1}{1}\sdr{6}{26}{1}{1}\sdr{22}{26}{2}{1}\sdr{5}{25}{1}{1}\sdr{14}{21}{1}{1}\sdr{21}{23}{1}{2}\sdr{25}{22}{1}{2}\sdr{7}{17}{1}{1}\sdr{11}{17}{1}{1}}%
\sdlevel{85}{\sdr{13}{26}{1}{1}\sdr{3}{23}{1}{1}\sdr{6}{28}{1}{1}\sdr{11}{26}{1}{1}\sdr{13}{23}{1}{1}\sdr{14}{24}{1}{1}\sdr{3}{20}{1}{1}\sdr{20}{23}{1}{1}\sdr{26}{22}{1}{2}\sdr{13}{22}{2}{1}\sdr{21}{22}{1}{1}\sdr{8}{18}{2}{1}\sdr{8}{17}{1}{1}\sdr{10}{17}{1}{1}}%
\sdlevel{90}{\sdr{3}{24}{1}{2}\sdr{10}{27}{1}{1}\sdr{5}{20}{1}{1}\sdr{7}{28}{5}{1}\sdr{10}{19}{1}{1}\sdr{3}{21}{1}{2}\sdr{14}{23}{1}{1}\sdr{20}{22}{1}{1}\sdr{5}{18}{1}{1}\sdr{7}{18}{1}{2}\sdr{9}{17}{1}{1}}%
\sdlevel{95}{\sdr{4}{24}{1}{1}\sdr{12}{24}{2}{1}\sdr{13}{25}{1}{1}\sdr{7}{27}{1}{1}\sdr{12}{26}{1}{1}\sdr{13}{20}{1}{1}\sdr{4}{21}{1}{1}\sdr{12}{21}{2}{1}\sdr{6}{18}{1}{1}\sdr{10}{18}{1}{1}}%
\sdlevel{100}{\sdr{4}{25}{1}{1}\sdr{12}{25}{1}{1}\sdr{5}{27}{2}{1}\sdr{11}{27}{1}{1}\sdr{5}{19}{2}{1}\sdr{11}{18}{1}{2}\sdr{4}{26}{2}{1}\sdr{4}{20}{1}{1}\sdr{12}{20}{1}{1}}%
}
\def\sdimgE{%
\sdlevel{5}{\sdr{2}{25}{1}{1}\sdr{2}{13}{1}{1}\sdr{2}{4}{1}{1}\sdr{6}{7}{1}{1}\sdr{4}{6}{1}{1}\sdr{12}{6}{3}{1}\sdr{12}{5}{2}{1}\sdr{12}{4}{6}{1}\sdr{19}{4}{1}{1}}%
\sdlevel{10}{\sdr{0}{30}{3}{2}\sdr{0}{28}{2}{2}\sdr{7}{24}{1}{2}\sdr{10}{24}{1}{1}\sdr{0}{27}{3}{1}\sdr{0}{21}{2}{6}\sdr{3}{28}{1}{1}\sdr{0}{18}{3}{3}\sdr{0}{17}{2}{1}\sdr{0}{14}{3}{3}\sdr{0}{13}{2}{1}\sdr{3}{17}{1}{1}\sdr{0}{12}{3}{1}\sdr{0}{10}{2}{2}\sdr{3}{10}{1}{1}\sdr{5}{23}{1}{1}\sdr{0}{9}{3}{1}\sdr{0}{7}{2}{1}\sdr{4}{12}{1}{1}\sdr{0}{8}{1}{1}\sdr{0}{4}{1}{1}\sdr{0}{6}{3}{1}\sdr{0}{0}{3}{2}\sdr{0}{5}{2}{1}\sdr{6}{29}{1}{1}\sdr{3}{7}{1}{1}\sdr{7}{21}{1}{1}\sdr{10}{21}{1}{1}\sdr{5}{7}{1}{1}\sdr{3}{13}{2}{1}\sdr{4}{5}{1}{1}\sdr{0}{2}{1}{1}\sdr{2}{8}{1}{1}\sdr{14}{7}{4}{1}\sdr{18}{6}{1}{1}\sdr{18}{5}{2}{1}}%
\sdlevel{15}{\sdr{2}{29}{1}{1}\sdr{8}{26}{2}{1}\sdr{4}{29}{1}{1}\sdr{3}{30}{4}{2}\sdr{6}{24}{1}{1}\sdr{8}{20}{2}{1}\sdr{8}{13}{2}{1}\sdr{2}{21}{1}{1}\sdr{5}{22}{3}{1}\sdr{8}{25}{1}{1}\sdr{2}{10}{1}{1}\sdr{3}{12}{1}{1}\sdr{1}{8}{1}{1}\sdr{1}{4}{1}{1}\sdr{3}{14}{4}{3}\sdr{3}{0}{4}{2}\sdr{2}{7}{1}{1}\sdr{4}{17}{1}{1}\sdr{13}{29}{1}{1}\sdr{18}{28}{2}{1}\sdr{8}{11}{2}{1}\sdr{8}{8}{1}{1}\sdr{6}{13}{1}{1}\sdr{5}{10}{3}{1}\sdr{11}{13}{1}{1}\sdr{3}{6}{1}{1}\sdr{9}{12}{2}{1}\sdr{2}{11}{4}{1}\sdr{4}{9}{2}{1}\sdr{7}{9}{2}{1}\sdr{1}{2}{1}{1}\sdr{3}{8}{3}{1}\sdr{19}{6}{1}{1}\sdr{2}{5}{2}{1}\sdr{0}{3}{3}{1}}%
\sdlevel{20}{\sdr{2}{26}{1}{1}\sdr{15}{29}{2}{1}\sdr{8}{24}{2}{1}\sdr{3}{27}{1}{1}\sdr{14}{27}{1}{1}\sdr{7}{30}{3}{2}\sdr{11}{30}{2}{2}\sdr{5}{29}{1}{1}\sdr{7}{29}{1}{1}\sdr{12}{29}{1}{1}\sdr{11}{14}{2}{3}\sdr{3}{18}{1}{2}\sdr{6}{21}{1}{1}\sdr{9}{25}{1}{1}\sdr{6}{11}{2}{1}\sdr{4}{10}{1}{1}\sdr{7}{13}{1}{1}\sdr{5}{13}{1}{1}\sdr{4}{7}{1}{1}\sdr{2}{24}{1}{1}\sdr{7}{14}{3}{3}\sdr{11}{12}{2}{1}\sdr{7}{0}{3}{2}\sdr{11}{0}{2}{2}\sdr{5}{5}{1}{1}\sdr{7}{7}{1}{1}\sdr{12}{13}{1}{1}\sdr{15}{22}{2}{1}\sdr{13}{28}{1}{1}\sdr{3}{9}{1}{1}\sdr{14}{19}{1}{1}\sdr{2}{2}{1}{1}\sdr{15}{25}{1}{2}\sdr{9}{9}{1}{1}\sdr{8}{21}{2}{1}\sdr{6}{23}{6}{1}\sdr{8}{22}{4}{1}\sdr{6}{9}{1}{1}\sdr{12}{17}{3}{1}\sdr{19}{17}{1}{1}\sdr{5}{12}{4}{1}\sdr{6}{8}{2}{1}\sdr{10}{11}{2}{1}\sdr{4}{4}{1}{1}\sdr{8}{10}{3}{1}\sdr{5}{6}{2}{1}}%
\sdlevel{25}{\sdr{10}{30}{1}{2}\sdr{15}{27}{3}{1}\sdr{15}{21}{3}{1}\sdr{10}{25}{1}{1}\sdr{13}{30}{2}{2}\sdr{17}{30}{3}{2}\sdr{11}{29}{1}{1}\sdr{10}{13}{1}{4}\sdr{13}{14}{2}{3}\sdr{17}{14}{3}{3}\sdr{10}{0}{1}{2}\sdr{13}{0}{2}{2}\sdr{17}{0}{3}{2}\sdr{11}{10}{1}{1}\sdr{17}{29}{1}{1}\sdr{20}{29}{1}{1}\sdr{25}{29}{1}{1}\sdr{14}{28}{1}{1}\sdr{15}{19}{3}{1}\sdr{16}{26}{4}{1}\sdr{16}{25}{2}{1}\sdr{15}{20}{2}{1}\sdr{18}{20}{1}{1}\sdr{14}{18}{7}{1}\sdr{2}{17}{1}{1}\sdr{5}{17}{1}{1}\sdr{13}{12}{1}{2}\sdr{21}{12}{2}{1}\sdr{22}{11}{1}{1}\sdr{9}{8}{1}{1}\sdr{3}{3}{1}{2}}%
\sdlevel{30}{\sdr{14}{29}{1}{1}\sdr{15}{28}{3}{1}\sdr{15}{30}{2}{2}\sdr{11}{24}{1}{1}\sdr{15}{23}{2}{1}\sdr{18}{29}{1}{1}\sdr{21}{29}{1}{1}\sdr{23}{30}{1}{2}\sdr{14}{26}{1}{1}\sdr{10}{29}{1}{1}\sdr{24}{29}{1}{1}\sdr{15}{14}{2}{4}\sdr{23}{14}{1}{3}\sdr{15}{0}{2}{2}\sdr{23}{0}{1}{2}\sdr{10}{9}{1}{1}\sdr{14}{13}{1}{1}\sdr{18}{17}{1}{1}\sdr{21}{17}{1}{1}\sdr{24}{8}{1}{1}\sdr{4}{28}{1}{1}\sdr{12}{28}{1}{1}\sdr{15}{24}{3}{1}\sdr{27}{28}{1}{1}\sdr{26}{27}{1}{1}\sdr{28}{25}{1}{1}\sdr{11}{21}{1}{1}\sdr{17}{20}{1}{1}\sdr{29}{20}{1}{1}\sdr{25}{17}{3}{1}\sdr{18}{13}{6}{1}\sdr{20}{12}{1}{1}\sdr{23}{11}{2}{2}\sdr{20}{11}{2}{1}\sdr{20}{10}{5}{1}\sdr{22}{9}{3}{1}\sdr{8}{7}{1}{1}\sdr{26}{7}{2}{1}\sdr{7}{6}{1}{1}\sdr{6}{5}{1}{1}}%
\sdlevel{35}{\sdr{28}{29}{4}{1}\sdr{28}{28}{1}{1}\sdr{18}{27}{3}{1}\sdr{28}{27}{4}{1}\sdr{20}{28}{1}{1}\sdr{20}{30}{3}{2}\sdr{24}{30}{8}{2}\sdr{28}{19}{1}{2}\sdr{29}{22}{1}{1}\sdr{25}{12}{1}{1}\sdr{25}{9}{1}{1}\sdr{27}{9}{2}{1}\sdr{20}{14}{3}{3}\sdr{24}{14}{8}{3}\sdr{20}{0}{3}{2}\sdr{24}{0}{8}{2}\sdr{28}{18}{4}{1}\sdr{20}{17}{1}{1}\sdr{28}{12}{1}{2}\sdr{18}{19}{3}{1}\sdr{27}{26}{1}{1}\sdr{6}{25}{1}{1}\sdr{29}{25}{2}{1}\sdr{29}{11}{1}{1}\sdr{7}{20}{1}{1}\sdr{10}{20}{1}{1}\sdr{30}{20}{2}{1}\sdr{26}{18}{1}{2}\sdr{17}{17}{1}{1}\sdr{24}{17}{1}{1}\sdr{29}{17}{3}{1}\sdr{24}{13}{3}{1}\sdr{18}{12}{2}{1}\sdr{25}{6}{1}{2}\sdr{12}{11}{1}{1}\sdr{25}{10}{2}{2}\sdr{28}{10}{3}{1}\sdr{27}{6}{2}{1}\sdr{22}{8}{2}{1}\sdr{25}{8}{3}{1}\sdr{23}{7}{1}{1}\sdr{28}{7}{2}{1}}%
\sdlevel{40}{\sdr{29}{28}{3}{1}\sdr{18}{25}{1}{1}\sdr{28}{26}{4}{1}\sdr{27}{27}{1}{1}\sdr{29}{23}{3}{2}\sdr{27}{18}{1}{3}\sdr{29}{19}{3}{1}\sdr{30}{22}{2}{1}\sdr{27}{13}{1}{1}\sdr{31}{25}{1}{1}\sdr{21}{28}{1}{1}\sdr{26}{29}{1}{1}\sdr{29}{12}{3}{2}\sdr{30}{11}{2}{1}\sdr{26}{9}{1}{1}\sdr{29}{9}{3}{1}\sdr{27}{11}{2}{1}\sdr{19}{29}{1}{1}\sdr{22}{29}{1}{1}\sdr{26}{6}{1}{1}\sdr{21}{9}{1}{1}\sdr{25}{28}{2}{1}\sdr{29}{6}{3}{1}\sdr{27}{10}{1}{1}\sdr{28}{21}{4}{1}\sdr{18}{21}{1}{1}\sdr{31}{10}{1}{1}\sdr{30}{7}{2}{1}\sdr{17}{13}{1}{1}\sdr{26}{12}{2}{1}\sdr{13}{11}{1}{1}\sdr{27}{4}{2}{1}\sdr{10}{8}{1}{1}\sdr{28}{8}{3}{1}\sdr{24}{6}{1}{2}\sdr{25}{5}{6}{1}\sdr{3}{2}{1}{1}}%
\sdlevel{45}{\sdr{27}{29}{1}{1}\sdr{28}{24}{1}{1}\sdr{24}{28}{1}{1}\sdr{31}{8}{1}{1}\sdr{26}{26}{1}{1}\sdr{20}{9}{1}{1}\sdr{29}{4}{3}{1}\sdr{23}{29}{1}{1}\sdr{27}{25}{1}{1}\sdr{24}{5}{1}{1}\sdr{26}{4}{1}{1}\sdr{28}{17}{1}{1}\sdr{17}{22}{1}{1}\sdr{25}{18}{1}{1}\sdr{22}{17}{1}{1}\sdr{15}{13}{2}{1}\sdr{19}{11}{1}{1}\sdr{20}{3}{1}{1}\sdr{21}{8}{1}{1}\sdr{31}{5}{1}{1}\sdr{18}{7}{1}{1}\sdr{7}{5}{1}{1}\sdr{29}{2}{3}{1}\sdr{28}{3}{3}{1}}%
\sdlevel{50}{\sdr{18}{24}{1}{1}\sdr{14}{11}{1}{2}\sdr{17}{12}{1}{1}\sdr{19}{20}{1}{1}\sdr{10}{6}{1}{1}\sdr{21}{18}{1}{1}\sdr{25}{27}{1}{1}\sdr{18}{11}{1}{1}\sdr{19}{7}{1}{1}\sdr{21}{3}{1}{1}\sdr{14}{8}{1}{1}\sdr{17}{8}{1}{1}\sdr{12}{10}{1}{1}\sdr{11}{9}{2}{1}\sdr{20}{8}{1}{1}\sdr{9}{7}{1}{1}\sdr{22}{7}{1}{1}\sdr{8}{5}{1}{2}\sdr{10}{3}{1}{1}\sdr{20}{6}{4}{1}\sdr{11}{5}{1}{1}\sdr{20}{4}{2}{2}\sdr{5}{4}{2}{1}\sdr{24}{4}{2}{1}\sdr{4}{3}{1}{1}\sdr{7}{2}{1}{2}\sdr{18}{3}{2}{1}\sdr{25}{3}{3}{1}\sdr{31}{3}{1}{1}\sdr{9}{2}{2}{1}\sdr{15}{2}{2}{1}\sdr{18}{2}{4}{1}\sdr{27}{2}{2}{1}}%
\sdlevel{55}{\sdr{15}{12}{2}{1}\sdr{11}{6}{1}{1}\sdr{22}{28}{1}{1}\sdr{17}{23}{1}{1}\sdr{22}{3}{1}{1}\sdr{17}{2}{1}{1}\sdr{25}{19}{1}{1}\sdr{13}{18}{1}{1}\sdr{15}{8}{2}{1}\sdr{15}{11}{3}{1}\sdr{13}{10}{7}{1}\sdr{15}{9}{5}{1}\sdr{12}{7}{1}{2}\sdr{19}{8}{1}{1}\sdr{10}{7}{1}{1}\sdr{20}{7}{1}{1}\sdr{9}{6}{1}{1}\sdr{11}{3}{1}{1}\sdr{9}{5}{2}{1}\sdr{22}{5}{2}{1}\sdr{7}{4}{4}{1}\sdr{5}{3}{2}{1}\sdr{8}{3}{2}{1}\sdr{14}{3}{4}{1}\sdr{24}{3}{1}{1}\sdr{4}{2}{2}{1}\sdr{8}{2}{1}{1}\sdr{25}{2}{2}{1}}%
\sdlevel{60}{\sdr{11}{7}{1}{2}\sdr{13}{7}{1}{1}\sdr{13}{9}{2}{1}\sdr{11}{4}{1}{1}\sdr{18}{8}{1}{1}\sdr{13}{3}{1}{1}\sdr{22}{4}{2}{1}\sdr{23}{3}{1}{1}\sdr{6}{2}{1}{1}\sdr{11}{2}{3}{1}\sdr{22}{2}{1}{1}\sdr{24}{2}{1}{1}}%
\sdlevel{65}{\sdr{27}{24}{1}{1}\sdr{21}{27}{1}{1}\sdr{21}{19}{1}{1}\sdr{23}{17}{1}{1}\sdr{5}{28}{1}{1}\sdr{21}{7}{1}{1}\sdr{27}{21}{1}{1}\sdr{19}{21}{1}{1}\sdr{24}{18}{1}{1}\sdr{23}{2}{1}{1}\sdr{13}{8}{1}{1}\sdr{12}{3}{1}{1}\sdr{14}{2}{1}{1}}%
\sdlevel{70}{\sdr{24}{27}{1}{1}\sdr{23}{28}{1}{1}\sdr{20}{26}{1}{1}\sdr{23}{24}{1}{1}\sdr{24}{25}{1}{1}\sdr{23}{21}{1}{1}\sdr{24}{20}{1}{1}\sdr{3}{26}{1}{1}\sdr{20}{21}{1}{1}\sdr{25}{26}{1}{1}\sdr{19}{25}{1}{1}\sdr{19}{24}{2}{1}\sdr{21}{23}{2}{1}\sdr{28}{22}{1}{1}\sdr{6}{20}{1}{1}\sdr{21}{20}{1}{1}\sdr{26}{20}{1}{1}\sdr{12}{18}{1}{1}\sdr{22}{18}{1}{1}}%
\sdlevel{75}{\sdr{22}{27}{2}{1}\sdr{22}{25}{2}{1}\sdr{11}{25}{1}{1}\sdr{25}{24}{2}{2}\sdr{21}{24}{2}{1}\sdr{19}{23}{2}{1}\sdr{23}{23}{4}{1}\sdr{28}{23}{1}{1}\sdr{18}{22}{10}{1}\sdr{21}{21}{1}{1}\sdr{26}{21}{1}{1}\sdr{22}{20}{1}{1}\sdr{25}{20}{1}{1}\sdr{23}{19}{2}{1}\sdr{23}{18}{1}{1}}%
\sdlevel{80}{\sdr{13}{27}{1}{1}\sdr{22}{21}{1}{1}\sdr{4}{25}{1}{1}\sdr{8}{29}{2}{1}\sdr{13}{19}{1}{1}\sdr{21}{26}{4}{1}\sdr{4}{18}{1}{1}\sdr{21}{25}{1}{1}\sdr{24}{24}{1}{1}\sdr{18}{23}{1}{1}\sdr{27}{23}{1}{1}\sdr{22}{19}{1}{1}\sdr{24}{21}{2}{1}\sdr{12}{20}{1}{1}\sdr{20}{20}{1}{1}\sdr{23}{20}{1}{1}\sdr{11}{18}{1}{1}\sdr{7}{17}{1}{1}}%
\sdlevel{85}{\sdr{7}{28}{1}{1}\sdr{10}{28}{1}{1}\sdr{3}{24}{1}{1}\sdr{8}{27}{1}{1}\sdr{10}{26}{1}{1}\sdr{7}{26}{1}{1}\sdr{10}{17}{1}{2}\sdr{8}{19}{1}{1}\sdr{20}{25}{1}{1}\sdr{3}{21}{1}{1}\sdr{5}{24}{1}{1}\sdr{2}{22}{1}{2}\sdr{14}{22}{1}{1}\sdr{13}{21}{2}{1}\sdr{4}{20}{1}{1}\sdr{13}{20}{1}{1}\sdr{11}{19}{1}{1}\sdr{6}{17}{1}{1}}%
\sdlevel{90}{\sdr{4}{27}{1}{1}\sdr{12}{26}{2}{1}\sdr{5}{25}{1}{2}\sdr{12}{25}{1}{1}\sdr{11}{27}{1}{1}\sdr{14}{24}{1}{2}\sdr{4}{22}{1}{2}\sdr{6}{27}{1}{2}\sdr{9}{27}{1}{1}\sdr{5}{21}{1}{1}\sdr{6}{19}{1}{1}\sdr{4}{19}{1}{1}\sdr{9}{19}{1}{1}\sdr{11}{20}{1}{1}\sdr{5}{18}{1}{1}\sdr{12}{22}{2}{2}\sdr{3}{25}{1}{1}\sdr{12}{21}{1}{1}\sdr{14}{20}{1}{1}\sdr{8}{17}{2}{1}}%
\sdlevel{95}{\sdr{11}{28}{1}{1}\sdr{4}{26}{1}{1}\sdr{4}{24}{1}{1}\sdr{5}{27}{1}{1}\sdr{7}{27}{1}{1}\sdr{10}{27}{1}{1}\sdr{12}{27}{1}{1}\sdr{11}{26}{1}{1}\sdr{9}{28}{1}{1}\sdr{11}{17}{1}{1}\sdr{5}{19}{1}{1}\sdr{7}{19}{1}{1}\sdr{10}{19}{1}{1}\sdr{12}{19}{1}{1}\sdr{4}{21}{1}{1}\sdr{6}{26}{1}{1}\sdr{13}{25}{1}{1}\sdr{12}{24}{2}{1}\sdr{3}{22}{1}{2}\sdr{14}{23}{1}{1}\sdr{6}{18}{4}{1}}%
\sdlevel{100}{\sdr{8}{28}{1}{1}\sdr{3}{20}{1}{1}\sdr{5}{20}{1}{1}}%
}
\def\sdpatternsv{(1,13.22) (2,3.97) (3,2.758) (4,1.74) (5,1.381) (6,1.174) (7,1.055) (8,0.8963) (9,0.7737) (10,0.6386) (11,0.6112) (12,0.5046) (13,0.4558) (14,0.4146) (15,0.3737) (16,0.3494) (17,0.3149) (18,0.2571) (19,0.2457) (20,0.194) (21,0.1725) (22,0.1193) (23,0.1135) (24,0.08176) (25,0.05574) (26,0.0279) (27,1.258e-15) (28,1.258e-15) (29,1.258e-15) (30,1.258e-15) (31,1.258e-15) (32,1.685e-16)}
\sdpanel{A}{\sdimg{\sdimgA}}\hfill
\sdpanel{B}{\sdimg{\sdimgB}}\hfill
\sdpanel{C}{\sdimg{\sdimgC}}
\sdpanelrow
\sdpanel{D}{\sdimg{\sdimgD}}\hfill
\sdpanel{E}{\sdimg{\sdimgE}}\hfill
\sdpanel{F}{\begin{tikzpicture}
\begin{axis}[sdplot, width=0.36\textwidth, height=0.36\textwidth, ymode=log, xmin=0, xmax=33, ymin=1e-17, ymax=100,
  xtick={1,8,16,24,32}, ytick={1e-15,1e-10,1e-5,1}, xlabel={$i$}, ylabel={$\sigma_i$}]
\addplot[black, only marks, mark=*, mark size=0.9pt] coordinates {\sdpatternsv};
\addplot[black, only marks, mark=o, mark size=2.2pt] coordinates {(2,3.970) (3,2.758) (5,1.381) (9,0.7737)};
\end{axis}
\end{tikzpicture}}
\caption{Low-rank approximation of a synthetic $32\times32$ pattern $A$ of rank $26$. (A) The pattern. (B) to (E) The truncated SVD $A_k$ for $k=1,2,4,8$, with entries outside $[0,1]$ shown as $0$ or $1$; the relative errors $\|A-A_k\|_{\mathrm F}/\|A\|_{\mathrm F}$ are $0.41$, $0.30$, $0.19$ and $0.11$, and the storage $k(m+n+1)$ is $6\%$, $13\%$, $25\%$ and $51\%$ of $mn$. (F) The singular values on a logarithmic scale; the circles mark $\sigma_{k+1}$, the spectral norm of the error of $A_k$, for the four values of $k$. The rank of $A$, computed exactly in rational arithmetic, is $26$; the six computed values below $10^{-14}$ are rounding errors of the computation.\label{fig:low-rank-pattern}}
\end{figure}

The background alone is a matrix of rank one, a constant column times a linearly increasing row, and so is a rectangle of constant intensity whose sides are parallel to the axes on a white ground. The first term $A_1$, like every matrix of rank one, has all its rows proportional to one row and all its columns to one column, and it shows only a pattern of horizontal and vertical bands. With $k=2$ the ring appears as a dark square block with a lighter centre; with $k=4$ the ring and the square are recognisable and the ellipse appears as a bright bar, but the triangle is only a blurred mound, streaked along the rows and columns; with $k=8$ all four shapes are recognisable, with faint stripes in the background. Shapes with slanted or curved boundaries are not sums of few outer products, and they emerge only gradually. The singular values in panel F decrease slowly, with no gap until $\sigma_{26}$, and for a matrix this small the saving in storage is modest. For a picture of $1000\times1000$ pixels approximated with $k=50$, the storage would be $10\%$ of $mn$; whether such an approximation is acceptable depends on how fast the singular values of the picture decay, which the truncated SVD makes visible but cannot change. \Cref{sec:svd-computation} confirms the identities \eqref{eq:truncation-error} for this pattern numerically.

\begin{note}{Singular values of integral operators}
The singular value decomposition was found twice, for bilinear forms by Beltrami and Jordan in 1873 and 1874~\cite{beltrami1873,jordan1874}, and for integral equations by Schmidt in 1907~\cite{schmidt1907}. Schmidt considered the operator $f\mapsto\int_a^bK(s,t)f(t)\,dt$ with a kernel $K$ that need not be symmetric, and showed that it has an expansion
\[
K(s,t)=\sum_{i\ge1}\sigma_i\,u_i(s)\,v_i(t)
\]
in the mean square sense, with orthonormal functions $u_i$ and $v_i$ and numbers $\sigma_1\ge\sigma_2\ge\cdots>0$ that tend to zero if there are infinitely many of them, the counterpart of \eqref{eq:svd-outer}; he obtained the $v_i$ as eigenfunctions of the symmetric kernel $\int K(r,s)K(r,t)\,dr$, as in our first proof, and he proved that the truncated sum is the best approximation of $K$ by a sum of $k$ products in the mean square sense, the counterpart of \cref{thm:eckart-young}~\cite{stewart1993}. Operators with such expansions, the compact operators, and those with $\sum_i\sigma_i^2<\infty$, the Hilbert--Schmidt operators, are treated in Volume~II.
\end{note}

\section{Perturbation of Singular Values}
\label{sec:svd-perturbation}

\subsection{A Min--Max Principle}

The eigenvalues of Hermitian matrices are perfectly conditioned (\cref{thm:weyl}), while those of general matrices can be extremely sensitive (\cref{tab:jordan-sensitivity}). Singular values behave like the former, for every matrix. The reason is that they are characterised by optimisation, through the Courant--Fischer theorem applied to $A^{\mathsf H}A$.

\begin{theorem}[Min--max characterisation of singular values]
\label{thm:svd-minmax}
Let $A\in\mathbb F^{m\times n}$, with $\sigma_k(A)\coloneqq0$ for $p<k\le n$. For $k=1,\ldots,n$,
\[
\sigma_k(A)=\max_{\dim S=k}\ \min_{\mathbf x\in S,\,\|\mathbf x\|=1}\|A\mathbf x\|=\min_{\dim S=n-k+1}\ \max_{\mathbf x\in S,\,\|\mathbf x\|=1}\|A\mathbf x\|,
\]
where $S$ runs through the subspaces of $\mathbb F^n$ of the dimension indicated; the extremes are attained, at $S=\operatorname{span}\{\mathbf v_1,\ldots,\mathbf v_k\}$ and at $S=\operatorname{span}\{\mathbf v_k,\ldots,\mathbf v_n\}$.
\end{theorem}

\begin{proof}
The Rayleigh quotient of the Hermitian matrix $A^{\mathsf H}A$ is $r(\mathbf x)=\|A\mathbf x\|^2/\|\mathbf x\|^2$, whose values on the unit vectors are the numbers $\|A\mathbf x\|^2$. By \cref{prop:singular-eigen}\,(i), $\sigma_k^2=\lambda_j(A^{\mathsf H}A)$ with $j=n+1-k$, and \cref{thm:courant-fischer} gives $\lambda_j$ as the minimum over subspaces of dimension $j=n-k+1$ of the maximum of $r$, and as the maximum over subspaces of dimension $n-j+1=k$ of the minimum of $r$. Taking square roots, which preserves the order, gives the two formulas. The proof of \cref{thm:courant-fischer} shows where the extremes are attained, with the eigenvectors $\mathbf v_i$ of $A^{\mathsf H}A$.
\end{proof}

For $k=1$ the second formula is $\sigma_1=\max\|A\mathbf x\|=\|A\|_2$, and for $k=n\le m$ the first is $\sigma_n=\min\|A\mathbf x\|$, the two parts of \cref{thm:two-norm}. In between, $\sigma_k$ is the largest number $s$ such that some $k$-dimensional subspace is stretched by at least the factor $s$ in every direction.

\subsection{Weyl's Inequalities for Singular Values}

\begin{theorem}[Perturbation of singular values]
\label{thm:weyl-singular}
Let $A,E\in\mathbb F^{m\times n}$. Then
\[
|\sigma_k(A+E)-\sigma_k(A)|\le\|E\|_2\qquad(k=1,\ldots,p),
\]
and more generally, for $B\in\mathbb F^{m\times n}$ and indices with $i+j-1\le p$,
\[
\sigma_{i+j-1}(A+B)\le\sigma_i(A)+\sigma_j(B).
\]
\end{theorem}

\begin{proof}
For every unit vector $\mathbf x$, $\|(A+E)\mathbf x\|\ge\|A\mathbf x\|-\|E\mathbf x\|\ge\|A\mathbf x\|-\|E\|_2$. So for every subspace $S$ of dimension $k$, the minimum of $\|(A+E)\mathbf x\|$ over the unit vectors of $S$ is at least the minimum of $\|A\mathbf x\|$ minus $\|E\|_2$, and taking the maximum over $S$ gives $\sigma_k(A+E)\ge\sigma_k(A)-\|E\|_2$ by \cref{thm:svd-minmax}. Applied to $A+E$ and $-E$ in place of $A$ and $E$, the same argument gives $\sigma_k(A)\ge\sigma_k(A+E)-\|E\|_2$.

For the second inequality let $A_{i-1}$ and $B_{j-1}$ be truncated SVDs of $A$ and $B$, which have rank at most $i-1$ and $j-1$ and satisfy $\|A-A_{i-1}\|_2=\sigma_i(A)$ and $\|B-B_{j-1}\|_2=\sigma_j(B)$ by \eqref{eq:truncation-error}, where a truncation index beyond the rank means the matrix itself. The matrix $C\coloneqq A_{i-1}+B_{j-1}$ has rank at most $i+j-2$ by \cref{prop:rank-sum}, so \cref{thm:eckart-young} with $k=i+j-2<p$ gives
\[
\sigma_{i+j-1}(A+B)\le\|A+B-C\|_2\le\|A-A_{i-1}\|_2+\|B-B_{j-1}\|_2=\sigma_i(A)+\sigma_j(B).\qedhere
\]
\end{proof}

This is the statement that \cref{sec:computing-rank} quoted: a perturbation of spectral norm $\|E\|_2$ moves no singular value by more than $\|E\|_2$. Together with \cref{thm:eckart-young} it explains why the singular values are the right tool for the rank of perturbed matrices. If $A$ has rank $r$, then $A+E$ has $p-r$ singular values of size at most $\|E\|_2$, which signal the rank deficiency, and $r$ singular values within $\|E\|_2$ of the positive ones of $A$; the gap between the two groups survives as long as $\|E\|_2$ is small compared with $\sigma_r(A)$.

The contrast with eigenvalues is sharp. The Jordan block $J_n(0)$ of \cref{sec:computing-eigenvalues}, with ones above the diagonal, has the eigenvalue $0$ of multiplicity $n$, and the perturbation $\varepsilon\mathbf e_n\mathbf e_1^{\mathsf T}$ of norm $\varepsilon$ moves it to the $n$ numbers of modulus $\varepsilon^{1/n}$, which is $0.1$ for $\varepsilon=10^{-16}$ and $n=16$. The singular values of $J_n(0)$ are $1$, $n-1$ times, and $0$, since the matrix maps $\mathbf e_{j+1}$ to $\mathbf e_j$ and $\mathbf e_1$ to $\mathbf 0$; those of the perturbed matrix, which maps $\mathbf e_{j+1}$ to $\mathbf e_j$ and $\mathbf e_1$ to $\varepsilon\mathbf e_n$, are $1$, $n-1$ times, and $\varepsilon$. The singular values moved by $\varepsilon$, as the theorem allows, while the eigenvalues moved by $\varepsilon^{1/n}$.

The Frobenius norm gives a bound for all the changes together.

\begin{theorem}[Mirsky]
\label{thm:mirsky}
For $A,E\in\mathbb F^{m\times n}$,
\[
\Bigl(\sum_{i=1}^p\bigl(\sigma_i(A+E)-\sigma_i(A)\bigr)^2\Bigr)^{1/2}\le\|E\|_{\mathrm F}.
\]
\end{theorem}

\begin{proof}[sketch]
The Hoffman--Wielandt theorem states that two Hermitian matrices $X$ and $Y$ of order $N$ satisfy $\sum_k(\lambda_k(X)-\lambda_k(Y))^2\le\|X-Y\|_{\mathrm F}^2$, with both sets of eigenvalues in increasing order; the standard proof writes $\|X-Y\|_{\mathrm F}^2$ in terms of the eigenvalues and the unitary matrix that relates the two eigenvector bases, observes that the squared moduli of the entries of this unitary matrix form a doubly stochastic matrix, and minimises a linear function over the doubly stochastic matrices, whose extreme points are the permutation matrices by Birkhoff's theorem~\cite{hoffman1953},~\cite[Sect.~6.3]{horn2013}. Applied to the Hermitian dilations $X=M_{A+E}$ and $Y=M_A$ of \cref{prop:dilation} below, whose eigenvalues in increasing order are $-\sigma_1,\ldots,-\sigma_p$, then $|m-n|$ zeros, then $\sigma_p,\ldots,\sigma_1$, the left side becomes twice the square of the left side of the theorem, and $\|M_E\|_{\mathrm F}^2=2\|E\|_{\mathrm F}^2$. Mirsky proved the inequality, and its analogue for every unitarily invariant norm, in 1960~\cite{mirsky1960}.
\end{proof}

\subsection{The Hermitian Dilation}

A rectangular matrix can be embedded in a Hermitian one of order $m+n$, and the singular values become eigenvalues.

\begin{proposition}[Hermitian dilation]
\label{prop:dilation}
For $A\in\mathbb F^{m\times n}$ the matrix
\[
M_A\coloneqq\begin{bmatrix}0&A\\A^{\mathsf H}&0\end{bmatrix}\in\mathbb F^{(m+n)\times(m+n)}
\]
is Hermitian, and its eigenvalues are $\pm\sigma_1,\ldots,\pm\sigma_p$ and $|m-n|$ zeros. For $i\le p$ the vectors $(\mathbf u_i,\mathbf v_i)/\sqrt2$ and $(\mathbf u_i,-\mathbf v_i)/\sqrt2$ are unit eigenvectors for $\sigma_i$ and $-\sigma_i$. In particular $\|M_A\|_2=\|A\|_2$.
\end{proposition}

\begin{proof}
$M_A^{\mathsf H}=M_A$ is clear. By \eqref{eq:singular-pairs}, $M_A(\mathbf u_i,\pm\mathbf v_i)=(\pm A\mathbf v_i,A^{\mathsf H}\mathbf u_i)=(\pm\sigma_i\mathbf u_i,\sigma_i\mathbf v_i)=\pm\sigma_i(\mathbf u_i,\pm\mathbf v_i)$. If $m>n$, the vectors $(\mathbf u_i,\mathbf 0)$ for $i>p$ satisfy $M_A(\mathbf u_i,\mathbf 0)=(\mathbf 0,A^{\mathsf H}\mathbf u_i)=\mathbf 0$, and if $n>m$, so do the vectors $(\mathbf 0,\mathbf v_i)$ for $i>p$. These $2p+|m-n|=m+n$ vectors are orthonormal, since $\langle(\mathbf u_i,\mathbf v_i),(\mathbf u_j,-\mathbf v_j)\rangle=\delta_{ij}-\delta_{ij}=0$ and the other inner products are computed in the same way, so they form an orthonormal basis of eigenvectors, and their eigenvalues are all the eigenvalues of $M_A$. The largest modulus is $\sigma_1$, so $\|M_A\|_2=\sigma_1$ by \cref{prop:hermitian-norm}.
\end{proof}

The dilation turns statements about Hermitian eigenvalues into statements about singular values. Since $M_{A+E}=M_A+M_E$ and $\|M_E\|_2=\|E\|_2$, Weyl's inequality for Hermitian matrices (\cref{thm:weyl}) gives the first part of \cref{thm:weyl-singular} again: the $k$th largest eigenvalue of $M_{A+E}$ is $\sigma_k(A+E)$ for $k\le p$, and it differs from $\sigma_k(A)$ by at most $\|E\|_2$. Volume~VIII uses the dilation to relate algorithms for the two problems.

\subsection{Interlacing}

Deleting a column or a row of a matrix interlaces its singular values, as deleting a row and a column of a Hermitian matrix interlaces its eigenvalues.

\begin{theorem}[Interlacing of singular values]
\label{thm:singular-interlacing}
Let $A\in\mathbb F^{m\times n}$, with $\sigma_k\coloneqq0$ beyond $p$ for every matrix.
\begin{enumerate}
\item If $B\in\mathbb F^{m\times(n-1)}$ is obtained from $A$ by deleting one column, then $\sigma_k(A)\ge\sigma_k(B)\ge\sigma_{k+1}(A)$ for $k=1,\ldots,n-1$.
\item If $B\in\mathbb F^{(m-1)\times n}$ is obtained from $A$ by deleting one row, then $\sigma_k(A)\ge\sigma_k(B)\ge\sigma_{k+1}(A)$ for $k=1,\ldots,m-1$.
\end{enumerate}
\end{theorem}

\begin{proof}
(i) Deleting the $j$th column of $A$ deletes the $j$th row and column of $A^{\mathsf H}A$, so $B^{\mathsf H}B$ is a principal submatrix of order $n-1$ of $A^{\mathsf H}A$. By \cref{thm:interlacing}, $\lambda_l(A^{\mathsf H}A)\le\lambda_l(B^{\mathsf H}B)\le\lambda_{l+1}(A^{\mathsf H}A)$ for $l=1,\ldots,n-1$. By \cref{prop:singular-eigen}\,(i), $\lambda_l(B^{\mathsf H}B)=\sigma_{n-l}(B)^2$, $\lambda_l(A^{\mathsf H}A)=\sigma_{n+1-l}(A)^2$ and $\lambda_{l+1}(A^{\mathsf H}A)=\sigma_{n-l}(A)^2$; with $k\coloneqq n-l$ this reads $\sigma_{k+1}(A)^2\le\sigma_k(B)^2\le\sigma_k(A)^2$. (ii) Apply (i) to $A^{\mathsf H}$, whose columns are the conjugated rows of $A$, and use \cref{prop:singular-eigen}\,(ii).
\end{proof}

Removing a variable from a least squares problem, a column from its matrix, can therefore only decrease the singular values, but by no more than to the next one: the smallest singular value of the reduced matrix is at least $\sigma_n$ of the full one when $m\ge n$. Conversely, by (i), appending a column to a matrix whose columns stay independent can only increase $\sigma_1$ and decrease the smallest singular value, so the condition number $\sigma_1/\sigma_n$ never decreases when a column is added; this is the monotone growth of $\kappa_2(A)$ with the degree in \cref{tab:polyfit}, where each design matrix contains the columns of the previous one.

\section{The Pseudoinverse}
\label{sec:pseudoinverse}

\subsection{Definition and the Penrose Conditions}

A matrix without an inverse still has, by \cref{thm:svd-subspaces}, an invertible part: it maps the row space bijectively onto the column space, $\mathbf v_i\mapsto\sigma_i\mathbf u_i$. Inverting this part and annihilating the left null space gives a matrix that serves as an inverse wherever one can be expected.

\begin{definition}[Pseudoinverse]
\label{def:pseudoinverse}
Let $A=U\Sigma V^{\mathsf H}\in\mathbb F^{m\times n}$ have rank $r$. Its \term{pseudoinverse}, or Moore--Penrose inverse, is
\[
A^+\coloneqq V\Sigma^+U^{\mathsf H}=\sum_{i=1}^r\frac1{\sigma_i}\,\mathbf v_i\mathbf u_i^{\mathsf H}\in\mathbb F^{n\times m},
\]
where $\Sigma^+\in\mathbb R^{n\times m}$ is the diagonal matrix with the entries $1/\sigma_1,\ldots,1/\sigma_r$ followed by zeros. In particular $0^+=0$.
\end{definition}

The definition uses a particular SVD, and \cref{prop:svd-unique} allows many. That the result does not depend on the choice follows from the next theorem, which characterises $A^+$ by four equations in which no decomposition appears.

\begin{theorem}[Penrose conditions]
\label{thm:penrose}
For $A\in\mathbb F^{m\times n}$, the matrix $X=A^+$ satisfies
\[
AXA=A,\qquad XAX=X,\qquad(AX)^{\mathsf H}=AX,\qquad(XA)^{\mathsf H}=XA ,
\]
and it is the only matrix $X\in\mathbb F^{n\times m}$ that satisfies these four equations. In particular $A^+$ does not depend on the singular value decomposition used to define it.
\end{theorem}

\begin{proof}
$\Sigma\Sigma^+\in\mathbb R^{m\times m}$ and $\Sigma^+\Sigma\in\mathbb R^{n\times n}$ are diagonal with $r$ ones followed by zeros, so $AA^+=U\Sigma\Sigma^+U^{\mathsf H}=U_rU_r^{\mathsf H}$ and $A^+A=V\Sigma^+\Sigma V^{\mathsf H}=V_rV_r^{\mathsf H}$ are Hermitian, and $\Sigma\Sigma^+\Sigma=\Sigma$, $\Sigma^+\Sigma\Sigma^+=\Sigma^+$ give the first two equations. Let $X$ and $Y$ both satisfy the four equations. Using alternately the third and the first equation for $X$ and for $Y$,
\[
\begin{aligned}
X&=XAX=X(AX)^{\mathsf H}=XX^{\mathsf H}A^{\mathsf H}=XX^{\mathsf H}(AYA)^{\mathsf H}=XX^{\mathsf H}A^{\mathsf H}(AY)^{\mathsf H}\\
&=X(AX)^{\mathsf H}(AY)^{\mathsf H}=XAXAY=XAY ,
\end{aligned}
\]
and in the same way, with the fourth equation,
\[
\begin{aligned}
Y&=YAY=(YA)^{\mathsf H}Y=A^{\mathsf H}Y^{\mathsf H}Y=(AXA)^{\mathsf H}Y^{\mathsf H}Y=(XA)^{\mathsf H}A^{\mathsf H}Y^{\mathsf H}Y\\
&=(XA)^{\mathsf H}(YA)^{\mathsf H}Y=XAYAY=XAY .
\end{aligned}
\]
So $X=Y$.
\end{proof}

The four equations are due to Penrose~\cite{penrose1955}, who proved that they have exactly one solution; Moore had introduced the same matrix in 1920~\cite{moore1920} and later characterised it by the property that $AX$ and $XA$ are the orthogonal projections onto the column spaces of $A$ and of $X$~\cite{bengreville2003}, which part (i) of the next proposition verifies for $X=A^+$. The equations make many properties easy to check: to show that a proposed matrix is $A^+$, it suffices to verify the four equations.

\begin{proposition}[Properties of the pseudoinverse]
\label{prop:pinv-properties}
Let $A\in\mathbb F^{m\times n}$ have rank $r$.
\begin{enumerate}
\item $AA^+$ is the orthogonal projection onto $\mathcal C(A)$ and $A^+A$ the orthogonal projection onto $\mathcal C(A^{\mathsf H})$; $I-AA^+$ and $I-A^+A$ are the orthogonal projections onto $\mathcal N(A^{\mathsf H})$ and $\mathcal N(A)$.
\item If $A$ has full column rank, $A^+=(A^{\mathsf H}A)^{-1}A^{\mathsf H}$, a left inverse of $A$; if $A$ has full row rank, $A^+=A^{\mathsf H}(AA^{\mathsf H})^{-1}$, a right inverse of $A$; if $A$ is square and invertible, $A^+=A^{-1}$.
\item $(A^+)^+=A$, $(A^{\mathsf H})^+=(A^+)^{\mathsf H}$, and $(cA)^+=c^{-1}A^+$ for $c\ne0$.
\item The positive singular values of $A^+$ are $1/\sigma_r\ge\cdots\ge1/\sigma_1$, and $\|A^+\|_2=1/\sigma_r$ if $A\ne0$.
\end{enumerate}
\end{proposition}

\begin{proof}
(i) $AA^+=U_rU_r^{\mathsf H}$ and $A^+A=V_rV_r^{\mathsf H}$, by the proof of \cref{thm:penrose}, and these are the projections of \cref{thm:svd-subspaces}. (ii) If $A$ has full column rank, $A^{\mathsf H}A$ is invertible (\cref{prop:gram-rank-complex}), and $X\coloneqq(A^{\mathsf H}A)^{-1}A^{\mathsf H}$ satisfies $XA=I_n$; then $AXA=A$, $XAX=X$, $XA=I$ is Hermitian, and $AX=A(A^{\mathsf H}A)^{-1}A^{\mathsf H}$ is Hermitian, so $X=A^+$ by \cref{thm:penrose}. The full row rank case follows by applying this to $A^{\mathsf H}$ and taking the conjugate transpose, using (iii); for invertible $A$ the inverse satisfies the four equations. (iii) $A^+=V\Sigma^+U^{\mathsf H}$ is, up to the order of the diagonal, a singular value decomposition with the diagonal entries $1/\sigma_i$, so its pseudoinverse is $U(\Sigma^+)^+V^{\mathsf H}=U\Sigma V^{\mathsf H}=A$; the other two statements follow from the four equations, which are preserved by conjugate transposition and by the scaling $X\mapsto c^{-1}X$. (iv) is read off from $V\Sigma^+U^{\mathsf H}$, with \cref{thm:two-norm}.
\end{proof}

Part (ii) is the statement promised after \cref{prop:gram-rank}: the left inverse of the normal equations and the right inverse $A^{\mathsf H}(AA^{\mathsf H})^{-1}$ are special cases of the pseudoinverse. Among the many left inverses of a matrix of full column rank, $A^+$ is the one that annihilates the left null space, $A^+\mathbf w=\mathbf 0$ for $\mathbf w\perp\mathcal C(A)$. For the column $\mathbf a=(1,2)$ of \cref{sec:inverse}, the SVD $\mathbf a=\mathbf u_1\sqrt5\,[1]$ with $\mathbf u_1=\mathbf a/\sqrt5$ gives
\[
\mathbf a^+=[1]\,\frac1{\sqrt5}\,\mathbf u_1^{\mathsf T}=\frac15\begin{bmatrix}1&2\end{bmatrix}=(\mathbf a^{\mathsf T}\mathbf a)^{-1}\mathbf a^{\mathsf T},
\]
the left inverse $(1,2)/5$ that \cref{sec:inverse} singled out, and the only left inverse that maps the orthogonal line spanned by $(-2,1)$ to $0$. In general $\mathbf a^+=\mathbf a^{\mathsf H}/\|\mathbf a\|^2$ for a nonzero column. The library routine \code{linalg.pinv} returned for this column the row $(0.19999999999999996,0.3999999999999999)$, within a few units of the last digit of $(1,2)/5$.

Two rules for inverses do not carry over. The pseudoinverse of a product is not, in general, the product of the pseudoinverses in reverse order: for the $1\times2$ matrix $A$ with the single row $(1,1)$ and the column $B=\mathbf e_1\in\mathbb R^2$, $AB=[1]$ and $(AB)^+=[1]$, while $A^+=(1,1)/2$ and $B^+=\mathbf e_1^{\mathsf T}$, so $B^+A^+=[1/2]$. And the pseudoinverse is not a continuous function of the matrix, as \cref{sec:pinv-continuity} shows.

\subsection{The Minimum Norm Least Squares Solution}

\begin{theorem}[Least squares and the pseudoinverse]
\label{thm:pinv-least-squares}
Let $A\in\mathbb F^{m\times n}$ have rank $r$ and $\mathbf b\in\mathbb F^m$. The vector
\[
\mathbf x^+=A^+\mathbf b=\sum_{i=1}^r\frac{\langle\mathbf u_i,\mathbf b\rangle}{\sigma_i}\,\mathbf v_i
\]
is the minimum norm least squares solution of \cref{prop:minimum-norm}: it minimises $\|A\mathbf x-\mathbf b\|$, and among all minimisers it has the smallest norm. The minimum of the residual is $\|\mathbf b-A\mathbf x^+\|^2=\sum_{i>r}|\langle\mathbf u_i,\mathbf b\rangle|^2$, and every least squares solution has the form $\mathbf x^++\sum_{i>r}c_i\mathbf v_i$.
\end{theorem}

\begin{proof}
By \cref{prop:pinv-properties}\,(i), $A\mathbf x^+=AA^+\mathbf b$ is the orthogonal projection of $\mathbf b$ onto $\mathcal C(A)$, so $\mathbf x^+$ is a least squares solution by \cref{thm:least-squares}\,(ii). It is a combination of $\mathbf v_1,\ldots,\mathbf v_r$, so it lies in $\mathcal C(A^{\mathsf H})$ (\cref{thm:svd-subspaces}), and by \cref{prop:minimum-norm} the only least squares solution in this subspace is the minimum norm solution. The residual is $(I-U_rU_r^{\mathsf H})\mathbf b=\sum_{i>r}\langle\mathbf u_i,\mathbf b\rangle\mathbf u_i$, and the least squares solutions form the flat $\mathbf x^++\mathcal N(A)$, with $\mathcal N(A)$ spanned by $\mathbf v_{r+1},\ldots,\mathbf v_n$.
\end{proof}

The theorem completes \cref{sec:least-squares}, which proved that the minimum norm solution exists and is a linear function of $\mathbf b$, and \cref{sec:systems-practice}, which asked for the shortest solution of an underdetermined system: if $A$ has full row rank, every system $A\mathbf x=\mathbf b$ is consistent, and its shortest solution is $\mathbf x^+=A^{\mathsf H}(AA^{\mathsf H})^{-1}\mathbf b$ by \cref{prop:pinv-properties}\,(ii). For the single equation $x_1+x_2=2$ of \cref{sec:least-squares}, the $1\times2$ matrix $A$ with the row $(1,1)$ has the pseudoinverse $A^+=(1,1)/2$, a column, and $\mathbf x^+=(1,1)$. The formula of the theorem also shows how the solution is computed in practice: rotate $\mathbf b$ into the coordinates $\langle\mathbf u_i,\mathbf b\rangle$, divide by the singular values, and rotate back with $V$. The library routine \code{linalg.lstsq} uses an SVD in this way, through the LAPACK routine \code{gelsd} according to its source code, and treats as zero the singular values below a tolerance relative to $\sigma_1$, as \code{matrix\_rank} does.

\subsection{Discontinuity and Continuity}
\label{sec:pinv-continuity}

The inverse of an invertible matrix depends continuously on the matrix (\cref{cor:perturbed-inverse}). The pseudoinverse does not. For
\[
A_\varepsilon=\begin{bmatrix}1&0\\0&\varepsilon\end{bmatrix},\qquad A_\varepsilon^+=\begin{bmatrix}1&0\\0&1/\varepsilon\end{bmatrix}\quad(\varepsilon\ne0),\qquad A_0^+=\begin{bmatrix}1&0\\0&0\end{bmatrix},
\]
$A_\varepsilon\to A_0$ as $\varepsilon\to0$, but $\|A_\varepsilon^+-A_0^+\|_2=1/|\varepsilon|\to\infty$. The trouble is the change of rank: the small singular value $\varepsilon$ of $A_\varepsilon$, which vanishes in the limit, becomes the large singular value $1/\varepsilon$ of the pseudoinverse. The next theorem shows that this is the only obstruction.

\begin{theorem}[Perturbation of the pseudoinverse]
\label{thm:pinv-perturbation}
Let $A,E\in\mathbb F^{m\times n}$ and $B\coloneqq A+E$.
\begin{enumerate}
\item $B^+-A^+=-B^+EA^+ +B^+(B^+)^{\mathsf H}E^{\mathsf H}(I-AA^+)+(I-B^+B)E^{\mathsf H}(A^+)^{\mathsf H}A^+$, and hence
\[
\|B^+-A^+\|_2\le\bigl(\|A^+\|_2\|B^+\|_2+\|B^+\|_2^2+\|A^+\|_2^2\bigr)\|E\|_2 .
\]
\item Let $A$ have rank $r\ge1$ and $\|E\|_2<\sigma_r(A)$. Then $\operatorname{rank}B\ge r$. If $\operatorname{rank}B=r$, then $\|B^+\|_2\le1/(\sigma_r(A)-\|E\|_2)$; if $\operatorname{rank}B>r$, then $\|B^+\|_2\ge1/\|E\|_2$.
\end{enumerate}
Consequently, if $A_k\to A$ and $\operatorname{rank}A_k=\operatorname{rank}A$ for all $k$, then $A_k^+\to A^+$; and if $\operatorname{rank}A_k\ne\operatorname{rank}A$ for infinitely many $k$, then $\|A_k^+\|_2$ is unbounded along these $k$, and $A_k^+$ does not converge to $A^+$.
\end{theorem}

\begin{proof}
(i) Write $T_1$, $T_2$, $T_3$ for the three terms on the right. By the first and the third Penrose equation for $A$, $A^{\mathsf H}(I-AA^+)=A^{\mathsf H}-(AA^+A)^{\mathsf H}=0$, and by the first and the fourth for $B$, $(I-B^+B)B^{\mathsf H}=B^{\mathsf H}-(BB^+B)^{\mathsf H}=0$. With $E=B-A$ therefore
\begin{align*}
T_1&=-B^+BA^++B^+AA^+,\\
T_2&=B^+(B^+)^{\mathsf H}B^{\mathsf H}(I-AA^+)=B^+(BB^+)^{\mathsf H}(I-AA^+)\\
&=B^+BB^+(I-AA^+)=B^+-B^+AA^+,\\
T_3&=-(I-B^+B)A^{\mathsf H}(A^+)^{\mathsf H}A^+=-(I-B^+B)(A^+A)^{\mathsf H}A^+\\
&=-(I-B^+B)A^+AA^+=-A^++B^+BA^+,
\end{align*}
and the three add up to $B^+-A^+$. For the bound, $I-AA^+$ and $I-B^+B$ are orthogonal projections, of norm at most $1$ (\cref{prop:projection-properties}\,(iv)), and $\|(B^+)^{\mathsf H}\|_2=\|B^+\|_2$.

(ii) By \cref{thm:weyl-singular}, $\sigma_r(B)\ge\sigma_r(A)-\|E\|_2>0$, so $B$ has at least $r$ positive singular values. If $\operatorname{rank}B=r$, then $\|B^+\|_2=1/\sigma_r(B)$ by \cref{prop:pinv-properties}\,(iv), which gives the bound. If $s\coloneqq\operatorname{rank}B>r$, then $\|B^+\|_2=1/\sigma_s(B)\ge1/\sigma_{r+1}(B)$, and $\sigma_{r+1}(B)\le\sigma_{r+1}(A)+\|E\|_2=\|E\|_2$ by the same theorem.

For the consequences, let $E_k\coloneqq A_k-A\to0$ and $r=\operatorname{rank}A$. If $r=0$, then $A^+=0$, and either $A_k=0$ or $\|A_k^+\|_2=1/\sigma_{\operatorname{rank}A_k}(A_k)\ge1/\|E_k\|_2$. If $r\ge1$, then $\|E_k\|_2<\sigma_r(A)/2$ for large $k$, and (ii) gives $\|A_k^+\|_2\le2/\sigma_r(A)$ when the ranks agree, so that the bound of (i) tends to zero, and $\|A_k^+\|_2\ge1/\|E_k\|_2\to\infty$ when they do not.
\end{proof}

The decomposition of (i) is due to Wedin~\cite{wedin1973}; Stewart's survey~\cite{stewart1977} treats the perturbation theory of pseudoinverses, projections and least squares problems in detail. For computation the theorem has a clear message. A matrix of low rank typically becomes a matrix of full rank when it is rounded, as \cref{sec:computing-rank} observed and \cref{prop:rank-perturbation} explains, and its computed pseudoinverse is then huge and meaningless: the tiny singular values that rounding has created are inverted. A pseudoinverse can only be computed reliably for a rank that is decided first, by a tolerance, and the routines \code{linalg.pinv} and \code{linalg.lstsq} do exactly this, inverting only the singular values above $\tau=\rho\sigma_1$ with a relative tolerance $\rho$ that the user may set. The result is the pseudoinverse of the truncated SVD $A_k$, where $k$ is the numerical rank of \cref{prop:numerical-rank}.

\subsection{Regularisation}

The same idea helps with matrices that are invertible but ill conditioned, such as the Hilbert matrices of \cref{sec:pivoting}. Suppose that $A\in\mathbb F^{n\times n}$ is invertible and that the right-hand side is known only with an error, $\mathbf b=A\mathbf x^*+\boldsymbol\delta$. The exact solution of the perturbed system is
\[
A^{-1}\mathbf b=\mathbf x^*+\sum_{i=1}^n\frac{\langle\mathbf u_i,\boldsymbol\delta\rangle}{\sigma_i}\,\mathbf v_i ,
\]
and the components of the error along the singular vectors for small $\sigma_i$ are amplified by the factor $1/\sigma_i$. The \term{truncated SVD solution} discards them,
\[
\mathbf x_k\coloneqq A_k^+\mathbf b=\sum_{i=1}^k\frac{\langle\mathbf u_i,\mathbf b\rangle}{\sigma_i}\,\mathbf v_i ,
\]
at the price of discarding the corresponding components of $\mathbf x^*$ as well.

\begin{proposition}[The error of the truncated SVD solution]
\label{prop:tsvd-error}
Let $A\in\mathbb F^{n\times n}$ be invertible, $\mathbf b=A\mathbf x^*+\boldsymbol\delta$, and $1\le k\le n$. Then
\[
\mathbf x_k-\mathbf x^*=-\sum_{i>k}\langle\mathbf v_i,\mathbf x^*\rangle\mathbf v_i+\sum_{i\le k}\frac{\langle\mathbf u_i,\boldsymbol\delta\rangle}{\sigma_i}\,\mathbf v_i,\qquad\|\mathbf x_k-\mathbf x^*\|^2\le\sum_{i>k}|\langle\mathbf v_i,\mathbf x^*\rangle|^2+\frac{\|\boldsymbol\delta\|^2}{\sigma_k^2}.
\]
\end{proposition}

\begin{proof}
By \eqref{eq:singular-pairs}, $\langle\mathbf u_i,A\mathbf x^*\rangle=\langle A^{\mathsf H}\mathbf u_i,\mathbf x^*\rangle=\sigma_i\langle\mathbf v_i,\mathbf x^*\rangle$, so $\mathbf x_k=\sum_{i\le k}\langle\mathbf v_i,\mathbf x^*\rangle\mathbf v_i+\sum_{i\le k}\sigma_i^{-1}\langle\mathbf u_i,\boldsymbol\delta\rangle\mathbf v_i$, and $\mathbf x^*=\sum_i\langle\mathbf v_i,\mathbf x^*\rangle\mathbf v_i$ (\cref{thm:orthonormal-expansion}). The terms of the error are orthogonal, and by Pythagoras' theorem its squared norm is $\sum_{i>k}|\langle\mathbf v_i,\mathbf x^*\rangle|^2+\sum_{i\le k}\sigma_i^{-2}|\langle\mathbf u_i,\boldsymbol\delta\rangle|^2$; the second sum is at most $\sigma_k^{-2}\sum_i|\langle\mathbf u_i,\boldsymbol\delta\rangle|^2=\sigma_k^{-2}\|\boldsymbol\delta\|^2$ by Parseval's identity (\cref{thm:orthonormal-expansion}).
\end{proof}

The first term, the error of truncation, decreases as $k$ grows, and the second, the propagated error of the data, grows, roughly like $\|\boldsymbol\delta\|/\sigma_k$. A good truncation index keeps the singular values that are large compared with the size of the error in the data and discards the others; \cref{sec:svd-computation} shows this for a Hilbert system. A smoother way to suppress the small singular values is to penalise large solutions.

\begin{theorem}[Tikhonov regularisation]
\label{thm:tikhonov}
Let $A\in\mathbb F^{m\times n}$, $\mathbf b\in\mathbb F^m$ and $\mu>0$. The function $\|A\mathbf x-\mathbf b\|^2+\mu\|\mathbf x\|^2$ has exactly one minimiser,
\[
\mathbf x_\mu=(A^{\mathsf H}A+\mu I)^{-1}A^{\mathsf H}\mathbf b=\sum_{i=1}^r\frac{\sigma_i}{\sigma_i^2+\mu}\,\langle\mathbf u_i,\mathbf b\rangle\,\mathbf v_i=\sum_{i=1}^rf_i\,\frac{\langle\mathbf u_i,\mathbf b\rangle}{\sigma_i}\,\mathbf v_i,\qquad f_i\coloneqq\frac{\sigma_i^2}{\sigma_i^2+\mu},
\]
which is the least squares solution of the system with the stacked matrix and right-hand side
\[
\begin{bmatrix}A\\\sqrt\mu\,I\end{bmatrix}\mathbf x\approx\begin{bmatrix}\mathbf b\\\mathbf 0\end{bmatrix}.
\]
Moreover $\|\mathbf x_\mu\|\le\|\mathbf b\|/(2\sqrt\mu)$, $\|\mathbf x_\mu\|$ does not increase as $\mu$ increases, and $\mathbf x_\mu\to A^+\mathbf b$ as $\mu\to0^+$.
\end{theorem}

\begin{proof}
The stacked matrix $C$ has independent columns, because of the block $\sqrt\mu I$, and $\|C\mathbf x-(\mathbf b,\mathbf 0)\|^2=\|A\mathbf x-\mathbf b\|^2+\mu\|\mathbf x\|^2$. By \cref{thm:least-squares} the minimiser is unique and solves the normal equations $C^{\mathsf H}C\mathbf x=C^{\mathsf H}(\mathbf b,\mathbf 0)$, which are $(A^{\mathsf H}A+\mu I)\mathbf x=A^{\mathsf H}\mathbf b$. With $A=U\Sigma V^{\mathsf H}$, $A^{\mathsf H}A+\mu I=V(\Sigma^{\mathsf T}\Sigma+\mu I)V^{\mathsf H}$ and $A^{\mathsf H}\mathbf b=V\Sigma^{\mathsf T}U^{\mathsf H}\mathbf b$, so $\mathbf x_\mu=V(\Sigma^{\mathsf T}\Sigma+\mu I)^{-1}\Sigma^{\mathsf T}U^{\mathsf H}\mathbf b$, whose $i$th coordinate in the basis $\mathbf v_i$ is $\sigma_i\langle\mathbf u_i,\mathbf b\rangle/(\sigma_i^2+\mu)$ for $i\le r$ and $0$ otherwise. Since $\sigma^2+\mu\ge2\sigma\sqrt\mu$, each coefficient $\sigma_i/(\sigma_i^2+\mu)$ is at most $1/(2\sqrt\mu)$, and it decreases as $\mu$ increases; Parseval's identity gives the first two claims about the norm. As $\mu\to0^+$, $\sigma_i/(\sigma_i^2+\mu)\to1/\sigma_i$ for each $i\le r$, which gives $A^+\mathbf b$ by \cref{thm:pinv-least-squares}.
\end{proof}

The numbers $f_i$ are \term{filter factors}: the regularised solution is the minimum norm solution with the $i$th term multiplied by $f_i$, and the truncated SVD solution has the filter factors $1$ for $i\le k$ and $0$ for $i>k$. \Cref{fig:filter-factors} compares the two filters. Tikhonov's filter passes the terms with $\sigma_i\gg\sqrt\mu$ almost unchanged and damps those with $\sigma_i\ll\sqrt\mu$ by the factor $\sigma_i^2/\mu$, so that they contribute at most $\sigma_i|\langle\mathbf u_i,\mathbf b\rangle|/\mu$; it does not need the SVD, since $\mathbf x_\mu$ can be computed from the stacked least squares problem by the QR factorisation. The method is named after Tikhonov~\cite{tikhonov1963}, and Phillips described the same idea for integral equations~\cite{phillips1962}; Hansen's book~\cite{hansen1998} treats both filters and the choice of $k$ and $\mu$. In statistics the same estimator is ridge regression~\cite{hoerl1970}, which Volume~XV treats, and Volume~XI applies regularisation to smoothing.

\begin{figure}
\sdfloat{%
\begin{tikzpicture}
\begin{axis}[sdplot, width=0.6\textwidth, height=0.38\textwidth, xmode=log, xmin=1e-6, xmax=10, ymin=0, ymax=1.12,
  xtick={1e-6,1e-4,1e-2,1}, ytick={0,0.5,1}, xlabel={$\sigma_i$}, ylabel={$f_i$}, legend pos=north west]
\addplot[black, line width=0.7pt, domain=-6:1, samples=201] ({10^x}, {10^(2*x)/(10^(2*x)+1e-4)});
\addlegendentry{Tikhonov, $\mu=10^{-4}$}
\addplot[black, dashed, line width=0.7pt] coordinates {(1e-6,0) (1e-2,0) (1e-2,1) (10,1)};
\addlegendentry{truncation at $\sigma=10^{-2}$}
\addplot[black!45, line width=0.7pt, domain=-6:1, samples=201] ({10^x}, {10^(2*x)/(10^(2*x)+1e-2)});
\addlegendentry{Tikhonov, $\mu=10^{-2}$}
\end{axis}
\end{tikzpicture}}{The filter factors $f_i=\sigma_i^2/(\sigma_i^2+\mu)$ of Tikhonov regularisation (\cref{thm:tikhonov}) as functions of the singular value, for $\mu=10^{-4}$ and $\mu=10^{-2}$, and the filter factors of the truncated SVD that keeps the singular values above $10^{-2}$. Both filters keep the terms with large singular values and suppress those with small ones; Tikhonov's filter equals $1/2$ at $\sigma_i=\sqrt\mu$ and rises from $0.01$ to $0.99$ between $\sqrt{\mu/99}$ and $\sqrt{99\mu}$, about two decades.\label{fig:filter-factors}}
\end{figure}

\section{Condition Numbers}
\label{sec:svd-condition}

\subsection{The Definition}

\Cref{sec:trace-norms} defined the condition number $\kappa(A)=\|A\|\|A^{-1}\|$ of an invertible matrix in any operator norm, and \cref{sec:householder-qr} defined $\kappa_2(A)$ provisionally, for matrices with independent columns, as the ratio of the longest to the shortest image of a unit vector. The singular values unify the two definitions and extend them to all nonzero matrices.

\begin{definition}[Condition number]
\label{def:kappa2}
The \term{condition number} of a nonzero matrix $A\in\mathbb F^{m\times n}$ of rank $r$ in the spectral norm is
\[
\kappa_2(A)\coloneqq\|A\|_2\|A^+\|_2=\frac{\sigma_1(A)}{\sigma_r(A)} .
\]
\end{definition}

The two expressions agree by \cref{prop:pinv-properties}\,(iv). For a matrix of full rank, $r=p$, the condition number is $\sigma_1/\sigma_p$; for a rank-deficient matrix it ignores the zero singular values. In the second case some authors call $\sigma_1/\sigma_p=\infty$ the condition number instead, and the library routine \code{linalg.cond}, which computes $\sigma_1/\sigma_p$ from the computed singular values, returned $1.39\times10^{17}$ for the singular matrix with all four entries $1$, whose computed second singular value is a rounding error.

\begin{proposition}[Properties of the condition number]
\label{prop:kappa-properties}
Let $A\in\mathbb F^{m\times n}$ be nonzero, of rank $r$.
\begin{enumerate}
\item If $A$ has full column rank, $\kappa_2(A)=\max_{\|\mathbf x\|=1}\|A\mathbf x\|/\min_{\|\mathbf x\|=1}\|A\mathbf x\|$. If $A$ is square and invertible, $\kappa_2(A)=\|A\|_2\|A^{-1}\|_2$.
\item $\kappa_2(A)\ge1$, and $\kappa_2(cA)=\kappa_2(A)$, $\kappa_2(QAZ)=\kappa_2(A)$, $\kappa_2(A^{\mathsf H})=\kappa_2(A^+)=\kappa_2(A)$ for $c\ne0$ and unitary $Q$ and $Z$.
\item $\kappa_2(A^{\mathsf H}A)=\kappa_2(AA^{\mathsf H})=\kappa_2(A)^2$.
\item $\kappa_2(A)=1$ if and only if all positive singular values are equal, that is, $A=cU_rV_r^{\mathsf H}$ with $c>0$. A matrix of full column rank has $\kappa_2(A)=1$ if and only if it is a positive multiple of a matrix with orthonormal columns, and an invertible square matrix if and only if it is a positive multiple of a unitary matrix.
\end{enumerate}
\end{proposition}

\begin{proof}
(i) is \cref{thm:two-norm}\,(i) and (ii) with $r=n=p$, and \cref{cor:distance-singular} for the inverse. (ii) $\sigma_1\ge\sigma_r$; the matrices $cA$, $QAZ$ and $A^{\mathsf H}$ have the singular values of $A$, multiplied by $|c|$ in the first case (\cref{prop:singular-eigen}\,(ii)), and $A^+$ has the reciprocals (\cref{prop:pinv-properties}\,(iv)), whose largest and smallest ratio is the same. (iii) $A^{\mathsf H}A=V\Sigma^{\mathsf T}\Sigma V^{\mathsf H}$ has the positive singular values $\sigma_1^2,\ldots,\sigma_r^2$ and rank $r$, so $\kappa_2(A^{\mathsf H}A)=\sigma_1^2/\sigma_r^2$; similarly for $AA^{\mathsf H}$. (iv) $\sigma_1=\sigma_r=c$ means $A=U_r(cI_r)V_r^{\mathsf H}$. For full column rank, $r=n$ and $V_r=V$ is unitary, so $U_rV^{\mathsf H}$ has orthonormal columns; conversely a matrix $Q$ with orthonormal columns has all singular values $1$, since $Q^{\mathsf H}Q=I$. For invertible square $A$, $r=n$ and $U_rV_r^{\mathsf H}=UV^{\mathsf H}$ is unitary, and a positive multiple $cQ$ of a unitary $Q$ has all singular values equal to $c$.
\end{proof}

Part (i) shows that the definition agrees with those of \cref{sec:trace-norms,sec:householder-qr}, as promised in \cref{ch:linear-systems,ch:inner-products}: for a square invertible matrix, $\kappa_2(A)=\sigma_1/\sigma_n$ is the ratio of the largest to the smallest singular value. Part (iii) is the statement that \cref{sec:least-squares} proved for the L\"auchli matrix only: forming $A^{\mathsf H}A$ squares the condition number, for every matrix. Part (iv) says that the perfectly conditioned matrices are the multiples of unitary ones, the transformations that \cref{sec:orthogonality-computation} found to be safe in floating-point arithmetic. For the matrix with the rows $(1,-2)$ and $(3,4)$,
\[
\kappa_2(A)=\Bigl(\frac{15+5\sqrt5}{15-5\sqrt5}\Bigr)^{1/2}=\Bigl(\frac{3+\sqrt5}{3-\sqrt5}\Bigr)^{1/2}=\frac{3+\sqrt5}2=\phi^2=2.618\ldots,
\]
since $(3+\sqrt5)(3-\sqrt5)=4$; for the matrix $B$ of \cref{sec:svd-geometry}, $\kappa_2(B)=\sigma_1/\sigma_2=\phi^2$ as well. For a positive definite matrix the condition number is $\lambda_n/\lambda_1$, which for the second difference matrix is $\kappa_2(T_n)=\cot^2(\pi/(2(n+1)))\approx4(n+1)^2/\pi^2$ (\cref{sec:rayleigh}), the value that \cref{ch:eigenvalues} anticipated after \cref{thm:second-difference-eigen}.

\subsection{Linear Systems}

\begin{theorem}[The condition number is attained]
\label{thm:kappa-attained}
Let $A\in\mathbb F^{n\times n}$ be invertible, with the singular value decomposition $A=U\Sigma V^{\mathsf H}$.
\begin{enumerate}
\item If $A\mathbf x=\mathbf b\ne\mathbf 0$ and $A(\mathbf x+\Delta\mathbf x)=\mathbf b+\Delta\mathbf b$, then
\[
\frac{\|\Delta\mathbf x\|}{\|\mathbf x\|}\le\kappa_2(A)\,\frac{\|\Delta\mathbf b\|}{\|\mathbf b\|},
\]
with equality when $\mathbf b$ is a multiple of $\mathbf u_1$ and $\Delta\mathbf b$ a multiple of $\mathbf u_n$.
\item The bound of \cref{prop:perturbed-solution} in the norm $\|\cdot\|_2$ is attained: for every $\delta>0$, the matrix $E\coloneqq-\delta\,\mathbf u_n\mathbf v_1^{\mathsf H}$, the vector $\hat{\mathbf x}\coloneqq\mathbf v_1$ and $\mathbf b\coloneqq(A+E)\hat{\mathbf x}$ satisfy $(A+E)\hat{\mathbf x}=\mathbf b$ and, with $\mathbf x\coloneqq A^{-1}\mathbf b$,
\[
\frac{\|\hat{\mathbf x}-\mathbf x\|}{\|\hat{\mathbf x}\|}=\kappa_2(A)\,\frac{\|E\|_2}{\|A\|_2}.
\]
\end{enumerate}
\end{theorem}

\begin{proof}
(i) $\Delta\mathbf x=A^{-1}\Delta\mathbf b$, so $\|\Delta\mathbf x\|\le\|\Delta\mathbf b\|/\sigma_n$, and $\|\mathbf b\|=\|A\mathbf x\|\le\sigma_1\|\mathbf x\|$; multiplying the two inequalities gives the bound. For $\mathbf b=\beta\mathbf u_1$ and $\Delta\mathbf b=\gamma\mathbf u_n$, $\mathbf x=A^{-1}\mathbf b=(\beta/\sigma_1)\mathbf v_1$ and $\Delta\mathbf x=(\gamma/\sigma_n)\mathbf v_n$ by \eqref{eq:singular-pairs}, so both inequalities are equalities. (ii) $\mathbf b=A\mathbf v_1+E\mathbf v_1=\sigma_1\mathbf u_1-\delta\mathbf u_n$, so $\mathbf x=A^{-1}\mathbf b=\mathbf v_1-(\delta/\sigma_n)\mathbf v_n$ and $\hat{\mathbf x}-\mathbf x=(\delta/\sigma_n)\mathbf v_n$. Since $\|\hat{\mathbf x}\|=1$ and $\|E\|_2=\delta$, the left side is $\delta/\sigma_n=(\sigma_1/\sigma_n)(\delta/\sigma_1)$.
\end{proof}

The bounds of \cref{prop:perturbed-solution} are therefore not pessimistic in general: the worst case occurs, and the theorem says where. The solution is most sensitive when the right-hand side points along the direction $\mathbf u_1$ that $A$ stretches most and the error along the direction $\mathbf u_n$ that $A^{-1}$ stretches most. For a particular $\mathbf b$ the error amplification may be much smaller than $\kappa_2(A)$; the bound describes the worst case over all right-hand sides and all perturbations. By \cref{cor:distance-singular}, the same number $1/\kappa_2(A)$ is the relative distance to the nearest singular matrix: an ill-conditioned matrix is one that is close to a singular one, relative to its size, and the condition number measures how close.

\subsection{Least Squares Problems}

For a least squares problem the residual enters the sensitivity of the solution, and the condition number may appear squared even when the normal equations are not formed. The first-order behaviour can be computed exactly.

\begin{proposition}[First-order sensitivity of least squares]
\label{prop:ls-first-order}
Let $A\in\mathbb F^{m\times n}$ have full column rank, $\mathbf b\in\mathbb F^m$, $\mathbf x=A^+\mathbf b$ the least squares solution and $\mathbf r=\mathbf b-A\mathbf x$ its residual. Let $E\in\mathbb F^{m\times n}$ and $\mathbf f\in\mathbb F^m$, and for real $t$ near $0$ let $\mathbf x(t)$ be the least squares solution of $(A+tE)\mathbf x\approx\mathbf b+t\mathbf f$. Then $\mathbf x(t)$ is differentiable at $t=0$, with
\[
\mathbf x'(0)=A^+(\mathbf f-E\mathbf x)+(A^{\mathsf H}A)^{-1}E^{\mathsf H}\mathbf r .
\]
Consequently, if $\|E\|_2\le\varepsilon\|A\|_2$ and $\|\mathbf f\|\le\varepsilon\|\mathbf b\|$, and $\mathbf x\ne\mathbf 0$,
\[
\frac{\|\mathbf x'(0)\|}{\|\mathbf x\|}\le\varepsilon\,\kappa_2(A)\Bigl(2+\bigl(\kappa_2(A)+1\bigr)\frac{\|\mathbf r\|}{\|A\|_2\|\mathbf x\|}\Bigr).
\]
\end{proposition}

\begin{proof}
By \cref{thm:weyl-singular}, $\sigma_n(A+tE)\ge\sigma_n(A)-|t|\|E\|_2>0$ for small $|t|$, so $A(t)\coloneqq A+tE$ has full column rank and $\mathbf x(t)=(A(t)^{\mathsf H}A(t))^{-1}A(t)^{\mathsf H}(\mathbf b+t\mathbf f)$ is unique (\cref{thm:least-squares}). By Cramer's rule (\cref{thm:cramer}) its components are quotients of polynomials in $t$ with a denominator that does not vanish near $0$, so they are differentiable. Differentiating the normal equations $A(t)^{\mathsf H}A(t)\mathbf x(t)=A(t)^{\mathsf H}(\mathbf b+t\mathbf f)$ at $t=0$ by the product rule gives
\[
(E^{\mathsf H}A+A^{\mathsf H}E)\mathbf x+A^{\mathsf H}A\,\mathbf x'(0)=E^{\mathsf H}\mathbf b+A^{\mathsf H}\mathbf f,
\qquad\text{so}\qquad
A^{\mathsf H}A\,\mathbf x'(0)=E^{\mathsf H}\mathbf r+A^{\mathsf H}(\mathbf f-E\mathbf x),
\]
and the formula follows with $A^+=(A^{\mathsf H}A)^{-1}A^{\mathsf H}$ (\cref{prop:pinv-properties}\,(ii)). For the bound, $\|A^+\|_2=1/\sigma_n$ and $\|(A^{\mathsf H}A)^{-1}\|_2=1/\sigma_n^2$, so
\[
\|\mathbf x'(0)\|\le\frac{\varepsilon\|\mathbf b\|+\varepsilon\sigma_1\|\mathbf x\|}{\sigma_n}+\frac{\varepsilon\sigma_1\|\mathbf r\|}{\sigma_n^2}.
\]
Divide by $\|\mathbf x\|$, use $\|\mathbf b\|\le\|A\mathbf x\|+\|\mathbf r\|\le\sigma_1\|\mathbf x\|+\|\mathbf r\|$ and $\sigma_1/\sigma_n=\kappa_2(A)$, and collect the terms.
\end{proof}

The bound has two parts. If the residual is zero, the problem behaves like a linear system, with the factor $2\kappa_2(A)$. If the residual is not small compared with $\|A\|_2\|\mathbf x\|$, the term $\kappa_2(A)^2\|\mathbf r\|/(\|A\|_2\|\mathbf x\|)$ dominates, and it comes from $(A^{\mathsf H}A)^{-1}E^{\mathsf H}\mathbf r$: a perturbation of the matrix tilts the column space, and the large residual, projected onto the tilted space, is amplified by $1/\sigma_n^2$. The effect is real, as the matrix with the rows $(1,0)$, $(0,\varepsilon)$, $(0,0)$ and the right-hand side $\mathbf b=(1,0,1)$ show. The least squares solution is $\mathbf x=(1,0)$, with the residual $\mathbf r=(0,0,1)$, and $\kappa_2(A)=1/\varepsilon$. Changing the entry $a_{32}$ from $0$ to $\delta$ leaves the first component of the solution unchanged and makes the second one minimise $(\varepsilon x_2)^2+(\delta x_2-1)^2$, which gives $x_2=\delta/(\varepsilon^2+\delta^2)$. For $\varepsilon=10^{-3}$ and $\delta=10^{-8}$ this is $0.009999\ldots$: a relative change of $10^{-8}$ in the matrix changed the solution by about $10^{-2}$, which is $\kappa_2(A)^2$ times the perturbation, while $\kappa_2(A)\delta=10^{-5}$. The first-order formula predicts exactly this, since $(A^{\mathsf H}A)^{-1}E^{\mathsf H}\mathbf r=\varepsilon^{-2}\mathbf e_2$ for $E=\mathbf e_3\mathbf e_2^{\mathsf T}$, while $E\mathbf x=\mathbf 0$. A complete bound, valid for finite perturbations, has the same form.

\begin{theorem}[Sensitivity of least squares]
\label{thm:ls-sensitivity}
Let $A\in\mathbb F^{m\times n}$ have full column rank, let $\mathbf x$ be the least squares solution of $A\mathbf x\approx\mathbf b$ with $\mathbf x\ne\mathbf 0$ and residual $\mathbf r$, and let $\hat{\mathbf x}$ be the least squares solution of $(A+E)\hat{\mathbf x}\approx\mathbf b+\mathbf f$, where $\|E\|_2\le\varepsilon\|A\|_2$, $\|\mathbf f\|\le\varepsilon\|\mathbf b\|$ and $\kappa_2(A)\varepsilon<1$. Then
\[
\frac{\|\hat{\mathbf x}-\mathbf x\|}{\|\mathbf x\|}\le\frac{\kappa_2(A)\varepsilon}{1-\kappa_2(A)\varepsilon}\Bigl(2+\bigl(\kappa_2(A)+1\bigr)\frac{\|\mathbf r\|}{\|A\|_2\|\mathbf x\|}\Bigr).
\]
\end{theorem}

\begin{proof}[sketch]
The proof follows the computation of \cref{prop:ls-first-order}, but expresses $\hat{\mathbf x}-\mathbf x$ exactly instead of to first order. One writes $\hat{\mathbf x}-\mathbf x=(A+E)^+(\mathbf b+\mathbf f)-A^+\mathbf b$, splits $\mathbf b=A\mathbf x+\mathbf r$, and uses the orthogonality relations of the Penrose conditions to separate a part that involves $E\mathbf x$ and $\mathbf f$, bounded by $\|(A+E)^+\|_2$, from a part that involves $\mathbf r$ and the tilt of the column space, bounded by $\|(A+E)^+\|_2^2\|E\|_2$; the assumption $\kappa_2(A)\varepsilon<1$ guarantees that $A+E$ has full column rank and $\|(A+E)^+\|_2\le1/(\sigma_n(1-\kappa_2(A)\varepsilon))$, by \cref{thm:weyl-singular}. The bound is due to Wedin~\cite{wedin1973}; Higham~\cite[Thm.~20.1]{higham2002accuracy} gives it in this form with a complete proof, which Volume~VIII presents.
\end{proof}

This is the bound that \cref{sec:orthogonality-computation} quoted. For a method that computes the exact solution of a nearby problem with $\varepsilon$ of the order of $u$, as Householder's QR method does~\cite[Chap.~20]{higham2002accuracy}, the theorem bounds the error by a multiple of $u\kappa_2(A)$ when the residual is small and of $u\kappa_2(A)^2$ only when it is large. The normal equations replace the problem by a linear system whose matrix has the condition number $\kappa_2(A)^2$ (\cref{prop:kappa-properties}\,(iii)), and the bound of \cref{prop:perturbed-solution} for that system is of order $u\kappa_2(A)^2$ whatever the residual; the measurements of \cref{tab:polyfit} were consistent with these orders.

\subsection{Examples}

The Hilbert matrices are the classical ill-conditioned matrices. \Cref{tab:hilbert-kappa} gives $\kappa_2(H_n)$ for the exact matrices, computed from their singular values with the ball arithmetic of python-flint, its growth from $n-1$ to $n$, and two numbers for the matrices with entries rounded to binary64: their exact condition number and the value computed by \code{linalg.cond}. In these rows the condition number grows by a factor between $29$ and $34$ per step, slowly increasing with $n$, and $\kappa_\infty(H_n)$ of \cref{tab:hilbert} is between $1.8$ and $2.5$ times as large. For $n\le12$ the rounded matrices have nearly the same condition number as the exact ones. For $n\ge13$ they do not: $\sigma_n(H_n)$ is then far below the norm of the difference $E$ between the stored and the exact matrix, about $6.5\times10^{-17}$ in the Frobenius norm (\cref{tab:hilbert-cholesky}), and \cref{thm:weyl-singular} only confines the smallest singular value of the stored matrix to the interval from $0$ to $\sigma_n(H_n)+\|E\|_2$. The computed condition number agrees with the exact one of the stored matrix to within $3\%$ for $n\le12$ and differs from it by factors between $1.4$ and $18$ for $n\ge13$; \cref{sec:svd-computation} relates this to the accuracy of the computed singular values.

\begin{table}
\sdfloat{%
\begin{sdtable}{colspec={rllll}, row{1}={font=\small\bfseries}}
\toprule
$n$ & $\kappa_2(H_n)$ & growth & stored, exact & \code{cond} \\
\midrule
$4$ & $1.5514\times10^{4}$ & $29.60$ & $1.551\times10^{4}$ & $1.551\times10^{4}$ \\
$6$ & $1.4951\times10^{7}$ & $31.37$ & $1.495\times10^{7}$ & $1.495\times10^{7}$ \\
$8$ & $1.5258\times10^{10}$ & $32.10$ & $1.526\times10^{10}$ & $1.526\times10^{10}$ \\
$10$ & $1.6026\times10^{13}$ & $32.50$ & $1.602\times10^{13}$ & $1.603\times10^{13}$ \\
$12$ & $1.7132\times10^{16}$ & $32.75$ & $1.682\times10^{16}$ & $1.719\times10^{16}$ \\
$13$ & $5.6279\times10^{17}$ & $32.85$ & $2.172\times10^{18}$ & $4.016\times10^{18}$ \\
$14$ & $1.8534\times10^{19}$ & $32.93$ & $2.898\times10^{17}$ & $4.049\times10^{17}$ \\
$16$ & $2.0223\times10^{22}$ & $33.06$ & $6.304\times10^{17}$ & $1.110\times10^{19}$ \\
\bottomrule
\end{sdtable}}{The condition number $\kappa_2(H_n)=\sigma_1/\sigma_n$ of the Hilbert matrix, computed in ball arithmetic; its growth $\kappa_2(H_n)/\kappa_2(H_{n-1})$; the exact condition number of the matrix with entries rounded to binary64; and the value computed for that matrix by \code{linalg.cond} in binary64.\label{tab:hilbert-kappa}}
\end{table}

The triangular matrix $U_{60}$ of \cref{sec:det-computation}, with $1$ on the diagonal and $-1$ above it, has the determinant $1$ and, computed in ball arithmetic, the smallest singular value $\sigma_{60}=2.6021\times10^{-18}$, below the $2^{-58}\approx3.47\times10^{-18}$ of the perturbation of one entry that made it singular, as \cref{cor:distance-singular} requires, and the condition number $\kappa_2(U_{60})=1.4323\times10^{19}$. The library routine \code{linalg.cond} returned $4.9\times10^{19}$; \cref{sec:svd-computation} explains why a condition number beyond $1/u\approx9\times10^{15}$ cannot be expected to come out accurately from a computation in binary64. The opposite example of that section, $0.1I_n$ with the tiny determinant $10^{-n}$, has $\kappa_2=1$. Determinant and condition number measure different things, and only the second measures nearness to singularity.

\subsection{What the Decomposition Gives}

\Cref{tab:svd-summary} collects the quantities that the singular value decomposition $A=\sum_{i\le r}\sigma_i\mathbf u_i\mathbf v_i^{\mathsf H}$ delivers, each with the result where it is proved.

\begin{table}
\sdfloat{%
\begin{sdtable}{colspec={lll}, row{1}={font=\small\bfseries}}
\toprule
quantity & in terms of the SVD & result \\
\midrule
rank & number of positive $\sigma_i$ & \cref{thm:svd-subspaces} \\
$\mathcal C(A)$, $\mathcal N(A^{\mathsf H})$ & spanned by $\mathbf u_1,\ldots,\mathbf u_r$ and by $\mathbf u_{r+1},\ldots,\mathbf u_m$ & \cref{thm:svd-subspaces} \\
$\mathcal C(A^{\mathsf H})$, $\mathcal N(A)$ & spanned by $\mathbf v_1,\ldots,\mathbf v_r$ and by $\mathbf v_{r+1},\ldots,\mathbf v_n$ & \cref{thm:svd-subspaces} \\
$\|A\|_2$, $\|A\|_{\mathrm F}$, $\|A\|_*$ & $\sigma_1$, $(\sum_i\sigma_i^2)^{1/2}$, $\sum_i\sigma_i$ & \cref{thm:two-norm} \\
distance to rank $<k$ & $\sigma_k$ (spectral norm), attained by $A_{k-1}$ & \cref{thm:eckart-young} \\
numerical rank at $\tau$ & number of $\sigma_i>\tau$ & \cref{prop:numerical-rank} \\
$A^+$, $\mathbf x^+=A^+\mathbf b$ & $\sum_{i\le r}\sigma_i^{-1}\mathbf v_i\mathbf u_i^{\mathsf H}$, $\sum_{i\le r}\sigma_i^{-1}\langle\mathbf u_i,\mathbf b\rangle\mathbf v_i$ & \cref{thm:pinv-least-squares} \\
$\kappa_2(A)$ & $\sigma_1/\sigma_r$ & \cref{def:kappa2} \\
$|\det A|$ (square $A$) & $\sigma_1\cdots\sigma_n$ & \cref{prop:singular-eigen} \\
\bottomrule
\end{sdtable}}{What the singular value decomposition of $A\in\mathbb F^{m\times n}$ of rank $r$ gives, with the result where each entry is proved.\label{tab:svd-summary}}
\end{table}

\section{Polar Decomposition and Principal Components}
\label{sec:polar-pca}

\subsection{The Polar Decomposition}

Every complex number is $z=e^{i\theta}|z|$, a number of modulus one times a nonnegative number. Matrices have the same factorisation, with a unitary factor in place of $e^{i\theta}$ and a positive semidefinite one in place of $|z|$.

\begin{theorem}[Polar decomposition]
\label{thm:polar}
Let $A\in\mathbb F^{m\times n}$ with $m\ge n$. Then
\[
A=WP,
\]
where $W\in\mathbb F^{m\times n}$ has orthonormal columns and $P\in\mathbb F^{n\times n}$ is positive semidefinite. The factor $P=(A^{\mathsf H}A)^{1/2}$ is unique. If $A$ has full column rank, $P$ is positive definite and $W=AP^{-1}$ is unique as well. If $A=U_n\Sigma_nV^{\mathsf H}$ is a thin singular value decomposition, one choice is
\[
W=U_nV^{\mathsf H},\qquad P=V\Sigma_nV^{\mathsf H}.
\]
For square $A$ the factor $W$ is unitary, and $A=P'W$ with $P'\coloneqq WPW^{\mathsf H}=(AA^{\mathsf H})^{1/2}$ as well. For real $A$ both factors can be chosen real.
\end{theorem}

\begin{proof}
With the thin SVD, $W^{\mathsf H}W=VU_n^{\mathsf H}U_nV^{\mathsf H}=I_n$, $P$ is Hermitian with the nonnegative eigenvalues $\sigma_i$, and $WP=U_nV^{\mathsf H}V\Sigma_nV^{\mathsf H}=U_n\Sigma_nV^{\mathsf H}=A$. If $A=WP$ is any such factorisation, then $A^{\mathsf H}A=PW^{\mathsf H}WP=P^2$, and by \cref{thm:square-root} the positive semidefinite square root of $A^{\mathsf H}A$ is unique, so $P=(A^{\mathsf H}A)^{1/2}$. If $A$ has full column rank, $A^{\mathsf H}A$ is positive definite (\cref{thm:pd-criteria}\,(iii)), hence so is $P$, and $W=AP^{-1}$. For square $A$, a square matrix with orthonormal columns is unitary, and $P'W=WPW^{\mathsf H}W=WP=A$; the matrix $WPW^{\mathsf H}$ is positive semidefinite by \cref{prop:pd-properties}\,(iii), and $(WPW^{\mathsf H})^2=WP^2W^{\mathsf H}=(WP)(WP)^{\mathsf H}=AA^{\mathsf H}$, so it is $(AA^{\mathsf H})^{1/2}$ by \cref{thm:square-root}. The real case uses a real SVD.
\end{proof}

The factor $P$ stretches $\mathbb F^n$ along the orthogonal axes $\mathbf v_i$ by the factors $\sigma_i$, and $W$ then moves the stretched space rigidly into $\mathbb F^m$: the decomposition regroups the three steps of \cref{fig:svd-action}: $P=V\Sigma V^{\mathsf H}$ performs the first two and then turns back with $V$, so that it stretches along the axes $\mathbf v_i$ without moving them, and $W=UV^{\mathsf H}$ combines this turning back with the last step $U$. For the matrix $B$ of \cref{sec:svd-geometry}, $U$ and $V$ are the rotations by $\alpha=\arctan\phi$ and $\beta=\arctan\phi^3$, so $W$ is the rotation by $\alpha-\beta$, which is $-\arctan(1/3)$, about $-18.43$ degrees, since $\tan(\beta-\alpha)=(\phi^3-\phi)/(1+\phi^4)=1/3$ by $\phi^2=\phi+1$; and
\[
B=\begin{bmatrix}1&1\\0&2\end{bmatrix}=\frac1{\sqrt{10}}\begin{bmatrix}3&1\\-1&3\end{bmatrix}\cdot\frac1{\sqrt{10}}\begin{bmatrix}3&1\\1&7\end{bmatrix}=WP ,
\]
as one checks by multiplying; $P$ is symmetric with $P^2=B^{\mathsf T}B$.

\Cref{ch:symmetric-matrices} observed that every square factor $X$ with $X^{\mathsf H}X=A$ of a positive definite matrix $A$ is $UA^{1/2}$ for a unitary $U$. This is the polar decomposition of $X$: its factor $P$ is $(X^{\mathsf H}X)^{1/2}=A^{1/2}$. In particular the Cholesky factor $R$ of $A$ has the polar decomposition $R=WA^{1/2}$. For the example of \cref{sec:cholesky},
\[
R=\frac1{\sqrt5}\begin{bmatrix}5&4\\0&3\end{bmatrix}=\frac1{\sqrt5}\begin{bmatrix}2&1\\-1&2\end{bmatrix}\begin{bmatrix}2&1\\1&2\end{bmatrix},
\]
where $A^{1/2}$ is the matrix with the rows $(2,1)$ and $(1,2)$ found there and $W=RA^{-1/2}$ is the rotation by $-\arctan(1/2)$.

\subsection{The Nearest Unitary Matrix}

The unitary factor of the polar decomposition is the unitary matrix closest to $A$, and a slight generalisation answers a question that arises whenever two sets of points are to be compared after a rotation.

\begin{theorem}[Nearest unitary matrix and orthogonal Procrustes problem]
\label{thm:procrustes}
\leavevmode
\begin{enumerate}
\item Let $A,B\in\mathbb F^{m\times n}$, and let $B^{\mathsf H}A=U\Sigma V^{\mathsf H}$ be a singular value decomposition. Among all unitary $Q\in\mathbb F^{n\times n}$, the matrix $Q=UV^{\mathsf H}$ minimises $\|A-BQ\|_{\mathrm F}$.
\item Let $A\in\mathbb F^{n\times n}$ have the polar decomposition $A=WP$. Then $\|A-W\|\le\|A-Q\|$ for every unitary $Q$, both in the Frobenius and in the spectral norm, and
\[
\|A-W\|_{\mathrm F}=\Bigl(\sum_{i=1}^n(\sigma_i-1)^2\Bigr)^{1/2},\qquad\|A-W\|_2=\max_i|\sigma_i-1| .
\]
\end{enumerate}
For real matrices the minimisers are real orthogonal.
\end{theorem}

\begin{proof}
(i) For unitary $Q$, $\|BQ\|_{\mathrm F}=\|B\|_{\mathrm F}$, so
\[
\|A-BQ\|_{\mathrm F}^2=\|A\|_{\mathrm F}^2+\|B\|_{\mathrm F}^2-2\operatorname{Re}\operatorname{tr}(A^{\mathsf H}BQ),
\]
and it suffices to maximise $\operatorname{Re}\operatorname{tr}(A^{\mathsf H}BQ)$. With $A^{\mathsf H}B=(B^{\mathsf H}A)^{\mathsf H}=V\Sigma^{\mathsf T}U^{\mathsf H}$ and the cyclic rule for traces, $\operatorname{tr}(A^{\mathsf H}BQ)=\operatorname{tr}(\Sigma^{\mathsf T}U^{\mathsf H}QV)=\sum_i\sigma_iz_{ii}$, where $Z\coloneqq U^{\mathsf H}QV$ is unitary, so that its columns are unit vectors and $|z_{ii}|\le1$. Hence $\operatorname{Re}\operatorname{tr}(A^{\mathsf H}BQ)\le\sum_i\sigma_i$, with equality for $Z=I$, that is, $Q=UV^{\mathsf H}$. (ii) For the Frobenius norm take $B=I$ in (i): $A=U\Sigma V^{\mathsf H}$ gives the minimiser $UV^{\mathsf H}$, which is the factor $W$ of \cref{thm:polar}, and $A-W=U(\Sigma-I)V^{\mathsf H}$; if $A$ is singular, $W$ is not unique, but every polar factor has this form: if $A=WP$, then $P=V\Sigma V^{\mathsf H}$ by \cref{thm:polar}, so $A=(WV)\Sigma V^{\mathsf H}$ is an SVD and $W=(WV)V^{\mathsf H}$, and (i) applies to it. For the spectral norm, every unitary $Q$ has all singular values $1$, so $\|A-Q\|_2\ge\max_i|\sigma_i(A)-1|$ by \cref{thm:weyl-singular}, and $\|A-W\|_2=\|\Sigma-I\|_2$ attains this bound.
\end{proof}

Part (i) is the \term{orthogonal Procrustes problem}, solved in this form by Sch\"onemann in 1966~\cite{schonemann1966}: given two configurations of points, the rows of $A$ and of $B$, find the rotation or reflection $Q$ that brings the second as close as possible to the first. It arises whenever two sets of corresponding points, measured in different coordinate systems, are to be matched (Volume~XVI); requiring $\det Q=1$, so that reflections are excluded, leads to a variant of the problem that Volume~XVI treats as well. Part (ii) is used to restore the orthogonality of a matrix that should be orthogonal but has been computed with errors: the polar factor is the closest orthogonal matrix, and it can be computed by iterations that converge quickly~\cite{higham2008}. By (ii), a square matrix is within $\varepsilon$ of a unitary matrix in the spectral norm if and only if all its singular values lie in $[1-\varepsilon,1+\varepsilon]$. This gives the loss of orthogonality $\|Q^{\mathsf H}Q-I\|$ of \cref{tab:orthogonality-loss} a precise meaning: if $\|Q^{\mathsf H}Q-I\|_2=\eta<1$, the singular values of $Q$ lie between $\sqrt{1-\eta}$ and $\sqrt{1+\eta}$, by \cref{prop:singular-eigen}\,(i) and \cref{thm:weyl}, and a square $Q$ is then within $\max(1-\sqrt{1-\eta},\sqrt{1+\eta}-1)\le\eta$ of a unitary matrix; since $\|\cdot\|_2\le\|\cdot\|_{\mathrm F}$, the Frobenius norms in the table bound $\eta$.

\subsection{Principal Components}

Let the $N$ rows $\mathbf x_1^{\mathsf T},\ldots,\mathbf x_N^{\mathsf T}$ of $X\in\mathbb R^{N\times n}$ be observations of $n$ variables, and assume that they are centred, $\sum_i\mathbf x_i=\mathbf 0$, as in \cref{sec:gaussian-uses}. The sample covariance matrix is $S=X^{\mathsf T}X/(N-1)$, and for a unit vector $\mathbf w$ the number $\langle\mathbf w,S\mathbf w\rangle=\|X\mathbf w\|^2/(N-1)$ is the sample variance of the projections $\langle\mathbf w,\mathbf x_i\rangle$ of the observations onto the line spanned by $\mathbf w$. \Cref{ch:symmetric-matrices} found the directions of largest variance among the eigenvectors of $S$; the singular value decomposition of $X$ itself delivers them without forming $S$.

\begin{theorem}[Principal components]
\label{thm:pca}
Let $X\in\mathbb R^{N\times n}$, $N\ge2$, have centred columns and the singular value decomposition $X=U\Sigma V^{\mathsf T}$ with singular values $\sigma_1\ge\cdots\ge\sigma_p$ and rank $r$, and let $S=X^{\mathsf T}X/(N-1)$.
\begin{enumerate}
\item $S=VDV^{\mathsf T}$ with $D=\Sigma^{\mathsf T}\Sigma/(N-1)$: the right singular vectors $\mathbf v_j$ are orthonormal eigenvectors of $S$, and the eigenvalue $\sigma_j^2/(N-1)$ is the sample variance of the projections onto $\mathbf v_j$.
\item The \term{scores} $X\mathbf v_j=\sigma_j\mathbf u_j$, the coordinates of the observations along $\mathbf v_j$, have mean zero, and the scores of different components are uncorrelated: $\langle X\mathbf v_j,X\mathbf v_l\rangle=0$ for $j\ne l$.
\item For $1\le k\le r$, among all $W\in\mathbb R^{n\times k}$ with orthonormal columns, $\operatorname{tr}(W^{\mathsf T}SW)$, the total variance of the projections onto the columns of $W$, is largest for $W=V_k=[\,\mathbf v_1\ \cdots\ \mathbf v_k\,]$, with the value $(\sigma_1^2+\cdots+\sigma_k^2)/(N-1)$; if $\sigma_k>\sigma_{k+1}$, only the columns of $V_k$ span a subspace with this property.
\item For $1\le k\le r$, among all orthogonal projections $P$ of $\mathbb R^n$ onto subspaces of dimension $k$, the sum $\sum_i\|\mathbf x_i-P\mathbf x_i\|^2$ of the squared distances of the observations from the subspace is smallest for $P=V_kV_k^{\mathsf T}$, with the value $\sigma_{k+1}^2+\cdots+\sigma_r^2$, and the projected data $XV_kV_k^{\mathsf T}$ are the truncated SVD $X_k$.
\end{enumerate}
\end{theorem}

\begin{proof}
(i) $X^{\mathsf T}X=V\Sigma^{\mathsf T}\Sigma V^{\mathsf T}$, and the eigenvalue for $\mathbf v_j$ is $\langle\mathbf v_j,S\mathbf v_j\rangle$, the sample variance of the projections onto $\mathbf v_j$. (ii) $X\mathbf v_j=\sigma_j\mathbf u_j$ by \eqref{eq:singular-pairs}; the sum of its entries is $\mathbf 1^{\mathsf T}X\mathbf v_j=0$ because the columns of $X$ are centred, and $\langle\sigma_j\mathbf u_j,\sigma_l\mathbf u_l\rangle=0$ for $j\ne l$. (iii) is \cref{lem:ky-fan} for $M=S$, whose $k$ largest eigenvalues are $\sigma_1^2/(N-1),\ldots,\sigma_k^2/(N-1)$. (iv) The $i$th row of $XP$ is $\mathbf x_i^{\mathsf T}P=(P\mathbf x_i)^{\mathsf T}$, since $P$ is symmetric, so $\sum_i\|\mathbf x_i-P\mathbf x_i\|^2=\|X-XP\|_{\mathrm F}^2$. The matrix $XP$ has rank at most $k$, so by \cref{thm:eckart-young} this is at least $\sum_{j>k}\sigma_j^2$. For $P=V_kV_k^{\mathsf T}$, $XP=\sum_{j\le k}X\mathbf v_j\mathbf v_j^{\mathsf T}=\sum_{j\le k}\sigma_j\mathbf u_j\mathbf v_j^{\mathsf T}=X_k$, which attains the bound.
\end{proof}

The vectors $\mathbf v_j$ are the \term{principal directions}, or principal axes, of the data, and the scores are the \term{principal components}; often the scores themselves are called the principal components and the $\mathbf v_j$ the loadings. The theorem says three things at once: the principal directions are the directions of largest variance, one after another and orthogonal to each other; the components are uncorrelated; and the first $k$ principal directions span the $k$-dimensional subspace that fits the observations best in the sense of perpendicular distances, the subspace that Pearson sought in 1901 as the plane of closest fit~\cite{pearson1901}. Hotelling introduced the name and the variance-maximising formulation in 1933~\cite{hotelling1933}. Part (iv) is \cref{thm:eckart-young} in the language of data: the best approximation of rank $k$ of the centred data matrix consists of the projections of the observations onto the first $k$ principal directions. The ratio $(\sigma_1^2+\cdots+\sigma_k^2)/(\sigma_1^2+\cdots+\sigma_r^2)$ is the \term{proportion of the variance} explained by the first $k$ components. \Cref{alg:pca} states the computation.

\begin{algorithm}
\begin{algorithmic}[1]
\Require data $X\in\mathbb R^{N\times n}$, one observation per row, $N\ge2$; the number $k$ of components
\Ensure the mean $\bar{\mathbf x}$, the principal directions $\mathbf v_1,\ldots,\mathbf v_k$, their variances $s_1,\ldots,s_k$ and the scores $Z\in\mathbb R^{N\times k}$
\Name{PrincipalComponents}{X,k}
\State $\bar{\mathbf x}\gets\bigl(\sum_{i=1}^N\mathbf x_i\bigr)/N$ \Comment{the column means}
\State $Y\gets X-\mathbf 1\bar{\mathbf x}^{\mathsf T}$ \Comment{subtract the mean from every row}
\State $U_n\Sigma_nV^{\mathsf T}\gets$ thin SVD of $Y$ \Comment{$S=Y^{\mathsf T}Y/(N-1)$ is never formed}
\For{$j=1,\ldots,k$}
  \State $s_j\gets\sigma_j^2/(N-1)$; $\mathbf z_j\gets\sigma_j\mathbf u_j$ \Comment{variance and scores of component $j$}
\EndFor
\State \Return $\bar{\mathbf x}$, $\mathbf v_1,\ldots,\mathbf v_k$, $s_1,\ldots,s_k$ and $Z=[\,\mathbf z_1\ \cdots\ \mathbf z_k\,]$
\end{algorithmic}
\caption{Principal component analysis by the singular value decomposition (\cref{thm:pca}). The approximation of the data by $k$ components is $\mathbf 1\bar{\mathbf x}^{\mathsf T}+ZV_k^{\mathsf T}$. The cost is dominated by the thin SVD, of the order of $Nn^2$ operations for $N\ge n$ (Volume~VIII).}
\label{alg:pca}
\end{algorithm}

\Cref{fig:pca} shows $N=40$ points in the plane, drawn with the random module of Python (seed $2111$) from the normal distribution with the mean $(2,1.5)$ and the standard deviations $1.6$ and $0.45$ along two perpendicular axes, the first at $32$ degrees to the horizontal; the sample mean is $(2.270,1.549)$. The centred data matrix has the singular values $8.683$ and $2.232$, so the two sample variances along the principal directions are $1.933$ and $0.128$, and the first component explains $94\%$ of the variance. The first principal direction is $\mathbf v_1=(0.8349,0.5505)$, at $33.4$ degrees. The line through the mean in this direction minimises the sum of the squared perpendicular distances of the points; the least squares line of \cref{sec:least-squares}, which predicts the second coordinate from the first, minimises the sum of the squared vertical distances instead and has the smaller slope $0.599$, against the slope $0.659$ of the principal axis. The two answer different questions: least squares treats the first coordinate as exact and the second as measured with errors, while the principal axis treats both coordinates alike.

\begin{figure}
\sdfloat{%
\begin{tikzpicture}
\begin{axis}[sdplot, width=0.62\textwidth, height=0.56\textwidth, axis equal image, xmin=-0.2, xmax=4.7, ymin=-0.3, ymax=3.9,
  xtick={0,1,2,3,4}, ytick={0,1,2,3}, xlabel={$x$}, ylabel={$y$}, legend pos=north west]
\addplot[black, line width=0.8pt, domain=-0.1:4.6, samples=2] {1.54894+0.65937*(x-2.27018)};
\addlegendentry{first principal axis}
\addplot[black, dashed, line width=0.7pt, domain=-0.1:4.6, samples=2] {1.54894+0.59863*(x-2.27018)};
\addlegendentry{least squares line}
\addplot[black!45, line width=0.5pt, domain=0:360, samples=120] ({2.27018+2.78089*cos(x)*0.83485-0.71467*sin(x)*0.55047}, {1.54894+2.78089*cos(x)*0.55047+0.71467*sin(x)*0.83485});
\addplot[black!45, densely dotted, line width=0.7pt] coordinates {(1.82980,2.21682) (2.71056,0.88106)};
\addplot[black, only marks, mark=o, mark size=1.3pt] coordinates {(1.058,1.332) (0.100,0.514) (0.458,0.671) (3.759,2.628) (2.823,1.394) (2.425,0.994) (4.069,2.954) (2.290,2.039) (3.120,2.324) (0.720,0.130) (2.847,1.445) (1.848,1.735) (4.267,2.911) (1.526,0.790) (1.463,1.797) (1.278,1.171) (3.421,1.597) (3.786,1.866) (1.999,0.740) (1.302,0.922) (4.328,2.362) (1.642,1.288) (2.001,1.382) (4.397,3.674) (2.362,1.933) (3.963,2.153) (0.343,0.156) (1.700,1.453) (3.636,2.639) (2.415,1.320) (0.711,0.401) (2.525,1.459) (1.544,0.559) (1.897,1.767) (1.690,1.277) (1.782,1.170) (1.962,0.830) (1.579,1.450) (3.687,3.022) (2.085,1.712)};
\addplot[black, line width=0.5pt] coordinates {(3.421,1.597) (3.09436,2.09238)};
\addplot[black, line width=0.5pt, dashed] coordinates {(3.421,1.597) (3.421,2.23786)};
\addplot[black, only marks, mark=*, mark size=1.4pt] coordinates {(2.27018,1.54894)};
\end{axis}
\end{tikzpicture}}{Forty points in the plane (circles) with their mean (dot), the first principal axis (solid), the second principal direction (dotted, through the mean), and the ellipse with the semi-axes $2\sigma_j/\sqrt{N-1}$ along the principal directions (grey). The least squares line (dashed) minimises the vertical distances, the principal axis the perpendicular distances; the two short segments show both distances for one point.\label{fig:pca}}
\end{figure}

\begin{note}{Total least squares}
Fitting a line by perpendicular distances is the simplest case of \emph{total least squares}, in which errors are allowed in the matrix as well as in the right-hand side of $A\mathbf x\approx\mathbf b$: one seeks the smallest perturbation $[\,E\ \ \mathbf f\,]$ in the Frobenius norm for which $(A+E)\mathbf x=\mathbf b+\mathbf f$ is consistent. By \cref{thm:eckart-young}, applied to the augmented matrix $C=[\,A\ \ \mathbf b\,]\in\mathbb F^{m\times(n+1)}$, every such perturbation, which must make $C$ rank deficient since $(\mathbf x,-1)$ lies in the null space of the perturbed matrix, has at least the norm $\sigma_{n+1}(C)$ when $m>n$. Write the right singular vector of $C$ for $\sigma_{n+1}$ as $\mathbf v_{n+1}=(\mathbf w,\gamma)$ with $\mathbf w\in\mathbb F^n$. If $\gamma\ne0$, the bound is attained by $-\sigma_{n+1}\mathbf u_{n+1}\mathbf v_{n+1}^{\mathsf H}$, and the solution is read off from that vector: $\mathbf x=-\mathbf w/\gamma$. Golub and Van Loan analysed the problem in 1980~\cite{golubvanloan1980}, and \cite[Sect.~6.3]{golub2013} gives the conditions under which the solution exists and is unique. For the points of \cref{fig:pca}, centred, this is the principal axis.
\end{note}

\section{The SVD in Computation: A First Look}
\label{sec:svd-computation}

The stardust library computes singular value decompositions with the routines \code{linalg.svd}, which returns $U$, the singular values and $V^{\mathsf H}$ in the thin form by default, and \code{linalg.svdvals}, which returns the singular values only; according to its source code they call the LAPACK routines \code{gesdd} and \code{gesvd} of the Accelerate library. The routines \code{linalg.matrix\_rank}, \code{linalg.pinv}, \code{linalg.lstsq} and \code{linalg.cond} are built on them, and \code{linalg.matrix\_norm} with \code{ord=2} returns $\sigma_1$. The algorithms, which reduce $A$ to bidiagonal form by Householder reflections and then compute the singular values of the bidiagonal matrix by an iteration or by divide and conquer~\cite{golubkahan1965,golubreinsch1970}, are the subject of Volume~VIII. What matters here is their accuracy: the computed singular values are the exact singular values of a matrix $A+E$ with $\|E\|_2$ bounded by a modest multiple of $u\|A\|_2$~\cite[Sect.~8.6]{golub2013}, so that by \cref{thm:weyl-singular} their \emph{absolute} errors are of the order of $u\sigma_1$, and a singular value much smaller than that has no correct digits. The experiments below test this and other statements of the chapter. The machine and the version of the library are those described in \cref{sec:systems-practice}, and exact reference values were computed in rational arithmetic or with the ball arithmetic of python-flint.

\subsection{Singular Values from the Normal Equations}

The first proof of \cref{thm:svd} suggests computing singular values as the square roots of the eigenvalues of $A^{\mathsf H}A$. \Cref{tab:lauchli-svd} shows what this costs for the L\"auchli matrix of \cref{sec:least-squares}, with the rows $(1,1)$, $(\varepsilon,0)$ and $(0,\varepsilon)$, whose singular values are exactly $(2+\varepsilon^2)^{1/2}$ and $\varepsilon$ for the stored binary64 number $\varepsilon$. We formed $A^{\mathsf T}A$ in binary64 with \code{matmul}, computed its eigenvalues with \code{linalg.eigvalsh} and took the square root of the smaller one, and compared the result, and the smaller value returned by \code{linalg.svdvals} for $A$ itself, with the exact $\sigma_2=\varepsilon$.

\begin{table}
\sdfloat{%
\begin{sdtable}{colspec={rll}, row{1}={font=\small\bfseries}}
\toprule
$\varepsilon$ & \code{svdvals} & $(\text{\code{eigvalsh}}(A^{\mathsf T}A))^{1/2}$ \\
\midrule
$10^{-2}$ & $0$ & $1.1\times10^{-13}$ \\
$10^{-4}$ & $1.4\times10^{-16}$ & $1.8\times10^{-9}$ \\
$10^{-6}$ & $0$ & $4.4\times10^{-5}$ \\
$10^{-7}$ & $1.3\times10^{-16}$ & $4.0\times10^{-4}$ \\
$10^{-8}$ & $1.7\times10^{-16}$ & $1$ (computed value $0$) \\
\bottomrule
\end{sdtable}}{Relative errors of the smaller singular value $\sigma_2=\varepsilon$ of the L\"auchli matrix, computed by \code{linalg.svdvals} from $A$ and as the square root of the smaller eigenvalue of $A^{\mathsf T}A$ formed in binary64.\label{tab:lauchli-svd}}
\end{table}

The routine that works with $A$ was accurate to within twice the unit roundoff in every row. The route through $A^{\mathsf T}A$ lost accuracy as $\varepsilon$ decreased and lost $\sigma_2$ completely at $\varepsilon=10^{-8}$, where $1+\varepsilon^2$ rounds to $1$ and the computed $A^{\mathsf T}A$ is the singular matrix with all entries $1$. The relative errors are consistent with the size of the perturbation: rounding the entries of $A^{\mathsf T}A$ perturbs it by about $u$, which may move its small eigenvalue $\varepsilon^2$ by that much by \cref{thm:weyl}, a relative error of up to about $u/\varepsilon^2$ in $\sigma_2^2$ and half of that in $\sigma_2$, that is, $5.5\times10^{-9}$, $5.5\times10^{-5}$ and $5.5\times10^{-3}$ for $\varepsilon=10^{-4},10^{-6},10^{-7}$; the measured errors are below these values. The larger singular value was accurate in both computations. The SVD routines never form $A^{\mathsf H}A$, for this reason.

\subsection{The Numerical Rank Revisited}

\Cref{tab:numerical-rank} returns to the experiment of \cref{tab:rank-experiment} and lists, for the same kinds of matrices, the largest singular value, the smallest one above the default tolerance $\tau$ of \code{linalg.matrix\_rank}, the largest one at or below it, both as computed by \code{linalg.svdvals} and exactly for the stored matrix, and the numerical rank returned. The tolerance is $\tau=\max(m,n)\,2^{-52}\sigma_1$ according to the source code of the library, in accordance with what \cref{sec:computing-rank} inferred from its behaviour. The products $XY$ were formed with new random integer factors, drawn as in that section with the generator of the library and the seed $20261126$; exact elimination over the integers confirmed that their ranks are $k$.

\begin{table}
\sdfloat{%
\begin{sdtable}{colspec={lllllr}, colsep=4pt, row{1-2}={font=\small\bfseries}}
\toprule
 & & & \SetCell[c=2]{l} next singular value & & \\
\cmidrule[lr]{4-5}
matrix & $\sigma_1$ & last above $\tau$ & computed & exact & rank \\
\midrule
$A$ & $10.12$ & $\sigma_2=1.912$  & $3.6\times10^{-16}$ & $0$ & $2$ \\
$A'$ & $10.12$ & $\sigma_2=1.912$  & $4.2\times10^{-16}$ & $2.6\times10^{-17}$ & $2$ \\
$1,\ldots,9$ & $16.85$ & $\sigma_2=1.068$  & $3.4\times10^{-16}$ & $0$ & $2$ \\
$0.1,\ldots,0.9$ & $1.685$ & $\sigma_2=0.1068$  & $6.9\times10^{-17}$ & $2.3\times10^{-17}$ & $2$ \\
$XY$, $k=3$ & $418.2$ & $\sigma_3=161.1$  & $2.9\times10^{-14}$ & $0$ & $3$ \\
$XY$, $k=10$ & $3991$ & $\sigma_{10}=1982$  & $4.7\times10^{-13}$ & $0$ & $10$ \\
$XY$, $k=50$ & $6003$ & $\sigma_{50}=511.1$  & $1.1\times10^{-12}$ & $0$ & $50$ \\
$H_8$ & $1.696$ & $\sigma_8=1.1\times10^{-10}$  & -- & -- & $8$ \\
$H_{10}$ & $1.752$ & $\sigma_{10}=1.1\times10^{-13}$  & -- & -- & $10$ \\
$H_{12}$ & $1.795$ & $\sigma_{11}=2.6\times10^{-14}$  & $1.04\times10^{-16}$ & $1.07\times10^{-16}$ & $11$ \\
$H_{14}$ & $1.831$ & $\sigma_{11}=3.6\times10^{-13}$  & $4.131\times10^{-15}$ & $4.132\times10^{-15}$ & $11$ \\
\bottomrule
\end{sdtable}}{Singular values and numerical rank of matrices stored in binary64: the matrix $A$ of \cref{sec:four-subspaces}, the same matrix $A'$ with $a_{34}=-1+2^{-52}$, the $3\times3$ matrices with the entries $1,\ldots,9$ and $0.1,\ldots,0.9$, the products $XY$ of order $10$ ($k=3$) and $100$, and Hilbert matrices. The columns give the largest singular value, the last one above the default tolerance $\tau=\max(m,n)\,2^{-52}\sigma_1$ of \code{linalg.matrix\_rank}, the next one as computed by \code{linalg.svdvals} and exactly for the stored matrix, and the rank returned by \code{matrix\_rank}. The tolerance $\tau$ ranges from $3.0\times10^{-15}$ for $H_8$ to $1.3\times10^{-10}$ for the largest product $XY$.\label{tab:numerical-rank}}
\end{table}

The table explains every entry of the last column of \cref{tab:rank-experiment}. For the matrices of exact rank $2$, $3$, $10$ and $50$ the computed singular values beyond the rank are of the order of $u\sigma_1$, as the accuracy of the routine allows, far below $\tau$ and far below the last genuine singular value: the gap is wide, and the numerical rank is the exact rank. The matrix $A'$ with $a_{34}=-1+2^{-52}$ has the exact rank $3$, but its third singular value is $2.6\times10^{-17}$, so by \cref{prop:numerical-rank} it lies within the distance $2.6\times10^{-17}$ of a matrix of rank $2$; the computed value, $4.2\times10^{-16}$, has no correct digit, as expected for a singular value below $u\sigma_1\approx1.1\times10^{-15}$, but it is still far below $\tau$, and the reported rank $2$ is the honest answer at this level of accuracy. The same holds for the matrix of the numbers $0.1,\ldots,0.9$, whose rounding raised the exact rank to $3$ with a third singular value of $2.3\times10^{-17}$. The Hilbert matrices are invertible, but $\sigma_{12}(H_{12})$ and $\sigma_{12}(H_{14})$ lie below the tolerance, the latter only by a factor of $1.4$; with no gap at the tolerance, the numerical rank $11$ of $H_{14}$ would change to $12$ if $\tau$ were halved, in accordance with the dependence on the tolerance that \cref{sec:computing-rank} observed. In the rows of $H_{12}$ and $H_{14}$ the computed and the exact singular values next to the tolerance agree to about $2\%$ and to three digits; their differences, $2.3\times10^{-18}$ and $10^{-18}$, are far below $u\sigma_1\approx2\times10^{-16}$.

\subsection{Regularising an Ill-Conditioned System}

The third experiment applies \cref{sec:pseudoinverse} to the Hilbert system of order $12$. We rounded $H_{12}$ to binary64, computed $\mathbf b=H_{12}\mathbf x^*$ for $\mathbf x^*=(1,\ldots,1)$ exactly and rounded it, and added a perturbation $\boldsymbol\delta$ whose entries were drawn uniformly from $[-1,1)$ by the generator of the library with the seed $20261126$, scaled to $\|\boldsymbol\delta\|=\rho\|\mathbf b\|$ with $\rho=10^{-10}$ and $\rho=10^{-6}$. \Cref{tab:tsvd} shows the relative errors $\|\mathbf x-\mathbf x^*\|/\|\mathbf x^*\|$ of the truncated SVD solutions $\mathbf x_k$, computed as $A_k^+\mathbf b$ with \code{linalg.pinv} and a tolerance between $\sigma_k$ and $\sigma_{k+1}$, and \cref{tab:tikhonov} those of the Tikhonov solutions, computed by \code{linalg.lstsq} from the stacked problem of \cref{thm:tikhonov}, together with the solution by \code{linalg.solve} and by \code{linalg.lstsq} with its default tolerance.

\begin{table}
\sdfloat{%
\begin{sdtable}{colspec={rlll}, row{1}={font=\small\bfseries}}
\toprule
$k$ & $\sigma_k$ & $\rho=10^{-10}$ & $\rho=10^{-6}$ \\
\midrule
$1$ & $1.80$ & $5.6\times10^{-1}$ & $5.6\times10^{-1}$ \\
$2$ & $3.80\times10^{-1}$ & $2.0\times10^{-1}$ & $2.0\times10^{-1}$ \\
$3$ & $4.47\times10^{-2}$ & $5.8\times10^{-2}$ & $5.8\times10^{-2}$ \\
$4$ & $3.72\times10^{-3}$ & $1.5\times10^{-2}$ & $1.5\times10^{-2}$ \\
$5$ & $2.33\times10^{-4}$ & $3.4\times10^{-3}$ & $3.6\times10^{-3}$ \\
$6$ & $1.12\times10^{-5}$ & $6.7\times10^{-4}$ & $2.0\times10^{-2}$ \\
$7$ & $4.08\times10^{-7}$ & $2.3\times10^{-4}$ & $1.6$ \\
$8$ & $1.12\times10^{-8}$ & $1.9\times10^{-3}$ & $6.2\times10^{1}$ \\
$9$ & $2.25\times10^{-10}$ & $3.4\times10^{-2}$ & $6.7\times10^{2}$ \\
$10$ & $3.11\times10^{-12}$ & $1.0\times10^{1}$ & $2.8\times10^{4}$ \\
$11$ & $2.65\times10^{-14}$ & $1.1\times10^{3}$ & $3.1\times10^{7}$ \\
$12$ & $1.04\times10^{-16}$ & $1.8\times10^{5}$ & $3.1\times10^{8}$ \\
\bottomrule
\end{sdtable}}{The Hilbert system of order $12$ with a perturbed right-hand side, $\|\boldsymbol\delta\|=\rho\|\mathbf b\|$: the computed singular values of the stored $H_{12}$ and the relative errors of the truncated SVD solutions $\mathbf x_k$.\label{tab:tsvd}}
\end{table}

\begin{table}
\sdfloat{%
\begin{sdtable}{colspec={lll}, row{1}={font=\small\bfseries}}
\toprule
method & $\rho=10^{-10}$ & $\rho=10^{-6}$ \\
\midrule
\code{solve} & $1.7\times10^{5}$ & $2.9\times10^{8}$ \\
\code{lstsq}, default tolerance & $1.1\times10^{3}$ & $3.1\times10^{7}$ \\
Tikhonov, $\mu=10^{-4}$ & $5.3\times10^{-2}$ & $5.3\times10^{-2}$ \\
Tikhonov, $\mu=10^{-6}$ & $1.5\times10^{-2}$ & $1.5\times10^{-2}$ \\
Tikhonov, $\mu=10^{-8}$ & $4.1\times10^{-3}$ & $4.5\times10^{-3}$ \\
Tikhonov, $\mu=10^{-10}$ & $1.6\times10^{-3}$ & $9.9\times10^{-3}$ \\
Tikhonov, $\mu=10^{-12}$ & $5.5\times10^{-4}$ & $2.3\times10^{-1}$ \\
Tikhonov, $\mu=10^{-14}$ & $1.8\times10^{-4}$ & $1.7$ \\
\bottomrule
\end{sdtable}}{The same problem as in \cref{tab:tsvd}: relative errors of the solution by $PA=LU$ (\code{solve}), by \code{lstsq} with its default tolerance, and of the Tikhonov solutions $\mathbf x_\mu$.\label{tab:tikhonov}}
\end{table}

The solution of the system as it stands is useless at both noise levels, with relative errors of $10^5$ and $10^8$, and \code{lstsq} with its default tolerance, which discarded only $\sigma_{12}$, is the truncated solution $\mathbf x_{11}$ and no better. The truncated solutions behave as \cref{prop:tsvd-error} describes. Their errors first decrease with $k$, as more of $\mathbf x^*$ is captured, and then increase, as the propagated error of the data grows like $\|\boldsymbol\delta\|/\sigma_k$; the smallest errors, $2.3\times10^{-4}$ for $\rho=10^{-10}$ at $k=7$ and $3.6\times10^{-3}$ for $\rho=10^{-6}$ at $k=5$, occurred where $\sigma_k$ exceeded $\|\boldsymbol\delta\|$, which is $5.2\times10^{-10}$ and $5.2\times10^{-6}$, by factors of about $780$ and $45$, and $\sigma_{k+1}$ by factors of about $21$ and $2$. For $k\le4$ the errors agree at both noise levels to the digits shown; by \cref{prop:tsvd-error} the propagated part of the relative error is then at most $\|\boldsymbol\delta\|/(\sigma_k\|\mathbf x^*\|)\le4.1\times10^{-4}$. The Tikhonov solutions show the same pattern in $\mu$: their errors were smallest, $1.8\times10^{-4}$ and $4.5\times10^{-3}$, for the values of $\mu$ in the table whose square roots, $10^{-7}$ and $10^{-4}$, lie between the same singular values as the optimal truncation. In both methods the best regularisation parameter depends on the size of the error in the data, which in practice must be estimated; how to choose $k$ or $\mu$ without knowing $\mathbf x^*$ is the subject of Volume~XI.

\subsection{The Theorem of Eckart and Young in Numbers}

The last experiment checks \eqref{eq:truncation-error} for the pattern of \cref{fig:low-rank-pattern}. We computed the SVD with \code{linalg.svd}, formed $A_k$ from its factors, and evaluated $\|A-A_k\|_2$ with \code{linalg.matrix\_norm} and \code{ord=2}, which computes the largest singular value of the difference, and $\|A-A_k\|_{\mathrm F}$ from the entries. \Cref{tab:eckart-young-check} compares them with $\sigma_{k+1}$ and with $(\sum_{i>k}\sigma_i^2)^{1/2}$, computed from the singular values returned by \code{linalg.svd}.

\begin{table}
\sdfloat{%
\begin{sdtable}{colspec={rllll}, row{1}={font=\small\bfseries}}
\toprule
$k$ & $\sigma_{k+1}$ & rel.\ difference & $(\sum_{i>k}\sigma_i^2)^{1/2}$ & rel.\ difference \\
\midrule
$1$ & $3.970487$ & $3.4\times10^{-16}$ & $5.848860$ & $0$ \\
$2$ & $2.758451$ & $1.6\times10^{-16}$ & $4.294694$ & $0$ \\
$4$ & $1.381066$ & $0$ & $2.794303$ & $3.2\times10^{-16}$ \\
$8$ & $0.7737125$ & $5.7\times10^{-16}$ & $1.614538$ & $1.4\times10^{-16}$ \\
$16$ & $0.3149458$ & $5.3\times10^{-16}$ & $0.5750852$ & $5.8\times10^{-16}$ \\
\bottomrule
\end{sdtable}}{The errors of the truncated SVD of the $32\times32$ pattern of \cref{fig:low-rank-pattern}: the singular value $\sigma_{k+1}$ and the relative difference between it and the computed $\|A-A_k\|_2$, and the tail $(\sum_{i>k}\sigma_i^2)^{1/2}$ and its relative difference from the computed $\|A-A_k\|_{\mathrm F}$.\label{tab:eckart-young-check}}
\end{table}

The identities held to within six units of the roundoff $u$ in every case, for the spectral norm, which the library computes by the same kind of algorithm as the singular values, and for the Frobenius norm, which it computes from the entries without any decomposition.

\bigskip
Every matrix $A\in\mathbb F^{m\times n}$ has a singular value decomposition $A=U\Sigma V^{\mathsf H}$ with unitary $U$ and $V$ and the singular values $\sigma_1\ge\cdots\ge\sigma_p\ge0$ on the diagonal of $\Sigma$; geometrically, the image of the unit sphere is an ellipsoid with the semi-axes $\sigma_i\mathbf u_i$, the images of the orthonormal vectors $\mathbf v_i$. The decomposition follows from the spectral theorem for $A^{\mathsf H}A$, or directly by maximising $\|A\mathbf x\|$; the singular values are unique and are the square roots of the eigenvalues of $A^{\mathsf H}A$, $\sigma_i^2=\lambda_{n+1-i}(A^{\mathsf H}A)$ in the two customary orders, and for normal matrices they are the moduli of the eigenvalues. The singular vectors are orthonormal bases of the four fundamental subspaces, the number of positive singular values is the rank, and the singular values are the complete invariants of unitary equivalence. The spectral norm is $\sigma_1$, the Frobenius norm is the Euclidean length of the vector of singular values, and the nuclear norm their sum. The truncated SVD $A_k$ is a best approximation of rank $k$ in both norms, with the errors $\sigma_{k+1}$ and $(\sum_{i>k}\sigma_i^2)^{1/2}$, so that $\sigma_n$ is the distance to the nearest singular matrix and the number of singular values above $\tau$ is the smallest rank within the distance $\tau$. Singular values satisfy a min--max principle, Weyl's and Mirsky's perturbation bounds and interlacing, and so are perfectly conditioned for every matrix. The pseudoinverse $V\Sigma^+U^{\mathsf H}$ is characterised by the four Penrose equations, gives the minimum norm least squares solution, and is continuous exactly where the rank does not change; truncation and Tikhonov regularisation filter out the small singular values. The condition number $\kappa_2=\sigma_1/\sigma_r$ is attained as an error amplification, is the reciprocal of the relative distance to singularity, is squared by the normal equations, and governs the sensitivity of least squares together with the residual. Finally, the polar decomposition $A=WP$ gives the nearest unitary matrix and solves the orthogonal Procrustes problem, and the SVD of a centred data matrix gives its principal components. \Cref{ch:canonical-forms} relates the singular values to the Schur form, through Schur's inequality and Weyl's inequalities between eigenvalues and singular values.

\chapternotes

Stewart~\cite{stewart1993} tells the early history of the singular value decomposition, from which the following account is taken. Beltrami in 1873~\cite{beltrami1873} and Jordan in 1874~\cite{jordan1874} found independently that a real bilinear form can be reduced to a diagonal form by two orthogonal substitutions, Beltrami for square nonsingular matrices with distinct singular values; Sylvester gave another derivation in 1889~\cite{sylvester1889}. Schmidt~\cite{schmidt1907} developed the decomposition for integral operators in 1907 and proved the approximation theorem that is now named after Eckart and Young, who stated it for matrices in 1936~\cite{eckart1936}; Weyl~\cite{weyl1912} proved in 1912, for integral equations with symmetric kernels, the inequalities that \cref{thm:weyl-singular} states for singular values, and remarked that his argument applies to unsymmetric kernels as well. Mirsky~\cite{mirsky1960} proved the approximation and perturbation theorems for every unitarily invariant norm, building on von Neumann's characterisation of these norms~\cite{vonneumann1937}, and Ky Fan proved his maximum principle in 1949~\cite{fan1949} and introduced the norms named after him in 1951~\cite{fan1951}. Moore~\cite{moore1920} and Penrose~\cite{penrose1955} defined the pseudoinverse, and Ben-Israel and Greville~\cite{bengreville2003} treat generalised inverses in depth; Wedin~\cite{wedin1973} and Stewart~\cite{stewart1977} developed the perturbation theory of pseudoinverses and least squares problems. Golub and Kahan~\cite{golubkahan1965} proposed the reduction to bidiagonal form on which the computation of the SVD rests, and Golub and Reinsch~\cite{golubreinsch1970} published an algorithm based on it; Golub and Van Loan~\cite{golubvanloan1980} analysed total least squares. Pearson~\cite{pearson1901} fitted lines and planes by perpendicular distances and Hotelling~\cite{hotelling1933} introduced principal components; Sch\"onemann~\cite{schonemann1966} solved the orthogonal Procrustes problem; and regularisation goes back to Phillips~\cite{phillips1962} and Tikhonov~\cite{tikhonov1963}, and in statistics to Hoerl and Kennard~\cite{hoerl1970}.

For the theory, Horn and Johnson~\cite[Sects.~2.6, 7.3 and~7.4]{horn2013} prove the SVD, the polar decomposition and their consequences, and their second book~\cite[Chap.~3]{horn1991topics} is devoted to singular value inequalities; Strang~\cite{strang2016} and Axler~\cite[Chap.~7]{axler2024} give accessible introductions, the second for operators. Trefethen and Bau~\cite[Lectures~4, 5 and~31]{trefethen1997} explain the geometry, the low-rank approximation and the computation of the SVD in a few pages, and Golub and Van Loan~\cite[Sects.~2.4 and~8.6]{golub2013} treat its theory and computation in depth. Stewart and Sun~\cite{stewartsun1990} is the standard reference for the perturbation theory of singular values, subspaces and pseudoinverses, Higham~\cite[Chap.~20]{higham2002accuracy} and Bj\"orck~\cite{bjorck1996} for least squares, and Hansen~\cite{hansen1998} for regularisation. Higham~\cite[Chap.~8]{higham2008} treats the polar decomposition and its computation, Jolliffe~\cite{jolliffe2002} principal component analysis, and Halko, Martinsson and Tropp~\cite{halko2011} the randomised algorithms that compute low-rank approximations of very large matrices from a few products with random vectors, which Volume~IX takes up.

\chapterreferences
